% main.tex -- C3 paper, version 1 (2026-11-08).
\documentclass[11pt]{article}
\usepackage[T1]{fontenc}
\usepackage[utf8]{inputenc}
\usepackage{lmodern}
\usepackage[margin=1in]{geometry}
\usepackage{amsmath,amssymb,amsthm}
\usepackage{booktabs,array}
\usepackage[hidelinks]{hyperref}

\newtheorem{theorem}{Theorem}[section]
\newtheorem{proposition}[theorem]{Proposition}
\newtheorem{lemma}[theorem]{Lemma}
\newtheorem{corollary}[theorem]{Corollary}
\theoremstyle{definition}
\newtheorem{definition}[theorem]{Definition}
\theoremstyle{remark}
\newtheorem{remark}[theorem]{Remark}

\newcommand{\lean}[1]{\footnote{\raggedright Lean: \texttt{#1}.}}

\title{Polynomial calculus degree of the onto functional pigeonhole principle
when the characteristic divides the surplus}
\author{Bj\"orn Lindheim}
\date{Version 1, November 2026}

\begin{document}
\maketitle

\begin{abstract}
The onto functional pigeonhole principle with $n+k$ pigeons and $n$ holes has a polynomial
calculus refutation of degree $1$ over $\mathbb F_p$ when $p\nmid k$, and one of degree
$p^{v_p(k)}$ in general, by a criterion of Riis.  We study the lower bound when $p$ divides
$k$.  In Lean~4 we prove that the degree is exactly $p$ when $p$ divides $k$ exactly once
($n\ge2p-3$), at least $p+1$ when $p^2$ divides $k$ ($n\ge2p-1$), and at least $5$ for $p=2$,
$k=8$ ($n\ge7$).  With written proofs that rest on named published theorems and have been
checked by AI reviewers but not by people, we show that the degree is exactly $p^{v_p(k)}$
for $n\ge2p^{v_p(k)}-1$, so that for surplus $p^\ell$ it grows as $p^\ell$; we give two
further mechanisms, one through Nullstellensatz designs (degree at least $2^\ell$ with
explicit thresholds) and one degree by degree (exact degree $p^2$ when $v_p(k)=2$); and we
record five errata in two of the sources these proofs use, among them a gap in the published
proof of a theorem we would otherwise have used.  Every statement carries its
evidence tier.  All results were produced by AI research agents under the author's direction.
\end{abstract}

% sec1-introduction.tex -- Section 1 of the C3 paper (draft 2026-11-08).
% Sources: OUTLINE.md; publish/summaries/c3-*.md; requests-20261106/c3-publication-prior-art-v2.md;
% notes-20261025/c3-sources.md; codex-20261104-runB-adjudication/00, 01, 05; notes-20261104/referee-runB-*.md;
% PLAN.md rows G-C3.1--G-C3.3.  Every tier and threshold below was checked against those files on 2026-11-08.

\section{Introduction}\label{sec:intro}

\subsection{The principle and the question}\label{sec:intro-question}

The onto functional pigeonhole principle $\mathrm{onto\text{-}FPHP}^{n+k}_n$ states that
$n+k$ pigeons can be mapped bijectively onto $n$ holes.  It is false for every $k\ge1$.
We use the \emph{linear encoding}: Boolean variables $x_{ij}$ (pigeon $i$ sits in hole $j$)
and the axioms
\[
  \textstyle\sum_{j}x_{ij}-1,\qquad \sum_{i}x_{ij}-1,\qquad x_{ij}x_{i'j}\ (i\ne i'),\qquad x_{ij}x_{ij'}\ (j\ne j'),
\]
saying that every pigeon sits in a hole, every hole holds exactly one pigeon, no two
pigeons share a hole and no pigeon sits in two holes.  This is the system of Beame and
Riis~\cite[Def.~7.2]{beame1998counting}, also used in Razborov's
survey~\cite{razborov2001php}.  All results of this paper are for this encoding.
Translating the CNF clause by clause instead turns each pigeon clause into a polynomial of
degree $n$~\cite[footnote~5]{razborov2001php}, so degree bounds below $n$ say nothing about
that encoding.

We measure refutations in the polynomial calculus (PC) over a field of characteristic
$p$: from the axioms and the Boolean axioms $x^2-x$, derive $1$ by linear combinations and
multiplication by a variable.  The \emph{degree} of a refutation is the largest degree of
a line, and the \emph{PC degree} of a system is the least degree of a refutation.
Section~\ref{sec:defs} fixes the conventions exactly.

If $p\nmid k$, summing the pigeon axioms and subtracting the hole axioms gives the constant
$-k$, a refutation of degree~$1$.  The question is what happens when $p$ divides $k$.
There the system has a refutation of degree $p^{v}$, where $v=v_p(k)$ is the exponent of
$p$ in $k$: Riis's criterion~\cite{riis1993independence} (Theorem~8 in
\cite{razborov2001php}) gives a refutation of degree $d$ whenever $\binom{n+k}{d}$ and
$\binom nd$ differ mod $p$, and by Lucas's theorem
$\binom{n+k}{p^v}-\binom{n}{p^v}\equiv k/p^v\not\equiv0\pmod p$.  The criterion cannot give
a smaller degree.  So for fixed $p$ and $k$ the PC degree is bounded independently of $n$,
and the questions are its exact value and how it grows with $v$.

\subsection{What was known}\label{sec:intro-known}

\paragraph{Nullstellensatz.}  Beame and Riis~\cite[Thm.~8.1, Lemma~8.7]{beame1998counting}
prove, for surplus $k=p^\ell$, that the Nullstellensatz degree is at least $2^\ell-1$ for
$n\ge((p+2)^\ell-p^\ell)/2$ (Theorem~8.1) and at most $p^\ell-1$ for $n\ge p^\ell$
(Lemma~8.7).  Their degree measures the multipliers.  PC line degree is at most
Nullstellensatz \emph{product} degree; for this encoding a certificate of multiplier degree
$d$ gives a PC refutation of degree at most $d+2$.  Nullstellensatz lower bounds alone
therefore do not give PC lower bounds.  Beame and Riis list non-trivial PC
lower bounds for this family as open~\cite[\S10]{beame1998counting}.

\paragraph{Razborov's Question~7.}  Razborov's survey~\cite[\S7, Question~7]{razborov2001php}
asks, for fixed $p$, whether $\mathrm{onto\text{-}FPHP}^{n+p^\ell}_n$ has no bounded-degree
PC refutations in characteristic $p$, noting that the Nullstellensatz analogue is known.
As printed, the question fixes $p$ and $\ell$ and lets $n$ grow.  For fixed $p$ and $\ell$,
Riis's upper bound $p^\ell$ above is independent of $n$.  We also consider the distinct
question of growth as $\ell$ varies.  We state both readings and say which one each result addresses; none of
our results settles the question in a sense we would call an answer, and we do not claim
one.

\paragraph{Riis and Sitharam.}  Theorem~33 of the technical report of Riis and
Sitharam~\cite{riis1997nonconstant} states a PC lower bound for the same family: degree at
least $2^\ell-1$ for surplus $p^\ell$ once
$N\ge(21+3p^2)(2^\ell-2)-p$.  This is growth in $\ell$.  We cite it for comparison only:
we have read its statement and the surrounding lemmas~\cite[pp.~36--38]{riis1997nonconstant}
but have not checked its proof, which delegates to an earlier tree construction.

\paragraph{Other work.}  General PC degree methods~\cite{miksa2015generalized} and the
cross-characteristic lower bounds for Tseitin and counting principles~\cite{bgip2001,alekhnovich2003nonbinomial,bikprs1996}
concern regimes in which the characteristic does not divide the relevant modulus;
Impagliazzo, Pudl\'ak and Sgall~\cite{ips1999} treat the injective principle without the
onto and functional axioms, where the bound does not depend on $p$.  None of them gives a
low-degree reduction to $\mathrm{onto\text{-}FPHP}^{n+k}_n$ with $p\mid k$.
Section~\ref{sec:literature} gives the details and the search behind this paragraph.

\subsection{Results}\label{sec:intro-results}

Table~\ref{tab:results} lists the results with the evidence behind each.
The tiers are defined in Section~\ref{sec:intro-tiers}; in short, T-LEAN means
kernel-checked in Lean~4, T-WRITTEN means a written proof checked by AI reviewers, and
T-MEASURED means an exact finite computation.  No person has reviewed any of the proofs.

\begin{table}[ht]
\centering\small
\renewcommand{\arraystretch}{1.25}
\begin{tabular}{@{}p{0.03\textwidth}p{0.57\textwidth}>{\raggedright\arraybackslash}p{0.24\textwidth}p{0.06\textwidth}@{}}
\toprule
 & Statement (characteristic $p$ prime, surplus $k$) & Tier & Where \\
\midrule
A & $k=pc$, $p\nmid c$, $n\ge2p-3$: PC degree exactly $p$.  Hence the same formula has PC
    degree $p$ over $\mathbb F_p$ and $1$ over $\mathbb F_q$ for every prime $q\nmid pc$.
  & T-LEAN & \S\ref{sec:lean-A} \\
B & $k=cp^2$, $c\ge1$, $n\ge2p-1$: no refutation of degree $\le p$, so PC degree
    $\ge p+1$.  At $p=2$, $k=4$, $n\ge5$: exactly $4$.
  & T-LEAN & \S\ref{sec:lean-B} \\
C & $p=2$, $k=8$, $n\ge7$: no refutation of degree $\le4$, so PC degree $\ge5$
    (the upper bound is $8$).
  & T-LEAN; certificate re-checked independently & \S\ref{sec:lean-C} \\
D & Growth.  $k=p^\ell$, $\ell\ge1$,
    $n\ge\frac{(p+2)^\ell-p^\ell}{2}+2^{\ell+1}-4$: PC degree $\ge2^\ell$, over every field of
    characteristic $p$.
  & T-WRITTEN; inputs \cite{churchellenberg,rhoades2024,beame1998counting} & \S\ref{sec:growth} \\
E & Exact degree.  For every $k\ge1$, with $q=p^{v_p(k)}$: PC degree exactly $q$ for
    $n\ge2q-1$; for $n\ge2q-2$ conditionally on \cite[Thm.~8.4.5]{hone2025}, whose
    published proof is incomplete (\S\ref{sec:errata}).  Mechanism: for $q>D$ and
    $m\equiv m'\pmod q$ large, degree-$\le D$ refutability of
    $\mathrm{onto\text{-}FPHP}^{m}_n$ and $\mathrm{onto\text{-}FPHP}^{m'}_n$ agree.
  & T-WRITTEN; inputs \cite{harman2015} and tilting theory (App.~\ref{app:tilting}) & \S\ref{sec:periodicity} \\
F & Bands.  $q=p^{v_p(k)}$, $1\le D<\min(q,p^2)$, $n\ge2D-1$: no refutation of degree
    $\le D$.  Hence, if $v_p(k)=2$ and $n\ge2p^2-3$: PC degree exactly $p^2$.
    Case $p=3$, $D=4$: $9\mid k$, $n\ge7$, no refutation of degree $\le4$ (a second
    route through finite certificates).
  & T-WRITTEN; inputs \cite{shareshianwachs,yalcin2016,miller2024};
    $p=3$, $D=4$: Lean conditional (K100) & \S\ref{sec:bands} \\
G & Three points in Harman~\cite{harman2015} and two in Hone et
    al.~\cite{hone2025}; our proofs avoid each.
  & T-WRITTEN & \S\ref{sec:errata} \\
\bottomrule
\end{tabular}
\caption{Results.  ``Inputs'' are published theorems the written proofs use and do not
re-prove.  Where the statement says ``no refutation'' the system is unsatisfiable and has a
refutation of degree $p^{v_p(k)}$ by Riis's criterion.  Throughout, lower-bound and
exact-degree statements assume $n\ge1$ and $k\ge1$.}
\label{tab:results}
\end{table}

\paragraph{The exact degree (E, with A and F).}  Result E says that Riis's upper bound is
attained: for every $k$, the PC degree is exactly $p^{v_p(k)}$ once $n\ge2p^{v_p(k)}-1$.  In
particular, for surplus $p^\ell$ it is $p^\ell$, so the degree grows with $\ell$ as fast as
it can.  E is the strongest statement in the paper and rests on the deepest inputs: Harman's
stable equivalences, with the repairs of Section~\ref{sec:errata}, and classical tilting
theory (Appendix~\ref{app:tilting}).  A shorter proof through the semisimplification
functors of Hone et al.\ would give the threshold one lower, but it uses their
Theorem~8.4.5, and while checking that theorem we found that its published proof is
incomplete (Sections~\ref{sec:periodicity} and~\ref{sec:errata}); we present that proof as
conditional.  The other results confirm parts of E on other inputs.  At the Lean tier,
A gives the exact degree $p$ when $v_p(k)=1$, from $n=2p-3$; B gives at least $p+1$ when
$p^2\mid k$, from $n=2p-1$; C gives at least $5$ for $p=2$, $k=8$, from $n=7$.  At the
written tier, F gives the exact degree $p^2$ when $v_p(k)=2$, from $n=2p^2-3$, using neither
of E's inputs.  For comparison, Theorem~33 of \cite{riis1997nonconstant} gives
$2^\ell-1$ for $N\ge(21+3p^2)(2^\ell-2)-p$.

\paragraph{Growth from Nullstellensatz designs (D).}  D gives the weaker bound $2^\ell$ with
explicit thresholds, by a transfer from Nullstellensatz designs to PC that uses neither E's
inputs nor F's.  Its degree bound is one more than the bound $2^\ell-1$ of
\cite[Thm.~33]{riis1997nonconstant}, and its threshold in $n$ is smaller for small $\ell$ and
larger for large $\ell$ (for $p=3$: $12$ against $93$ at $\ell=2$, $1501$ against $1437$ at
$\ell=5$).  At $p=2$ its bound meets Riis's, so D gives the exact degree $2^\ell$ over
$\mathbb F_2$ independently of E, from $n\ge(4^\ell-2^\ell)/2+2^{\ell+1}-4$.  None of D, E, F
is formalised in Lean.

\paragraph{The bands (F).}  F is not a proof of growth: a bound that stops at $p^2$
cannot give unbounded growth at fixed $p$.  It gives the fixed-valuation and fixed-degree
statements above.  Its F$_3$ case ($9\mid k$, degree $\le4$, all $n\ge7$) has a Lean
counterpart that is conditional on a finite-looking spanning statement for every board
(Section~\ref{sec:bands}); that condition is verified by exact computation on the smallest
board at two of the three levels it needs.

\paragraph{Errata (G).}  Checking the sources of E led to five points in two recent
sources: in Harman's paper, a statement that is false at $p=2$, a misprinted modulus and an
omitted justification; in Hone et al., a theorem that is false at the boundary of its
hypotheses, and a lemma that is false as stated, on which the proof of their Theorem~8.4.5
depends.  We sent the first four to the authors on 8 November 2026.  Our proofs avoid each
defect.

\subsection{Evidence tiers}\label{sec:intro-tiers}

Every statement in this paper carries one of three tiers.
\begin{description}
\item[T-LEAN] The stated formal theorem is kernel-checked in Lean~4, using only the axioms
  \texttt{propext}, \texttt{Classical.choice} and \texttt{Quot.sound}.  Additional model
  re-proofs and independent certificate checks are described separately for each result
  (Section~\ref{sec:process}).
\item[T-WRITTEN] There is a written proof.  At least one AI model other than its author
  reconstructed it, and a separate AI adjudicator (Astra, a Codex model that took no part
  in the run) checked it against the primary sources.  Author and reconstructor are named
  per result; they are not always of different vendors.  Published theorems that the proof uses are named as
  inputs and are not re-proved.
\item[T-MEASURED] An exact finite computation, with a reproducible script; a second,
  independently written implementation where we say so.
\end{description}
The tiers are not equally strong.  A T-LEAN statement is as reliable as the Lean kernel
and the match between the Lean definitions and the mathematics, which
Section~\ref{sec:defs} states in full.  A T-WRITTEN statement is as reliable as its
model review and its inputs; we name, for each such result, the input we consider most
exposed.  No human expert has reviewed any proof in this paper.

The results were produced by AI research agents working in a harness that records every
claim, review and computation; Section~\ref{sec:process} describes this and what was
checked.  The Lean library and the checks are public with the four findings behind A, B, C and the
$\mathbb F_3$ target (Zenodo DOIs 10.5281/zenodo.23996564, 10.5281/zenodo.23996585,
10.5281/zenodo.23996615 and 10.5281/zenodo.24056819).

\subsection{Organisation}

Section~\ref{sec:defs} defines the two formulations of PC used here and proves that
they have the same refutations, and sets up the duality between refutations and designs.
Section~\ref{sec:lean} gives prose proofs of A, B and C.  Sections~\ref{sec:growth},
\ref{sec:periodicity} and~\ref{sec:bands} prove D, E and F.  Section~\ref{sec:errata}
records the errata, Section~\ref{sec:literature} the relation to the literature, and
Section~\ref{sec:process} how the results were produced and checked.
Section~\ref{sec:open} lists open problems.

% sec2-definitions.tex -- Section 2 of the C3 paper (draft 2026-11-08).
% The framework formerly at the top of sec3-lean-results.tex, moved here and expanded:
% full proof of the equivalence of the two PC formulations (from StdPC.lean, StdPCConv.lean),
% degree monotonicity (maintenance-20261108/DegreeMono.lean), the upper bound, and the
% matching/design machinery shared by Sections 3--6.  Labels are unchanged, so references in
% sec3 still resolve.  Requires \lean as defined at the top of sec3 (or the preamble).

\section{Definitions and the two formulations of polynomial calculus}\label{sec:defs}

\begingroup\sloppy
Lean names are given in footnotes, with namespaces abbreviated as in
Section~\ref{sec:lean}: \texttt{[0459]} stands for
\texttt{Programme.\guillemotleft 20261025-0459-proof-complexity\guillemotright}, and so on;
\texttt{[0635]} for \texttt{Programme.\guillemotleft 20261102-0635-proof-complexity\guillemotright};
\texttt{[maint]} stands for
\texttt{Programme.\guillemotleft maintenance-20261108\guillemotright}.

\subsection{The formula}\label{sec:defs-formula}

Fix $n,k\in\mathbb N$.  There are $n+k$ pigeons $0\le i<n+k$ and $n$ holes
$0\le j<n$; the variable $x_{ij}$ has index $in+j$, so the variables are
$x_0,\dots,x_{N-1}$ with $N=(n+k)n$.  The system $\mathrm{ontoFPHP}(n,k)$ is the list of
integer polynomials
\[
  \textstyle P_i=\sum_{j<n}x_{ij}-1,\quad H_j=\sum_{i<n+k}x_{ij}-1,\quad
  x_{ij}x_{i'j}\ (i<i'),\quad x_{ij}x_{ij'}\ (j<j').
\]
We call them pigeon, hole, collision and functional axioms.\lean{[0459] OntoFPHP.ontoFPHP}
They are the axioms of \cite[Def.~7.2]{beame1998counting}, with integer coefficients that
are read mod $p$ below.  Lean represents a polynomial by a list of terms (an integer coefficient and a list of
variables) and uses its raw syntactic degree, the maximum length of a term's variable list;
rule (eq) below permits collecting coefficients and removing zero coefficients.  For the
displayed axioms, whose representations are normalised, this presents ordinary polynomial
calculus over $\mathbb F_p$ with the same degree bound.  No reduction by $x^2=x$ is applied
unless we say so.

\subsection{Standard polynomial calculus}\label{sec:defs-std}

\begin{definition}[standard PC]\label{def:stdpc}
Let $p\ge2$, let $A$ be a list of integer polynomials and $D\in\mathbb N$.
$\mathrm{Der}_p(A,D)$ is the least set of integer polynomials such that
\begin{itemize}
\item[(ax)] $P\in\mathrm{Der}_p(A,D)$ for every $P\in A$ with $\deg P\le D$;
\item[(bool)] $x^2-x\in\mathrm{Der}_p(A,D)$ for every variable $x$, provided $D\ge2$;
\item[(lin)] $uP+vQ\in\mathrm{Der}_p(A,D)$ for $P,Q\in\mathrm{Der}_p(A,D)$ and $u,v\in\mathbb Z$;
\item[(mul)] $xP\in\mathrm{Der}_p(A,D)$ for $P\in\mathrm{Der}_p(A,D)$ with $\deg P+1\le D$;
\item[(eq)] $Q\in\mathrm{Der}_p(A,D)$ whenever $P\in\mathrm{Der}_p(A,D)$, $\deg Q\le D$ and
  every monomial has congruent coefficients mod $p$ in $P$ and $Q$.
\end{itemize}
A \emph{degree-$D$ refutation} of $A$ over $\mathbb F_p$ exists, written
$\mathrm{Ref}_p(A,D)$, if $1\in\mathrm{Der}_p(A,D)$.\lean{[1156] StdPC.Der, StdPC.Refutes, StdPC.PEq, StdPC.boolAx}
\end{definition}

For prime $p$ this is the usual polynomial calculus over $\mathbb F_p$ with Boolean
axioms: rule (eq) identifies an integer polynomial with its image in $\mathbb F_p[x]$, and
every line has degree at most $D$ (by induction; (lin) does not raise degree).  Two points
of convention matter for comparisons.  First, the degree is the degree of a \emph{line},
measured syntactically, before any reduction by $x^2=x$: rule (mul) is guarded by
$\deg P+1\le D$ on the unreduced polynomial $P$.  Second, Nullstellensatz degree in
\cite{beame1998counting} is the degree of the multipliers; we quote their bounds in
multiplier degree and state any conversion to product or PC line degree explicitly.

\begin{lemma}[monotonicity]\label{lem:mono}
If $D\le D'$ then $\mathrm{Der}_p(A,D)\subseteq\mathrm{Der}_p(A,D')$.  Hence
$\neg\mathrm{Ref}_p(A,D')$ implies $\neg\mathrm{Ref}_p(A,D)$ for every $D\le D'$.\lean{[maint] DegreeMono.Der.mono, Refutes.mono}
\end{lemma}

\begin{proof}
Every rule instance of degree $D$ is a rule instance of degree $D'$: each guard
$\deg\le D$ or $\deg+1\le D$ (and $D\ge2$) implies the same guard with $D'$.
\end{proof}

The parameter $D$ is an upper bound on line degree.  Monotonicity allows a lower bound
at $D$ to be restated for all smaller bounds; the library states those forms
separately.\lean{[maint] DegreeMono.uniform\_no\_degree\_le\_p, uniform\_no\_standard\_degree\_le\_p, no\_deg\_le4, deg\_ge5}

\subsection{Semantic polynomial calculus on the cube}\label{sec:defs-sem}

The Lean library proves its lower bounds for a formulation in which lines are functions on
the Boolean cube.  This removes the Boolean axioms and makes congruence a rule, which is
convenient for the duality arguments of Section~\ref{sec:defs-designs}.

A \emph{point} is a map $a\colon\mathbb N\to\{0,1\}$; a \emph{line} is a function
$g$ from points to $\mathbb Z$.  Each integer polynomial $P$ defines a line $\hat P$ by
evaluation.  A line $g$ has \emph{cube degree at most $d$} if $\hat Q(a)\equiv g(a)\pmod p$
for all points $a$, for some integer polynomial $Q$ with $\deg Q\le d$.\lean{[1633] PCTseitin.CubeDeg}

\begin{definition}[semantic PC]\label{def:ppc}
The lines of a degree-$D$ derivation over $\mathbb Z/p$ from $A$ form the least set
$\mathrm{PPC}_p(A,D)$ of lines containing
(ax) $\hat P$ for each $P\in A$ with $\deg P\le D$;
(lin) $uf+vg$ for $f,g\in\mathrm{PPC}_p(A,D)$ and $u,v\in\mathbb Z$;
(mul) $x\cdot g$ for a variable $x$ and $g\in\mathrm{PPC}_p(A,D)$ of cube degree at most
$d$ with $d+1\le D$;
(congr) every $f$ with $f(a)\equiv g(a)\pmod p$ for all points $a$, for some
$g\in\mathrm{PPC}_p(A,D)$.
A \emph{degree-$D$ refutation} is a derivation of the constant line $1$; we write
$\mathrm{PPCRef}_p(A,D)$.\lean{[0459] OntoFPHP.PPCDerivable, OntoFPHP.PPCRefutation}
\end{definition}

Here $p$ is any integer $\ge2$, not necessarily prime; all our results take $p$ prime.
Degree monotonicity holds as in Lemma~\ref{lem:mono}.\lean{[maint] DegreeMono.PPCDerivable.mono, PPCRefutation.mono}

\begin{theorem}\label{thm:refutes-iff}
For every $p\ge2$, list $A$ and $D$: $\mathrm{Ref}_p(A,D)\iff\mathrm{PPCRef}_p(A,D)$.\lean{[1156] StdPCConv.refutes\_iff}
\end{theorem}

The two directions are Propositions~\ref{prop:std-to-sem} and~\ref{prop:sem-to-std}.
Consequently every lower bound in this paper holds verbatim for standard PC over
$\mathbb F_p$ with Boolean axioms.

\begin{proposition}\label{prop:std-to-sem}
If $p\ge2$ and $\mathrm{Ref}_p(A,D)$, then $\mathrm{PPCRef}_p(A,D)$.\lean{[1156] StdPC.refutes\_ppc}
\end{proposition}

\begin{proof}
We show by induction on $\mathrm{Der}_p(A,D)$ that for every $P\in\mathrm{Der}_p(A,D)$
either $\hat P(a)\equiv0\pmod p$ at every point $a$, or
$\hat P\in\mathrm{PPC}_p(A,D)$.\lean{[1156] StdPC.der\_sem}  The first alternative is
needed because the Boolean axioms have no semantic counterpart: their value is $0$, and
the zero line need not be in $\mathrm{PPC}_p(A,D)$ (it is empty when $A$ has no axiom of
degree $\le D$).
(ax) gives the second alternative and (bool) the first.  (lin): if both $\hat P,\hat Q$
are derivable, so is $u\hat P+v\hat Q$; if one vanishes, $u\hat P+v\hat Q$ is congruent
to a multiple of the other and is derivable by (lin) and (congr), or vanishes.  (eq):
polynomials with coefficients congruent mod $p$ take congruent values, so (congr)
applies.  (mul): $\deg P+1\le D$ means $\hat P$ has cube degree at most $D-1$, so semantic
(mul) applies to $\hat P$, and $\widehat{xP}=x\cdot\hat P$; if $\hat P$ vanishes, so does
$x\cdot\hat P$.

Now let $1\in\mathrm{Der}_p(A,D)$.  The first case is impossible: $1\not\equiv0\pmod p$ at
any point since $p\ge2$.  So $\hat1=1\in\mathrm{PPC}_p(A,D)$.
\end{proof}

For the converse the only difficulty is (mul): a semantic line $g$ of cube degree $\le D-1$
may be represented in the standard derivation by a polynomial $Q$ of degree $D$, which may
not be multiplied.  We must derive a low-degree representative $R$ of $g$ instead.  Two
lemmas do this.  For a monomial $m$ let $\mathrm{ml}(m)$ be its multilinear form (the set of
its variables), and extend $\mathrm{ml}$ linearly to polynomials.

\begin{lemma}[multilinear reduction]\label{lem:ml-reduce}
If $\mathrm{Der}_p(A,D)$ contains some polynomial and $m$ is a monomial with $|m|\le D$,
then $m-\mathrm{ml}(m)\in\mathrm{Der}_p(A,D)$.\lean{[1156] StdPCConv.der\_dd}
\end{lemma}

\begin{proof}
If $m$ is multilinear, $m-\mathrm{ml}(m)$ has all coefficients $0$: derive zero as
$0\cdot Q+0\cdot Q$ by (lin) from any derived line $Q$, and apply (eq).  Otherwise some variable $x$ occurs twice, $m=x^2u$ with
$|u|=|m|-2$, and in particular $D\ge|m|\ge2$.  Then $m-xu=u\cdot(x^2-x)$, obtained from the
Boolean axiom by $|u|$ applications of (mul); the intermediate products have degree
$\le|m|\le D$, so every guard holds.  Since $\mathrm{ml}(xu)=\mathrm{ml}(m)$ and $|xu|<|m|$,
induction on $|m|$ gives $xu-\mathrm{ml}(m)$, and (lin) adds the two.
\end{proof}

\begin{lemma}[multilinear vanishing]\label{lem:ml-vanish}
If a multilinear integer polynomial $S=\sum_T c_Tx_T$ satisfies $\hat S(a)\equiv0\pmod p$
at every point $a$, then $c_T\equiv0\pmod p$ for every $T$.\lean{[1156] StdPCConv.cf\_ml\_dvd}
\end{lemma}

\begin{proof}
By induction on $|T|$.  Evaluate at the indicator point $a_T$ of $T$:
$\hat S(a_T)=\sum_{T'\subseteq T}c_{T'}$.  By induction $c_{T'}\equiv0$ for $T'\subsetneq T$,
so $c_T\equiv\hat S(a_T)\equiv0$.
\end{proof}

\begin{lemma}\label{lem:der-of-cube}
Let $Q\in\mathrm{Der}_p(A,D)$, and let $R$ be an integer polynomial with $\deg R\le D$ and
$\hat R(a)\equiv\hat Q(a)\pmod p$ at every point.  Then $R\in\mathrm{Der}_p(A,D)$.\lean{[1156] StdPCConv.der\_of\_cube, der\_of\_zero, der\_of\_ml}
\end{lemma}

\begin{proof}
Put $S=R-Q$, a polynomial of degree $\le D$ that vanishes mod $p$ on the cube.  Its
multilinear form $\mathrm{ml}(S)$ takes the same values, so by
Lemma~\ref{lem:ml-vanish} all its coefficients are $\equiv0$; by (eq) applied to $0\cdot Q$,
$\mathrm{ml}(S)\in\mathrm{Der}_p(A,D)$.  Adding the multiples $c\,(m-\mathrm{ml}(m))$ of
Lemma~\ref{lem:ml-reduce} for the terms $cm$ of $S$ gives $S$.  Then $R=S+Q$ by (lin).
\end{proof}

\begin{proposition}\label{prop:sem-to-std}
If $\mathrm{PPCRef}_p(A,D)$, then $\mathrm{Ref}_p(A,D)$.\lean{[1156] StdPCConv.ppc\_refutes, ppc\_std}
\end{proposition}

\begin{proof}
By induction on $\mathrm{PPC}_p(A,D)$, every semantic line $g$ has some
$Q\in\mathrm{Der}_p(A,D)$ with $\hat Q\equiv g$ pointwise mod $p$.  (ax), (lin) and (congr)
are immediate.  For (mul), let $g$ have cube degree $\le D-1$, witnessed by $R$ with
$\deg R\le D-1$, and let $Q$ be the derived representative of $g$.  By
Lemma~\ref{lem:der-of-cube}, $R\in\mathrm{Der}_p(A,D)$, so $xR\in\mathrm{Der}_p(A,D)$ by
(mul), and $\widehat{xR}\equiv x\cdot g$.  Finally, if $1$ is a semantic line, its
representative $Q$ agrees with the constant polynomial $1$ on the cube, and
Lemma~\ref{lem:der-of-cube} gives $1\in\mathrm{Der}_p(A,D)$.
\end{proof}

\subsection{The upper bound}\label{sec:defs-upper}

\begin{proposition}[Riis~\cite{riis1993independence}; Theorem~8 of \cite{razborov2001php}]\label{prop:riis}
Let $p$ be prime.  If $\binom{n+k}{d}\not\equiv\binom nd\pmod p$ for some $d\ge1$, then
$\mathrm{onto\text{-}FPHP}^{n+k}_n$ has a PC refutation of degree $d$ over $\mathbb F_p$.  For
$k=p^vc$ with $p\nmid c$, the choice $d=p^v$ applies by Lucas's theorem, giving degree
$p^{v}$ for every $n$.
\end{proposition}

We use this only as an upper bound and do not re-prove it in general.  In the case needed
for Theorem~A ($v=1$) the library proves the degree-$p$ refutation directly
(Proposition~\ref{prop:A-upper}).  The criterion cannot give less than $p^{v}$: for
$d<p^v$, Lucas's theorem gives $\binom{n+k}{d}\equiv\binom nd$.

\subsection{Matchings, balance and extension}\label{sec:defs-matchings}

A \emph{matching} is a set of variables no two of which share a pigeon or a hole;
Lean represents it by its decreasing list (\texttt{CM}).  For
$f\colon\{\text{matchings}\}\to\mathbb Z$, a matching $g$ and a pigeon $i$ free in
$g$, the \emph{row sum} is $\mathrm{row}_f(g,i)=\sum_j f(g\cup\{x_{ij}\})$ over the
holes $j$ for which $g\cup\{x_{ij}\}$ is a matching; the \emph{column sum}
$\mathrm{col}_f(g,j)$ at a free hole $j$ is defined likewise over pigeons.
Write $|g|$ for the number of edges.

\begin{definition}\label{def:bal}
$\mathrm{Bal}_p(f,c,e)$ holds if for every matching $g$ with $|g|=e$, every free
row sum and every free column sum of $f$ at $g$ is $\equiv c(g)\pmod p$.  For
$p=0$ this is exact equality.  We call $T$ a \emph{window up to level $d$} if
$\mathrm{Bal}_p(T,T,e)$ for all $e\le d$, and write $\mathcal W_d$ for the set of
these.\lean{[1402] ExtLemma.Bal, rowS, colS; [0607] Peel.Wd}
\end{definition}

The parameters $n,k$ are implicit in $\mathrm{Bal}$; we write
$\mathrm{Bal}_p^{(n,k)}$ when they vary.

\begin{definition}\label{def:ext}
$\mathrm{EXT}_p(n,k,d)$: for all $f,c$, if $d=0$ or $\mathrm{Bal}_p(f,c,d-1)$, then
there is $y$ with $\mathrm{Bal}_p(y,f,d)$.\lean{[1402] ExtLemma.EXT}
\end{definition}

For a set $G$ of functions, $\mathrm{Span}_e(G)$ is the least set containing $G$
and closed under sums, integer multiples, and replacement of $f$ by any $f'$ with
$f'(g)\equiv f(g)\pmod p$ for all matchings $g$ with
$|g|=e$.\lean{[1402] ExtLin.InSpan}

\begin{lemma}\label{lem:ext-span}
$\mathrm{Bal}_p(y,f,e)$ is linear in the pair $(y,f)$ and depends on $f$ only
through its residues on matchings of size $e$.  Hence, if every $f\in G$ has some
$y$ with $\mathrm{Bal}_p(y,f,e)$, so does every $f\in\mathrm{Span}_e(G)$.  In
particular, if every $f$ with $\mathrm{Bal}_p(f,c,d-1)$ (any $c$) lies in
$\mathrm{Span}_d(G)$, then $\mathrm{EXT}_p(n,k,d)$.\lean{[1402] ExtLin.Bal\_add, Bal\_smul, Bal\_congr, ext\_of\_span, EXT\_of\_span}
\end{lemma}

\begin{proof}
Row and column sums are linear in $f$; the three closure rules of
$\mathrm{Span}_e$ are matched by adding extensions, scaling them, and keeping the
same $y$.
\end{proof}

\begin{lemma}[flag lemma]\label{lem:flag}
If $\mathrm{Bal}_p(f,c,e+1)$, then at every matching $g$ with $|g|=e$, every free
row sum of $c$ is congruent to every free column sum of $c$.  Consequently there
is $c'$ with $\mathrm{Bal}_p(c,c',e)$.\lean{[0607] Flag.row\_col\_gen, Flag.flag; also [1434] RowColG.link\_bal}
\end{lemma}

\begin{proof}
Let $i$ be a free pigeon and $b$ a free hole of $g$.  For a free hole $j\ne b$,
$c(g+x_{ij})\equiv\mathrm{col}_f(g+x_{ij},b)=\sum_{i'\ne i}f(g+x_{ij}+x_{i'b})$, so
$\mathrm{row}_c(g,i)\equiv c(g+x_{ib})+\sum_{j\ne b}\sum_{i'\ne i}f(g+x_{ij}+x_{i'b})$.
Expanding $\mathrm{col}_c(g,b)$ through the row sums of $f$ at $g+x_{i'b}$ gives
$c(g+x_{ib})$ plus the same double sum; the sums run over free $j$ and free
$i'$.  Put $c'(g)$ equal to the column sum of $c$ at a chosen free hole of $g$
($0$ if none).
\end{proof}

\subsection{Designs and the extension criterion}\label{sec:defs-designs}

Lower bounds are proved by exhibiting dual objects, \emph{designs}, which vanish on every
line of a low-degree derivation but not on $1$.

\begin{definition}\label{def:design}
Let $N=(n+k)n$.  A \emph{functional} $E$ assigns an integer to each line, is
linear mod $p$, respects congruence mod $p$, and reads its argument only at
points whose variables of index $\ge N$ are $0$.  It is a \emph{degree-$D$ design}
for $A$ if $E(x_m\cdot P)\equiv0\pmod p$ for every $P\in A$ and every list $m$ of
variables with $|m|+\deg P\le D$, where $x_m=\prod_{v\in m}x_v$.  $E$
\emph{extends} $E'$ to degree $d$ if $E(x_m)\equiv E'(x_m)$ for all $|m|\le
d$.\lean{[1627] DesignExt.Fn, IsDesign, Extends}
\end{definition}

\begin{lemma}\label{lem:ext-sound}
Let $D\ge1$.  Suppose every degree-$(D-1)$ design has an extension to degree $D-1$
that is a degree-$D$ design, and some degree-$(D-1)$ design $E_0$ has
$E_0(1)\not\equiv0$.  Then there is no degree-$D$ refutation of
$A$.\lean{[1627] DesignExt.ext\_sound}
\end{lemma}

\begin{proof}
We show that every degree-$D$ design $E$ kills every line $g$ of a degree-$D$
derivation, by induction on the derivation.  (ax) $E(1\cdot P)\equiv0$ since
$\deg P\le D$.  (lin), (congr) by linearity and congruence.  (mul) Let $g$ have
cube degree $\le D-1$.  The functional $h\mapsto E(x\cdot h)$ is a degree-$(D-1)$
design;\lean{[1627] DesignExt.xFn, xFn\_design} by hypothesis it has an extension
$E_3$ to degree $D-1$ that is a degree-$D$ design.  By induction $E_3(g)\equiv0$;
two functionals that agree on monomials of length $\le D-1$ agree on lines of
cube degree $\le D-1$,\lean{[1627] DesignExt.agree} so
$E(x\cdot g)\equiv E_3(g)\equiv0$.  Finally extend $E_0$ to a degree-$D$ design
$E$: then $E(1)\equiv E_0(1)\not\equiv0$, so $1$ is not a line.
\end{proof}

The hypothesis of Lemma~\ref{lem:ext-sound} is stronger than the existence of one
degree-$D$ design with $E(1)\not\equiv0$, which is the Nullstellensatz-type condition.
The extra requirement, that \emph{every} lower-degree design extend, is what makes the
lemma sound for PC, where multiplication may follow linear combination.

\begin{lemma}[extension step]\label{lem:ext-step}
Let $n>0$, assume $\mathrm{EXT}_p(n,k,d)$, and let $F$ be a degree-$d$ design for
$\mathrm{ontoFPHP}(n,k)$.  Then some degree-$(d+1)$ design extends $F$ to degree
$d$.\lean{[1402] ExtLemma.ext\_step}
\end{lemma}

\begin{proof}
Put $f(t)=F(x_t)$ for matchings $t$.  (1) A degree-$d$ design kills every monomial
of length $\le d$ that contains two distinct variables sharing a pigeon or a hole, since
such a monomial is a multiple of a collision or functional
axiom.\lean{[1402] ExtLemma.conf\_ax, kill}
(2) $\mathrm{Bal}_p(f,f,d-1)$ when $d\ge1$: for a matching $g$ with $|g|=d-1$ and a
free pigeon $i$, $F(x_g P_i)\equiv\sum_j F(x_{g}x_{ij})-F(x_g)$, and the terms with
$g\cup\{x_{ij}\}$ not a matching vanish by (1), so $F(x_gP_i)\equiv
\mathrm{row}_f(g,i)-f(g)$; the design property gives $F(x_gP_i)\equiv0$.  Columns
likewise.\lean{[1402] ExtLemma.rowId, termsum, line\_bal}
(3) $\mathrm{EXT}_p(n,k,d)$ yields $y$ with $\mathrm{Bal}_p(y,f,d)$.  Let $E$ be the
moment functional whose value on $x_S$ is $F(x_S)$ for $|S|\le d$, $y(S)$ for
matchings with $|S|=d+1$, and $0$ for other sets with $|S|=d+1$.\lean{[1402] ExtLemma.W, Ext; [1627] MomGen.sumWFn}
It extends $F$ to degree $d$.\lean{[1402] ExtLemma.Ext\_extends}
(4) $E$ is a degree-$(d+1)$ design.  For $|m|+\deg P\le d$ this follows from $F$
being a design and the agreement.\lean{[1627] DesignExt.low\_kill}
For a collision or functional axiom with $|m|=d-1$, $x_mP$ is a monomial of length
$\le d+1$ with a conflicting pair: its value is $0$ (non-matching at the top
level) or killed by (1).\lean{[1402] ExtLemma.Ext\_kill}
For a pigeon axiom with $|m|=d$ (holes are symmetric), let $s$ be the set of $m$.  If
$|s|<d$, then $x_m=x_s$ on the cube and $x_sP_i$ has degree at most $d$, so
$E(x_mP_i)\equiv F(x_sP_i)\equiv0$; we may therefore assume $|s|=d$.
If $s$ is not a matching, every term of $E(x_s P_i)$ contains a conflict.  If some
$x_{ij}\in s$, all terms but one conflict and the remaining one cancels $E(x_s)$.
Otherwise the line is free and $E(x_sP_i)\equiv \mathrm{row}_y(s,i)-f(s)\equiv0$ by
$\mathrm{Bal}_p(y,f,d)$.\lean{[1402] ExtLemma.Ext\_line, line\_conf, line\_occ}
\end{proof}

\begin{theorem}[EXT criterion]\label{thm:ext-criterion}
Let $p>1$, $n_0>0$, $D\ge1$, and suppose $\mathrm{EXT}_p(n_0,k,d)$ for every $d<D$.
Then for every $n\ge n_0$ there is no degree-$D$ refutation of
$\mathrm{ontoFPHP}(n,k)$ over $\mathbb Z/p$.  No hypothesis on $k$ is
needed.\lean{[1402] ExtLemma.no\_refutation\_of\_ext, no\_refutation\_of\_ext\_ge}
\end{theorem}

\begin{proof}
At $n=n_0$: the functional with moment $1$ at $\emptyset$ is a degree-$0$ design
(vacuously, every axiom has degree $\ge1$) with value $1$ at $1$.  Lemma
\ref{lem:ext-step} applied $D-1$ times gives a degree-$(D-1)$ design $E_0$ with
$E_0(1)\equiv1$, and once more shows that every degree-$(D-1)$ design
extends.\lean{[1402] ExtLemma.designs}  Lemma~\ref{lem:ext-sound} applies.
For $n>n_0$ use restriction: substituting $x_{n+k,n}:=1$ and $0$ for the other
variables of pigeon $n+k$ and of hole $n$ maps every axiom of
$\mathrm{ontoFPHP}(n+1,k)$ to an axiom of $\mathrm{ontoFPHP}(n,k)$ or to a line
derivable from them, so a degree-$D$ refutation for $n+1$ yields one for
$n$.\lean{[0459] OntoSubst.ontoFPHP\_restrict, PPCDerivable.subst}
\end{proof}

The restriction step also shows that refutability is monotone in $n$ for fixed $k$ and
$D$: lower bounds proved at one $n_0$ hold for all larger $n$.

\paragraph{Units.}  For prime $p$ every $1\le s\le p-1$ is invertible mod $p$
(B\'ezout), hence so is $s!$ for $s\le p-1$.\lean{[0607] Peel.units\_of\_prime, fact\_unit}
Lean's \texttt{IsPrime p} is: $p\ge2$ and every divisor of $p$ is $1$ or
$p$.\lean{[0607] Peel.IsPrime}

\subsection{The matching algebra}\label{sec:defs-written}

Sections~\ref{sec:growth}--\ref{sec:bands} work over an arbitrary field $F$ and use a
linear-algebra form of the same duality.  We fix a board with a set $P$ of $m=n+k$ pigeons
and a set $S$ of $n$ holes, and write $\mathrm{PC}_F$ for Definition~\ref{def:stdpc} with
coefficients in $F$ in place of integers read mod $p$.

\paragraph{Reduction.}  Let $\mathcal B=F[x]/(x_{ij}^2-x_{ij},\ \text{collision and
functional axioms})$ be the \emph{Boolean matching algebra}, and $\rho\colon F[x]\to\mathcal B$
the quotient map.  The monomials $x_M=\prod_{(i,j)\in M}x_{ij}$, $M$ a matching, form a
basis of $\mathcal B$; $\rho$ sends a monomial to $x_M$ if its set of variables is a matching
$M$ and to $0$ otherwise.  Since $\rho$ is a ring homomorphism, $\rho(xf)=\rho(x)\rho(f)$.
Let $F_D\subseteq\mathcal B$ be the span of the $x_M$ with $|M|\le D$; then
$\rho(f)\in F_D$ whenever $\deg f\le D$.

\paragraph{The relation spaces.}  For a matching $N$ and a pigeon $i$ free in $N$, the
\emph{row relation} at $(N,i)$ is
$\rho(x_N\,(\sum_j x_{ij}-1))=\sum_{j}x_{N\cup\{(i,j)\}}-x_N$, the sum over holes $j$
free in $N$; column relations at $(N,j)$ are defined likewise.  (If $i$ is not free in
$N$, the reduced product is $0$.)  Let
\[
  W_D=\operatorname{span}\{\text{row and column relations at bases $N$ with } |N|\le D-1\}
  \subseteq F_D,
\]
so $W_0=0$ and $W_{D-1}\subseteq W_D\cap F_{D-1}$.  We use two conditions:
\[
  \mathrm{U}_D:\ 1\notin W_D,\qquad
  \mathrm{ND}_D:\ W_D\cap F_{D-1}=W_{D-1}.
\]

\paragraph{Designs.}  A \emph{unit degree-$D$ design} on the board is a linear functional
$L$ on $F_D$ with $L(W_D)=0$ and $L(1)=1$.  In coordinates it is a function on matchings of
size at most $D$ with value $1$ at $\emptyset$ such that, at every base $N$ with $|N|<D$,
every free row sum and every free column sum equals the value at $N$; this is a
$D$-design in the sense of Beame and Riis~\cite[Def.~8.2]{beame1998counting}, with values
in $F$.  By finite-dimensional duality, $\mathrm{U}_D$ holds if and only if a unit
degree-$D$ design exists, and $\mathrm{ND}_D$ holds if and only if every linear functional
on $F_{D-1}$ that vanishes on $W_{D-1}$ extends to one on $F_D$ that vanishes on $W_D$.
(If every such functional extends, each of them vanishes on $W_D\cap F_{D-1}$, so
$W_D\cap F_{D-1}\subseteq W_{D-1}^{\perp\perp}=W_{D-1}$; the converse is the same
argument read backwards.)

$\mathrm{U}_D$ is the Nullstellensatz condition: reducing a Nullstellensatz certificate
$\sum_r q_rA_r=1$ of product degree at most $D$ under $\rho$ kills the Boolean, collision
and functional terms and turns each row or column term into a combination of relations
at bases of size at most $D-1$, so the certificate exists if and only if $1\in W_D$.  This
is the product-degree convention; it differs from the multiplier degree of
\cite{beame1998counting}, and we convert between them only where stated.

\begin{proposition}[from $\mathrm{ND}_D$ and $\mathrm{U}_D$ to PC]\label{prop:nd-pc}
Let $D\ge1$.  If $\mathrm{ND}_D$ and $\mathrm{U}_D$ hold on a board, then
$\mathrm{onto\text{-}FPHP}$ on that board has no $\mathrm{PC}_F$ refutation of degree at
most $D$.
\end{proposition}

\begin{proof}
We show that $\rho(f)\in W_D$ for every line $f$ of a degree-$D$ derivation, by induction.
Pigeon and hole axioms reduce to the relations at $N=\emptyset$, which lie in $W_D$ since
$D\ge1$; Boolean, collision and functional axioms reduce to $0$.  Linear combinations, and
replacement of coefficients by equal ones in $F$, preserve membership.  For
multiplication, the guard of rule (mul) gives $\deg f\le D-1$ for the unreduced line, so
$\rho(f)\in W_D\cap F_{D-1}=W_{D-1}$ by $\mathrm{ND}_D$.  Each relation at a base $N$ with
$|N|\le D-2$, multiplied by a variable $x$ and reduced, is $0$ or a relation at the base
$N\cup\{x\}$ (or at $N$ itself, if $x\in N$), of size at most $D-1$; so
$\rho(xf)=\rho(x)\rho(f)\in W_D$.  Since $1\notin W_D$ by $\mathrm{U}_D$, the line $1$ is
not derived.
\end{proof}

The induction is checked in Lean in abstract form: given any predicate on polynomials with
the closure properties used above, no refutation exists.\lean{[0635] StepSix.Invariant, StepSix.lines\_mem, StepSix.no\_refutation}
Its instantiation with $\rho^{-1}(W_D)$ is the argument above and is not formalised.  We use a
representative of degree at most $D-1$ before multiplying; the bounded normalisation
argument of Lemma~\ref{lem:der-of-cube} justifies passing to such a representative.  Checking
only the reduced degree of the product would not justify an ordinary PC multiplication
step.

\paragraph{Balanced arrays.}  For $e\ge1$ let $H_e(P,S)$ be the space of functions $f$ on
the $e$-matchings of the board all of whose free row and column sums vanish: for every
$(e-1)$-matching $N$ and free pigeon $i$, $\sum_j f(N\cup\{(i,j)\})=0$, and likewise for
free holes.  If a functional $L$ on $F_t$ vanishes on $W_t$ and on all matchings of size
$<e$, with $1\le e\le t$, then its layer $f(M)=L(x_M)$, $|M|=e$, lies in $H_e(P,S)$: the
relation at a base of size $e-1\le t-1$ lies in $W_t$ and says that the free row (or column)
sum equals the value at the base, which is $0$.

\endgroup

% sec3-lean-results.tex -- Section 3 of the C3 paper (draft 2026-11-08).
% Human-readable proofs of results A, B, C, written from the Lean 4 sources.
% Requires in the preamble: amsmath, amssymb, amsthm with environments
% theorem, proposition, lemma, corollary, definition, remark.
\providecommand{\lean}[1]{\footnote{\raggedright Lean: \texttt{#1}.}}

\section{Results proved in Lean}\label{sec:lean}

\begingroup\sloppy % local to this section: long Lean names
This section reconstructs the programme's Lean proofs of A, B and C in
mathematical prose, with the unreconstructed dependencies marked.  The proofs are Lean~4
developments in the programme library (directory
\texttt{programmes/\allowbreak proof-complexity/\allowbreak library}).  The prose follows the Lean
proofs lemma by lemma; the footnote at each statement names the Lean
declaration that proves it.  Namespaces are abbreviated by the run they come
from: \texttt{[0459]} stands for
\begin{center}
\texttt{Programme.\guillemotleft 20261025-0459-proof-complexity\guillemotright},
\end{center}
and likewise \texttt{[1627]}, \texttt{[1402]}, \texttt{[0607]}, \texttt{[1434]},
\texttt{[0841]}, \texttt{[1156]}, \texttt{[0634]}, \texttt{[1249]} and \texttt{[1633]} for
the runs 20261025-1627, 20261026-1402, 20261027-0607, 20261027-1434, 20261028-0841,
20261029-1156, 20261030-0634, 20261102-1249 and 20261024-1633.  Names under \texttt{[1249]} live in
the sub-namespace \texttt{UniformC3} (and \texttt{UniformC3.Block}).

Where a Lean step is a finite computation evaluated by the kernel
(\texttt{decide}, \texttt{decide +kernel}) or a case split closed by automation
(\texttt{simp}, \texttt{omega}, \texttt{grind}), we say so and do not supply a
conceptual argument in its place.  Lemmas from library files that we cite without reconstructing their proofs are marked
``(stated, not reconstructed: \ldots)'' with their Lean names.

\paragraph{Framework.}
The formula $\mathrm{ontoFPHP}(n,k)$, the two formulations of PC and their equivalence
(Theorem~\ref{thm:refutes-iff}), degree monotonicity (Lemma~\ref{lem:mono}), balance and
extension (Definitions~\ref{def:bal}, \ref{def:ext}), designs and the EXT criterion
(Theorem~\ref{thm:ext-criterion}) and the flag lemma (Lemma~\ref{lem:flag}) are in
Section~\ref{sec:defs}.  All three proofs below go through the EXT criterion or the
peeling of Section~\ref{sec:lean-A}.

%======================================================================
\subsection{A. Exact degree $p$ when $p$ divides the surplus exactly once}
\label{sec:lean-A}

\begin{theorem}[A]\label{thm:A}
Let $p$ be prime, $c\in\mathbb N$ with $p\nmid c$, and $n\ge 2p-3$.  For every
$D\in\mathbb N$, $\mathrm{ontoFPHP}(n,pc)$ has a degree-$D$ semantic PC refutation
over $\mathbb Z/p$ if and only if $p\le D$.\lean{[0607] Lucas.pc\_degree\_exact}
\end{theorem}

Here $2p-3$ is natural-number subtraction (for $p=2$ the threshold is $n\ge1$).
By Theorem~\ref{thm:refutes-iff} the same holds verbatim for standard PC.\lean{[1156] StdPCExact.std\_pc\_degree\_exact}
Theorem~A combines an upper bound (Proposition~\ref{prop:A-upper},
Lemma~\ref{lem:A-binom}) with the lower bound:

\begin{theorem}[A, lower bound]\label{thm:A-lower}
Let $p$ be prime, $p\mid k$ and $n\ge2p-3$.  Then $\mathrm{ontoFPHP}(n,k)$ has no
degree-$(p-1)$ refutation over $\mathbb Z/p$.\lean{[0607] Peel.pc\_degree\_ge\_p}
\end{theorem}

\begin{corollary}[separation]\label{cor:A-sep}
Let $p,q$ be primes, $p\nmid c$, $q\nmid pc$, $n\ge2p-3$.  Then
$\mathrm{PPCRef}_p(\mathrm{ontoFPHP}(n,pc),D)\iff p\le D$ for all $D$, and
$\mathrm{PPCRef}_q(\mathrm{ontoFPHP}(n,pc),1)$.\lean{[0607] Lucas.separation}
\end{corollary}

\begin{proof}
The first part is Theorem~A.  For the second, $pc>0$ and there is $u$ with
$upc\equiv1\pmod q$ (B\'ezout).  The degree-$1$ line
$-u\sum_iP_i+u\sum_jH_j=-u\bigl(S-(n+pc)\bigr)+u(S-n)=upc$, where
$S=\sum_{i,j}x_{ij}$, is $\equiv1$.\lean{[0459] OntoFPHP.ontoFPHP\_deg1; [0607] Lucas.unit\_of\_ndvd}
\end{proof}

\begin{proof}[Proof of Theorem~A]
If $D<p$ and a degree-$D$ refutation existed, it would be a degree-$(p-1)$
refutation (Lemma~\ref{lem:mono}), contradicting Theorem~\ref{thm:A-lower} (note $p\mid pc$).
If $D\ge p$, Proposition~\ref{prop:A-upper} with $t=p$ and Lemma~\ref{lem:A-binom}
give a degree-$p$ refutation, which is one of degree $D$.
\end{proof}

\subsubsection*{A.1 The upper bound}

Let $r_i=\sum_{j<n}x_{ij}$ (pigeon $i$) and $c_j=\sum_{i<n+k}x_{ij}$ (hole $j$),
and let $e_t(f_0,\dots,f_{s-1})$ be the elementary symmetric polynomial.

\begin{lemma}\label{lem:esym}
Let $D\ge1$ and assume the zero line is derivable in degree $D$.  For $s\le n+k$ and $t\le D$, the line
$e_t(r_0,\dots,r_{s-1})-\binom st$ is derivable in degree $D$ from the pigeon axioms;
likewise $e_t(c_0,\dots,c_{s-1})-\binom st$ for $s\le n$ from the hole
axioms.\lean{[0459] OntoLucas.esym\_der}
\end{lemma}

\begin{proof}
For $s=0$ or $t=0$ the line is zero.  Otherwise induction on $s$, using
\[
e^{(s+1)}_t-\tbinom{s+1}t=\bigl(e^{(s)}_t-\tbinom st\bigr)
+r_s\bigl(e^{(s)}_{t-1}-\tbinom s{t-1}\bigr)+\tbinom s{t-1}(r_s-1).
\]
The middle term is obtained by multiplying the line
$e^{(s)}_{t-1}-\binom s{t-1}$, of cube degree $t-1\le D-1$, by each variable
$x_{sj}$ (rule (mul)) and summing; the last term is a multiple of the axiom
$P_s$.
\end{proof}

\begin{lemma}\label{lem:quot}
Let every quadratic monomial axiom $x_ux_w$ of $A$ have $u\ne w$, and call a Boolean
point \emph{good} if it satisfies all of them.  If the zero line is derivable in
degree $D$, then every line of cube degree $\le D$ that vanishes mod $p$ at all good
points is derivable in degree $D$.\lean{[0459] MatchQuot.quot}
\end{lemma}

\begin{proof}[Proof (as in the source)]
Write $Q=Q|_{x_v=0}+x_v\,\partial_vQ$ variable by variable, carrying a monomial
prefix $S$ of variables already set to $1$.  When no variable is left, the line is
$x_S$ times a constant.  If $S$ contains a forbidden pair, $x_S$ is a monomial axiom
times $|S|-2$ variables; otherwise the indicator point of $S$ is good, so the
constant is $\equiv0$.  
In detail, for a list $S$ of variables write $x_S=\prod_{s\in S}x_s$ (a line), and for a
polynomial $Q$ and a variable $v$ let $Q|_{v=0}$ consist of the terms of $Q$ not containing
$v$, and $\partial_vQ$ of the terms containing $v$ with every occurrence of $v$ deleted.  On
the cube $\hat Q=\widehat{Q|_{v=0}}+x_v\,\widehat{\partial_vQ}$; moreover
$\deg Q|_{v=0}\le\deg Q$, $\deg\partial_vQ\le d$ whenever $\deg Q\le d+1$, and
$\widehat{Q|_{v=0}}(a)=\hat Q(a[v:=0])$.\lean{[0459] MatchQuot.pz, pd, eval\_split, deg\_pz, deg\_pd, eval\_pz}
The Lean proof establishes, by induction on a duplicate-free list $V$ of variables:

\smallskip\noindent\emph{Claim.}  Let $Q$ have all its variables in $V$, let $S$ be a list of
variables disjoint from $V$, and let $d$ satisfy $\deg Q\le d$ and $|S|+d\le D$.  If
$\hat Q(a)\equiv0\pmod p$ at every good point $a$ with $a_s=1$ for all $s\in S$, then
$x_S\cdot\hat Q$ is derivable in degree $D$.\lean{[0459] MatchQuot.quot\_aux}

\smallskip
The lemma is the case $S=\emptyset$, $d=D$, $V=\{0,1,\dots,m\}$ ($m$ the largest variable of
$Q$), where $Q$, $\deg Q\le D$, witnesses the cube degree of the given line $g$; $Q$ vanishes
mod $p$ at good points because $g$ does, and (congr) replaces $\hat Q$ by
$g$.\lean{[0459] MatchQuot.quot}
The only applications of (mul) are made in the following auxiliary fact: \emph{for an
axiom $P\in A$ and a list $m$ of variables with $\deg P+|m|\le D$, the line $x_m\hat P$ is
derivable.}\lean{[0459] MatchQuot.mulMono}  Indeed, for $m=\emptyset$ this is (ax), as
$\deg P\le D$; and $x_{x::m}\hat P=x\cdot(x_m\hat P)$, where $x_m\hat P$ has cube degree at
most $\deg P+|m|$ (witness: $P$ with $m$ prepended to each monomial) and
$\deg P+|m|+1\le D$, so (mul) applies to the line given by induction.  Every intermediate
line $x_{m'}\hat P$ ($m'$ a suffix of $m$) thus has cube degree at most
$\deg P+|m'|<D$ when it is multiplied.

\emph{Base case} $V=\emptyset$.  Then $\hat Q$ is a constant $c$.  (i) If $S$ contains $u,w$
with $x_ux_w\in A$, then $u\ne w$ by hypothesis; removing one occurrence of each from $S$
leaves $S'$ with $|S'|=|S|-2$, and $x_S=x_{S'}\,x_ux_w$ on the cube.  The auxiliary fact with
$P=x_ux_w$, $m=S'$ applies because $2+|S'|=|S|\le|S|+d\le D$; then
$c\cdot x_{S'}x_ux_w+0\cdot0$ (rule (lin), using the zero line) is $x_S\hat Q$.  (ii) Otherwise
the indicator point of $S$ is good and is $1$ on $S$, so $p\mid c$; hence $x_S\hat Q\equiv0$
pointwise, and (congr) applied to the zero line gives it.

\emph{Step} $V=v::V'$.  Here $v\notin V'$ and $v\notin S$; $Q|_{v=0}$ and $\partial_vQ$ have their
variables in $V'$.
\begin{itemize}
\item $x_S\cdot\widehat{Q|_{v=0}}$: the induction hypothesis with the same $S$ and $d$.  Its
vanishing hypothesis holds because $\widehat{Q|_{v=0}}(a)=\hat Q(a[v:=0])$ and setting $x_v$
to $0$ keeps a point good and keeps it $1$ on $S$ ($v\notin S$).
\item $x_{v::S}\cdot\widehat{\partial_vQ}$: if $d=0$, all terms of $Q$ are constants,
$\widehat{\partial_vQ}=0$, and the line is the zero line.  If $d=d'+1$, use the induction
hypothesis with prefix $v::S$ and bound $d'$: $\deg\partial_vQ\le d'$ and
$|v::S|+d'=|S|+d\le D$.  For its vanishing hypothesis let $a$ be good with $a_v=1$ and
$a=1$ on $S$; then $\hat Q(a)=\hat Q(a[v:=0])+\widehat{\partial_vQ}(a)$ and both values of
$\hat Q$ are $\equiv0$, so $\widehat{\partial_vQ}(a)\equiv0$.
\end{itemize}
Rule (lin) with coefficients $1,1$ gives $x_S\widehat{Q|_{v=0}}+x_vx_S\widehat{\partial_vQ}=x_S\hat Q$.

\emph{Degree bookkeeping.}  The invariant $|S|+d\le D$ is preserved: passing to
$\partial_vQ$ moves one unit from $d$ to $|S|$, and passing to $Q|_{v=0}$ changes neither.
The prefix $x_S$ is never multiplied onto a derived line: the recursion only combines lines
by (lin) and (congr), which have no degree guard, and every product with a variable is
formed inside the auxiliary fact, building $x_{S'}x_ux_w$ upward from the axiom $x_ux_w$
with all intermediate lines of cube degree $<|S|\le D$.  The zero line is used in case (ii),
in the case $d=0$, and as the second argument of (lin) in case (i).

\end{proof}

\begin{proposition}\label{prop:A-upper}
Let $t\ge2$ and $n+k>0$.  If $\binom{n+k}t-\binom nt$ is a unit mod $p$, then
$\mathrm{ontoFPHP}(n,k)$ has a degree-$t$ refutation over
$\mathbb Z/p$.\lean{[0459] OntoLucas.ontoFPHP\_upper}
\end{proposition}

\begin{proof}
The zero line is $0\cdot P_0$.  At a good point (collision and functional axioms
hold) every $r_i$ and every $c_j$ is $0$ or $1$ and $\sum_ir_i=\sum_jc_j=N$, so
$e_t(r)=\binom Nt=e_t(c)$.\lean{[0459] OntoLucas.good\_eq}  By
Lemma~\ref{lem:quot} (here $t\ge2$ makes the quadratic axioms available)
$e_t(r)-e_t(c)$ is derivable in degree $t$.  With Lemma~\ref{lem:esym},
$\bigl(e_t(r)-\binom{n+k}t\bigr)-\bigl(e_t(c)-\binom nt\bigr)-\bigl(e_t(r)-e_t(c)\bigr)
=\binom nt-\binom{n+k}t$ is derivable; multiply by $-u$, where $u$ inverts
$\binom{n+k}t-\binom nt$ mod $p$, and apply (congr).
\end{proof}

\begin{lemma}\label{lem:A-binom}
Let $p$ be prime.  Then $p\mid\binom pr$ for $0<r<p$; $\binom{m+p}r\equiv\binom mr$
for $r<p$; $\binom{m+p}p\equiv\binom mp+1$; hence
$\binom{n+pc}p-\binom np\equiv c\pmod p$, which is a unit when
$p\nmid c$.\lean{[0607] Lucas.dvd\_bin, shift, shiftp, telescope, unit\_c}
\end{lemma}

\begin{proof}
$p\mid\binom pr$: from the absorption identity $r\binom pr=p\binom{p-1}{r-1}$ and
Euclid's lemma, since $p\nmid r$.  The two shift congruences follow by induction on
$m$ from Pascal's rule, the base case $m=0$ being the first statement (and
$\binom pp=1$).  Telescoping $c$ times gives the congruence; B\'ezout gives the
inverse of $c$.
\end{proof}

\begin{remark}
This congruence is what Lucas's theorem gives here, but the Lean proof does not
invoke Lucas's theorem (the file is called \texttt{Lucas.lean}); it proves the
shift congruences directly as above.
\end{remark}

\subsubsection*{A.2 The lower bound: peeling}

Theorem~\ref{thm:A-lower} follows from Theorem~\ref{thm:ext-criterion} with
$D=p-1$ and $n_0=2p-3$ once $\mathrm{EXT}_p(2p-3,k,d)$ holds for all $d\le p-2$.
That is the case $t=p-1$ of:

\begin{proposition}\label{prop:ext-d}
Let $p\mid k$, $2d+1\le n$, and suppose $s$ is invertible mod $p$ for
$1\le s\le d+1$.  Then $\mathrm{EXT}_p(n,k,d)$.\lean{[0607] Peel.ext\_d}
\end{proposition}

For $d\le p-2$ we have $2d+1\le 2p-3$ and $d+1\le p-1$, so the hypotheses hold for
prime $p$; this proves Theorem~\ref{thm:A-lower}.\lean{[0607] Peel.no\_refutation\_units}
The rest of A.2 proves Proposition~\ref{prop:ext-d}.

\paragraph{Rectangle products.}
Let $j\ge0$, let $P=((a_s,a'_s))_{s<j}$ be pigeon pairs and $Q=((b_s,b'_s))_{s<j}$
hole pairs, with the $2j$ pigeons distinct and the $2j$ holes distinct.  Put
$\rho_s(x_{ib})=([i=a_s]-[i=a'_s])([b=b_s]-[b=b'_s])$, the signed indicator of the
rectangle $\{a_s,a'_s\}\times\{b_s,b'_s\}$.  For a matching $g$ with $|g|=j$,
$\mathrm{RP}_j(P,Q)(g)=\sum_\sigma\prod_{s<j}\rho_s(\sigma(s))$, the sum over the
bijections $\sigma$ from slots to the edges of $g$ (a permanent; $\mathrm{RP}_0=1$).\lean{[0607] SJ.RP, SJ.jp, LemK.VP}

\begin{lemma}[$S_j$]\label{lem:SJ}
Let $2(j+1)\le n$ and let $(j+1)!$ be invertible mod $p$.  If
$\mathrm{Bal}_p(f,0,j)$ (``$f$ has zero links at level $j$''), then $f$ lies in
$\mathrm{Span}_{j+1}$ of the rectangle products
$\mathrm{RP}_{j+1}(P,Q)$.\lean{[0607] SJ.SJ}
\end{lemma}

\begin{proof}[Proof outline (from the Lean structure)]
For an injective pigeon tuple $u$ and hole tuple $w$ of length $j+1$, put
$F(u,w)=f(\{x_{u_sw_s}\})$.  Zero links translate into vanishing slot sums:
the pigeon-slot sums of $F$ are column sums of $f$ and the hole-slot sums are row
sums.  \emph{Lemma~K} (characteristic-free): a function on injective $j$-tuples
from a ground set with at least $2j$ points whose slot sums all vanish is, on
valid tuples, a finite sum $\sum_P \mathrm{DP}_P(u)\,c_P(w)$ with
$\mathrm{DP}_P(u)=\prod_s([u_s=P_s^1]-[u_s=P_s^2])$ and each $c_P$ an integer
combination of the functions $F(u',\cdot)$.\lean{[0607] LemK.lemK (file LemKMain.lean)}
Applying it over pigeons and then over holes gives
$F\equiv\sum d\,\mathrm{DP}_P(u)\mathrm{DP}_Q(w)$.\lean{[0607] SJ.SJ\_tuple}
Summing over the $(j+1)!$ orderings of a matching $g$ turns
$\mathrm{DP}_P(u)\mathrm{DP}_Q(w)$ into $\mathrm{RP}_{j+1}(P,Q)(g)$, so
$(j+1)!\,f\equiv\sum d\,\mathrm{RP}_{j+1}(P,Q)$ on matchings of size
$j+1$;\lean{[0607] SJ.SJ\_avg} multiply by the inverse of $(j+1)!$.

It remains to prove Lemma~K, which we restate in the form proved in Lean.  Let $X$ be a
finite set (Lean: a Boolean predicate read below a bound $M$) and $W$ any set of parameters.
A \emph{valid} $j$-tuple is $u=(u_0,\dots,u_{j-1})$ of distinct points of $X$.  For
$F\colon\{\text{valid $j$-tuples}\}\times W\to\mathbb Z$, the \emph{slot-$t$ sum} at a valid
$u$ is $\sigma^X_tF(u,w)=\sum_aF(u[t:=a],w)$ over the $a\in X$ different from all $u_s$,
$s\ne t$.  Write $\mathrm{KC}_m(X,j,F)$ if $m\mid\sigma^X_tF(u,w)$ for all valid $u$, $t<j$,
$w$ ($m\in\mathbb Z$ arbitrary, $m=0$ meaning exactly).  A pair list
$P=((P^1_s,P^2_s))_{s<j}$ is \emph{valid} if its $2j$ points are distinct points of $X$, and
$\mathrm{DP}_P(u)=\prod_{s<j}\delta_{P_s}(u_s)$ with $\delta_q(x)=[x=q^1]-[x=q^2]$.  Finally
$\mathrm{LC}_X(F)$ is the set of functions $c(w)=\sum_\ell a_\ell F(u^{(\ell)},w)$, a finite
integer combination with valid $u^{(\ell)}$.\lean{[0607] LemK.Valid, sig, KC, DP, VP, LC}

\smallskip\noindent\emph{Lemma K.}  If $|X|\ge2j$ and $\mathrm{KC}_m(X,j,F)$, there is a finite
list of pairs $(P,c_P)$, $P$ valid and $c_P\in\mathrm{LC}_X(F)$, with
$F(u,w)\equiv\sum_P\mathrm{DP}_P(u)\,c_P(w)\pmod m$ for all valid $u$ and all
$w$.\lean{[0607] LemK.LK, lemK (LemKMain.lean)}

\smallskip
Three facts are used.  (K1) \emph{Difference products have zero slot sums}: if the pairs of
$P$ are pairwise disjoint and the two points of $P_t$ are distinct and lie in $Y$, then
$\sigma^Y_t\mathrm{DP}_P(u)=0$ exactly.  For $\mathrm{DP}_P(u[t:=a])=\delta_{P_t}(a)\,R$ with
$R$ the product of the other factors, so the sum is
$R\,([P^1_t\text{ free}]-[P^2_t\text{ free}])$; if $P^1_t$ (say) is not free, it equals some
$u_s$, $s\ne t$, and $\delta_{P_s}(u_s)=0$ by disjointness, so $R=0$.\lean{[0607] LemK.sig\_DP}
(K2) \emph{Vanishing off an anchor}: if $z\in X$, $\mathrm{KC}_m(X,j+1,H)$, and $H\equiv0$ at
all valid tuples not containing $z$, then $H\equiv0$ at all valid tuples.  If $u_t=z$, the
term $a=z$ of $\sigma^X_tH(u)$ is $H(u)$ itself and every other term is at a tuple without
$z$, so $H(u)\equiv\sigma^X_tH(u)\equiv0$.\lean{[0607] LemK.zero\_of\_del}
(K3) $\mathrm{LC}$ is closed under sums and integer multiples, and is transitive: a combination
of functions each in $\mathrm{LC}_X(F)$ lies in $\mathrm{LC}_X(F)$.\lean{[0607] LemK.LC\_add, LC\_smul, LC\_tr}

\smallskip\noindent\emph{Proof of Lemma K.}  Induction on $j$.  For $j=0$ the only valid tuple
is empty; take the single pair $(\emptyset,F(\emptyset,\cdot))$, with
$\mathrm{DP}_\emptyset=1$.\lean{[0607] LemK.lemK\_zero}
For the step let $|X|\ge2j+2$ and $\mathrm{KC}_m(X,j+1,F)$.\lean{[0607] LemK.lemK\_step}
\begin{enumerate}
\item \emph{Anchor.}  Fix $z\in X$ and put $X'=X\setminus\{z\}$, so $|X'|\ge2j+1$.
\item \emph{Slot defects.}  For $t\le j$ and a valid $j$-tuple $v$ over $X'$ let
$G_t(v,w)=-F(\mathrm{ins}_t(z,v),w)$, where $\mathrm{ins}_t(z,v)$ inserts $z$ at slot $t$
(a valid $(j+1)$-tuple over $X$).  A point can fill slot $t'$ of $v$ over $X'$ if and only if
it can fill the corresponding slot ($t'$ if $t'<t$, else $t'+1$) of $\mathrm{ins}_t(z,v)$
over $X$, so the slot sums of $G_t$ are minus slot sums of $F$ and
$\mathrm{KC}_m(X',j,G_t)$.\lean{[0607] LemK.free\_ins, ins\_upd, valid\_ins}  The induction
hypothesis ($|X'|\ge2j$) gives a list $B_t$ of pairs $(Q,c)$, $Q$ valid over $X'$, with
$G_t(v,w)\equiv\sum_{B_t}\mathrm{DP}_Q(v)c(w)$; by (K3) each $c$ lies in
$\mathrm{LC}_X(F)$ (at length $j+1$).
\item \emph{Correction.}  For each such $Q$ choose $w_Q\in X'$ outside the $2j$ points of
$Q$, possible as $|X'|>2j$.\lean{[0607] LemK.fresh}  Let $\varepsilon(u,w)$ be the sum, over
$t\le j$ and $(Q,c)\in B_t$, of $\mathrm{DP}_{P}(u)c(w)$ with $P=\mathrm{ins}_t((w_Q,z),Q)$
(the pair $(w_Q,z)$ inserted at slot $t$; $P$ is valid over $X$), and put
$\psi=F-\varepsilon$.\lean{[0607] LemK.insP, VP\_insP}
\item \emph{$\psi$ has zero slot sums over $X'$.}  Let $u$ be valid over $X'$ and $s\le j$.
Since $z$ is not in $u$, it is free at slot $s$, and $u[s:=z]=\mathrm{ins}_s(z,\mathrm{skip}_su)$
($\mathrm{skip}_s$ deletes slot $s$); hence
$\sigma^{X'}_sF(u)=\sigma^X_sF(u)+G_s(\mathrm{skip}_su)$.  In $\sigma^{X'}_s\varepsilon(u)$ the
terms with $t\ne s$ vanish by (K1), as the slot-$s$ pair of $P$ is a pair of $Q$, inside $X'$.
For $t=s$, $\mathrm{DP}_P(u[s:=a])=\delta_{(w_Q,z)}(a)\,\mathrm{DP}_Q(\mathrm{skip}_su)$; summing
over free $a\in X'$ ($z\notin X'$) gives $\mathrm{DP}_Q(\mathrm{skip}_su)$ if $w_Q$ is free, and
otherwise $0=\mathrm{DP}_Q(\mathrm{skip}_su)$, because then $w_Q$ is an entry of
$\mathrm{skip}_su$ lying off the pairs of $Q$.\lean{[0607] LemK.DP\_insP, DP\_off}  So
\[
\sigma^{X'}_s\psi(u)=\sigma^X_sF(u)+\Bigl(G_s(\mathrm{skip}_su)-\sum_{B_s}\mathrm{DP}_Q(\mathrm{skip}_su)\,c\Bigr)\equiv0+0 .
\]
\item \emph{Slices.}  For $a\in X'$ let $\psi_a(v)=\psi((a,v))$ on valid $j$-tuples $v$ over
$X''=X'\setminus\{a\}$, where $(a,v)$ prepends $a$.  Slot $s$ of $v$ is slot $s+1$ of
$(a,v)$, so $\mathrm{KC}_m(X'',j,\psi_a)$; and $|X''|=|X|-2\ge2j$.  The induction hypothesis
gives $B_a$, with coefficients in $\mathrm{LC}_X(F)$ by (K3) ($\psi_a(v)$ is $F((a,v))$ minus a
combination of the coefficients of $\varepsilon$).  Let $B_\psi$ consist of the pairs
$(\mathrm{ins}_0((a,z),Q),c)$ for $a\in X'$, $(Q,c)\in B_a$.  For $u$ valid over $X'$,
$u_0\ne z$ gives $\mathrm{DP}_{\mathrm{ins}_0((a,z),Q)}(u)=[u_0=a]\,\mathrm{DP}_Q(\mathrm{skip}_0u)$, so
$\sum_{B_\psi}\mathrm{DP}\,c=\sum_{B_{u_0}}\mathrm{DP}_Q(\mathrm{skip}_0u)\,c\equiv\psi_{u_0}(\mathrm{skip}_0u)=\psi(u)$.\lean{[0607] LemK.valid\_cns, free\_cns, cns\_skip}
\item \emph{Assembly.}  Take the pairs of $\varepsilon$ followed by $B_\psi$, and let
$H=F-\sum\mathrm{DP}_P\,c_P$ over this list.  $\mathrm{KC}_m(X,j+1,H)$ holds since $F$ satisfies
it and each $\mathrm{DP}_P$ ($P$ valid over $X$) has zero slot sums over $X$ by (K1).  At
tuples avoiding $z$, $H=\psi-\sum_{B_\psi}\mathrm{DP}\,c\equiv0$ by step~5.  By (K2),
$H\equiv0$ at every valid tuple.
\end{enumerate}
The size hypothesis is used four times: the anchor exists ($|X|>0$), the induction
hypothesis applies on $X'$ and on $X''$ ($\ge2j$ points), and the fresh points $w_Q$ exist
($|X'|>2j$).  No division occurs, which is why Lemma~K holds for every modulus $m$.

\end{proof}

\paragraph{Windows.}  The next lemma realises every rectangle product, up to a
factorial, as a window.

\begin{lemma}[windows for rectangle data]\label{lem:window}
Let $p\mid k$, $j\le d$, $2d+1\le n$, and $P,Q$ as above with $j$ pairs each.  There
is $W\in\mathcal W_d$ with $W(g)=(d-j+1)!\,\mathrm{RP}_j(P,Q)(g)$ for every matching
$g$ with $|g|=j$.\lean{[0607] WinMain.window}
\end{lemma}

\begin{proof}
Split pigeons and holes into blocks: the $j$ rectangle blocks
$\{a_s,a'_s\}\times\{b_s,b'_s\}$; $T=n-j-d-1$ blocks $\{x_t\}\times\{h_t\}$ on fresh
pigeons and holes; and a rest block with $m=n-2j-T=d-j+1$ holes and
$n+k-2j-T\equiv m\pmod p$ pigeons (here $p\mid k$ is used).\lean{[0607] WinMain.good\_of, Blocks.fresh\_fn}
$W$ is the block product: on a matching with an edge joining two different blocks
it is $0$; otherwise it is the product of local factors evaluated on the parts of
the matching inside the blocks.\lean{[0607] WinJ.Wn, BProd.PB}  The local factors
are: on a rectangle block, $A_s$ with $A_s(\emptyset)=0$, $A_s=\rho_s$ on single
edges, and $A_s=x_{a_sb_s}x_{a'_sb'_s}-x_{a_sb'_s}x_{a'_sb_s}$ (the determinant)
on its two $2$-matchings;\lean{[1402] Win.Aw} the constant $1$ on each $1\times1$
block; and on the rest block the symmetric design of degree $D=d-j+1$, whose value
at a local $s$-matching is $\mathrm{cv}(s)=(m-s)(m-s-1)\cdots(m-D+1)$, so
$\mathrm{cv}(s)=(m-s)!$ as $m=D$.\lean{[1402] Win.cvf, SymDesign.cv}

\emph{Balance} for $e<D+j=d+1$.\lean{[0607] WinJ.win\_bal}
A free line at a matching lies in one block, and its sum factors as the local line
sum times the other factors.\lean{[0607] BProd.line\_PB, bal\_winB}  Each local
factor satisfies the local window identity wherever a free line exists:
for the symmetric design at a local $s$-matching with $s<D$ the free row sum is
$(m-s)\,\mathrm{cv}(s+1)=\mathrm{cv}(s)$ exactly and the free column sum is
$(\#\text{pigeons}-s)\,\mathrm{cv}(s+1)\equiv\mathrm{cv}(s)$ mod
$p$;\lean{[1402] SymDesign.cv\_step; [0607] LocWin.sym\_rowD, sym\_colD} the constant
block is trivial;\lean{[0607] LocWin.unit\_row, unit\_col} and for $A_s$ the
identity can be checked by hand: at $\emptyset$ the row sum at $a_s$ is
$\rho_s(x_{a_sb_s})+\rho_s(x_{a_sb'_s})=1-1=0=A_s(\emptyset)$ (other rows and the
columns alike), and at the edge $x_{a_sb_s}$ the only free cell is
$x_{a'_sb'_s}$, with $A_s(\{x_{a_sb_s},x_{a'_sb'_s}\})=1=A_s(x_{a_sb_s})$ (the other
three edges alike, with signs).\lean{[0607] LocWin.rect\_row, rect\_col via [1402] Win.blk0\_row, blk0\_col}
A rest line at a matching in which every rectangle block is occupied sees at most
$e-j<D$ rest edges;\lean{[0607] Blocks.cnt\_filter} otherwise some rectangle block
is empty and its factor $A_s(\emptyset)=0$ makes both sides vanish.

\emph{Value.} At a matching $g$ with $|g|=j$ and $W(g)\ne0$ every rectangle block
carries exactly one edge, so the rest block is empty and contributes
$\mathrm{cv}(0)=D!$; the permanent of functions with pairwise disjoint supports is
the product of their sums,\lean{[0607] Blocks.jp\_prod; WinJ.win\_val} giving
$W(g)=(d-j+1)!\,\mathrm{RP}_j(P,Q)(g)$.
\end{proof}

\begin{corollary}\label{cor:rp-span}
If moreover $(d-j+1)!$ is invertible mod $p$, every $\mathrm{RP}_j(P,Q)$ lies in
$\mathrm{Span}_j(\mathcal W_d)$.\lean{[0607] Peel.rp\_span}
\end{corollary}

\begin{proof}
$\mathrm{RP}_j\equiv u\,W$ on $j$-matchings, $u$ the inverse of $(d-j+1)!$.
\end{proof}

\paragraph{Peeling.}
Call $f$ \emph{admissible at level $0$} always, and \emph{at level $\ell+1$} if
$\mathrm{Bal}_p(f,c,\ell)$ for some $c$ admissible at level $\ell$.  By the flag
lemma (Lemma~\ref{lem:flag}) iterated, $\mathrm{Bal}_p(f,c,\ell)$ already implies
that $f$ is admissible at level $\ell+1$.\lean{[0607] Peel.Adm, adm\_of\_bal}
$\mathcal W_d$ is closed under sums and multiples, so every member of
$\mathrm{Span}_\ell(\mathcal W_d)$ agrees mod $p$ on $\ell$-matchings with a single
$T\in\mathcal W_d$.\lean{[0607] Peel.span\_rep}

\begin{proposition}[peeling]\label{prop:peel}
Under the hypotheses of Proposition~\ref{prop:ext-d}, for every $\ell\le d$ every
$f$ admissible at level $\ell$ lies in
$\mathrm{Span}_\ell(\mathcal W_d)$.\lean{[0607] Peel.claim}
\end{proposition}

\begin{proof}
Induction on $\ell$.  $\ell=0$: Lemma~\ref{lem:window} with $j=0$ gives a window $W$
with $W(\emptyset)=(d+1)!$ (the \emph{type-$0$ window}: $n-d-1$ unit blocks and a
rest block with $d+1$ holes); $f\equiv f(\emptyset)\,u\,W$ at $\emptyset$, $u$ the
inverse of $(d+1)!$.
$\ell\to\ell+1$: let $\mathrm{Bal}_p(f,c,\ell)$ with $c$ admissible.  By induction
$c\equiv T$ on $\ell$-matchings for some $T\in\mathcal W_d$.  Since
$\mathrm{Bal}_p(T,T,\ell)$, the residual $f-T$ satisfies
$\mathrm{Bal}_p(f-T,\,c-T,\,\ell)$, i.e.\ it has zero links at level $\ell$.  By
Lemma~\ref{lem:SJ} (with $(\ell+1)!$ invertible and $2(\ell+1)\le n$) it lies in the
span of the $\mathrm{RP}_{\ell+1}$, which by Corollary~\ref{cor:rp-span} lies in
$\mathrm{Span}_{\ell+1}(\mathcal W_d)$.  Add $T$ back.
\end{proof}

\begin{proof}[Proof of Proposition~\ref{prop:ext-d}]
Each $T\in\mathcal W_d$ has the extension $T$ itself ($\mathrm{Bal}_p(T,T,d)$).
An $f$ with $\mathrm{Bal}_p(f,c,d-1)$ is admissible at level $d$, so lies in
$\mathrm{Span}_d(\mathcal W_d)$ by Proposition~\ref{prop:peel};
Lemma~\ref{lem:ext-span} concludes.  All factorials used are at most $(d+1)!$.
\end{proof}

\paragraph{Mechanism of A.}
The upper bound is Riis's symmetric-function refutation at $t=p$, made a unit by the
Lucas-type congruence $\binom{n+pc}p-\binom np\equiv c$.  The lower bound is
\emph{peeling}: a balanced function at level $d$ is written, level by level, as a
window plus a zero-link residual; zero-link residuals are spans of rectangle
products ($S_j$, via Lemma~K), and rectangle products are windows up to the factorial
$(d-j+1)!$.  The only arithmetic input is that $1!,\dots,(d+1)!$ are units, which for
$d=p-2$ is exactly why the bound stops at $p-1$: at $d=p-1$ the type-$0$ window has
value $p!\equiv0$.  The board must hold the blocks, which gives $n\ge 2d+1=2p-3$.

%======================================================================
\subsection{B. Degree at least $p+1$ when $p^2$ divides the surplus}
\label{sec:lean-B}

\begin{theorem}[B]\label{thm:B}
Let $p$ be prime and $c,n\in\mathbb N$ with $n\ge2p-1$.  Then
$\mathrm{ontoFPHP}(n,cp^2)$ has no degree-$p$ semantic PC refutation over
$\mathbb Z/p$, and no degree-$p$ standard PC refutation
over~$\mathbb F_p$.\lean{[1249] uniform\_no\_degree\_p\_unconditional, uniform\_no\_standard\_degree\_p\_unconditional}
\end{theorem}

The Lean statement allows $c=0$ (then the formula is satisfiable and the claim is
trivial) and states the bound $D=p$; by monotonicity
(Lemma~\ref{lem:mono}) there is no
refutation of degree at most $p$ (for $c\ge1$),\lean{[maint] DegreeMono.uniform\_no\_degree\_le\_p, uniform\_no\_standard\_degree\_le\_p} so the PC degree is at least $p+1$.  The surplus
is written $(c\cdot p)\cdot p$.  The proof produces, as an intermediate result:

\begin{theorem}[unit design]\label{thm:B-design}
For prime $p$ and every $c$ there is $T$ with
$\mathrm{Bal}^{(2p-1,\,cp^2)}_p(T,T,e)$ for all $e\le p-1$ and
$T(\emptyset)\equiv1\pmod p$.\lean{[1249] uniform\_unit\_design\_unconditional}
\end{theorem}

The proof has four steps: a criterion turning such a $T$ into the lower bound
(B.1); the construction of $T$ from a \emph{rectangular lift} (B.2); the rectangular
lift from a \emph{square lift} on the $(2p-1)\times(2p-1)$ board (B.3); and an
explicit square lift for every prime (B.4).  In Lean, B.2--B.3 are proved first
under the hypothesis that a square lift exists
(\texttt{SquareLiftAssumption}),\lean{[1249] Main.lean: uniform\_unit\_design, uniform\_no\_degree\_p}
and B.4 discharges it.\lean{[1249] squareLiftAssumption (Unconditional.lean)}

\subsubsection*{B.1 The design-to-lower-bound criterion}

\begin{theorem}\label{thm:B-criterion}
Let $p$ be prime, $p\mid k$, $n_0\ge2p-1$, and let $T_0$ satisfy
$\mathrm{Bal}^{(n_0,k)}_p(T_0,T_0,e)$ for all $e\le p-1$ and $T_0(\emptyset)\equiv1$.
Then for every $n\ge n_0$, $\mathrm{ontoFPHP}(n,k)$ has no degree-$p$ refutation
over $\mathbb Z/p$.\lean{[0607] VUnit.no\_refutation\_of\_unit}
\end{theorem}

\begin{proof}
Apply Theorem~\ref{thm:ext-criterion} with $D=p$; we need
$\mathrm{EXT}_p(n_0,k,d)$ for $d\le p-1$.  For $d\le p-2$ this is
Proposition~\ref{prop:ext-d} ($2d+1\le2p-3\le n_0$, $d+1\le p-1$).  For $d=p-1$ run
the peeling of Proposition~\ref{prop:peel} with $T_0$ in place of the type-$0$
window: at level $0$, $f\equiv f(\emptyset)T_0$ at $\emptyset$ because
$T_0(\emptyset)\equiv1$; at levels $\ell+1\ge1$ the argument is unchanged and uses
only the units $(\ell+1)!$ with $\ell+1\le p-1$ (Lemma~\ref{lem:SJ}) and
$(d-j+1)!$ with $j\ge1$, i.e.\ at most $(p-1)!$
(Corollary~\ref{cor:rp-span}).\lean{[0607] VUnit.claimU, rp\_span', ext\_unit}
The board condition is $2d+1=2p-1\le n_0$.  The type-$0$ window would have needed
$(d+1)!=p!$, which is $0$ mod $p$; $T_0$ is exactly what replaces it.
\end{proof}

\subsubsection*{B.2 From a rectangular lift to a unit design}

\begin{definition}\label{def:rect-lift}
A \emph{rectangular lift} for $(p,k)$ is a function $H$ on matchings of the board
with $2p-1$ holes and $2p-1+k$ pigeons such that $H(g)\equiv1\pmod p$ for every
matching with $|g|=p$, and at every matching with $|g|=p-1$ every free row sum and
every free column sum of $H$ is exactly $0$ (as an
integer).\lean{[1249] RectangularLift (Boundary.lean)}
\end{definition}

\begin{lemma}[top layer]\label{lem:top}
Let $p$ be prime and $H$ a rectangular lift for $(p,cp^2)$.  Put $U(g)=(H(g)-1)/p$ for
$|g|=p$.  Then $\mathrm{Bal}^{(2p-1,cp^2)}_p(U,-1,p-1)$.\lean{[1249] RectangularLift.top, top\_bal}
\end{lemma}

\begin{proof}
Let $|g|=p-1$.  At a free pigeon there are $(2p-1)-(p-1)=p$ free holes, so
$0=\mathrm{row}_H(g,i)=p+p\,\mathrm{row}_U(g,i)$ and $\mathrm{row}_U(g,i)=-1$.  At a
free hole there are $(2p-1+k)-(p-1)=p+k$ free pigeons, so
$0=(p+k)+p\,\mathrm{col}_U(g,j)$, i.e.\ $\mathrm{col}_U(g,j)=-1-k/p=-1-cp\equiv-1$.
This is the step that uses $p^2\mid k$.\lean{[1249] RectangularLift.line\_eq; [1402] SymDesign.rowS\_const, colS\_const}
\end{proof}

\begin{lemma}[symmetric boundary]\label{lem:boundary}
Let $n>0$, $p\mid k$, and $\mathrm{Bal}^{(n,k)}_p(U,-1,d)$.  Define
$T(g)=-\mathrm{cv}_{n,d}(|g|)=-(n-|g|)(n-|g|-1)\cdots(n-d+1)$ for $|g|\le d$ and
$T(g)=U(g)$ for $|g|=d+1$.  Then $T(\emptyset)=-n(n-1)\cdots(n-d+1)$ and
$\mathrm{Bal}_p(T,T,e)$ for every $e\le d$.\lean{[0634] FactorBoundary.design\_of\_top, lower\_bal, lower\_top; TopSplice.splice\_bal}
\end{lemma}

\begin{proof}
For $e<d$ the free row sum is $-(n-e)\mathrm{cv}(e+1)=-\mathrm{cv}(e)=T(g)$ and the
free column sum is $-(n+k-e)\mathrm{cv}(e+1)\equiv-\mathrm{cv}(e)=T(g)$, using $p\mid k$.  For
$e=d$ the line sums are sums of $U$, which are $\equiv-1=T(g)$ since
$\mathrm{cv}(d)=1$.
\end{proof}

\begin{proposition}\label{prop:B-design}
If $p$ is prime and a rectangular lift for $(p,cp^2)$ exists, then the conclusion
of Theorem~\ref{thm:B-design} holds.\lean{[1249] unit\_design\_of\_rectangular\_lift}
\end{proposition}

\begin{proof}
Lemma~\ref{lem:top} and Lemma~\ref{lem:boundary} with $n=2p-1$, $d=p-1$ give $T$ with
$T(\emptyset)=-(2p-1)(2p-2)\cdots(p+1)$.  Each factor $p+r$ ($1\le r\le p-1$) is
$\equiv r$, a unit, so $T(\emptyset)$ is a unit;\lean{[0634] FactorUnit.ff\_unit, empty\_unit}
scale $T$ by its inverse.
\end{proof}

Theorem~B follows from Proposition~\ref{prop:B-design} and
Theorem~\ref{thm:B-criterion} with $n_0=2p-1$ ($p\mid cp^2$), and the standard form
from Theorem~\ref{thm:refutes-iff}.\lean{[1249] no\_refutation\_of\_rectangular\_lift; uniform\_no\_standard\_degree\_p}

\subsubsection*{B.3 From the square to the rectangle}

\begin{definition}\label{def:square-lift}
A \emph{square lift} (integral lift) for $p$ is a function $H$ on matchings of the
square board, $2p-1$ pigeons and $2p-1$ holes ($k=0$), with $H(g)\equiv1$ for
$|g|=p$ and all free row and column sums exactly $0$ at matchings with
$|g|=p-1$.\lean{[1249] IntegralLift (Lift.lean)}
\end{definition}

\begin{proposition}\label{prop:transfer}
For prime $p$, a square lift for $p$ and any $c$ yield a rectangular lift for
$(p,cp^2)$.\lean{[1249] rectangleFromSquare (Transfer.lean)}
\end{proposition}

\begin{proof}
Let $k=cp^2$.  For each set $A$ of $2p-1$ pigeons among $2p-1+k$, let $H_A$ be $H$
transported to the pigeons $A$ by a relabelling and extended by $0$ to every
matching that uses a pigeon outside $A$.\lean{[1249] Selection.embed, Selection.perm, Pad.pad}
Put $\widetilde H=\sum_AH_A$.\lean{[1249] transferred}

\emph{Exact harmonicity.}  Padding by zeros keeps all free line sums at
$(p-1)$-matchings exactly $0$: a row of an outside pigeon sums zeros, a line of an
inside pigeon or of a hole only gains zero terms, and at a matching that already
uses an outside pigeon every term is $0$.  Relabelling preserves $\mathrm{Bal}$, and
a sum of exactly harmonic functions is exactly harmonic.\lean{[1249] Pad.pad\_bal, embedded\_bal, embed\_harmonic, bal\_sum; [1402] Relabel.bal\_relab}

\emph{Unit.}  For $|g|=p$, $H_A(g)\equiv[\text{the pigeons of $g$ lie in $A$}]$.
The $p$ pigeons of $g$ are distinct, so the number of such $A$ is
$\binom{2p-1+k-p}{(2p-1)-p}=\binom{k+p-1}{p-1}$.\lean{[1249] embed\_unit, selection\_count, SubsetCount.coverCount\_eq, transferred\_congr\_count}
Finally $\binom{cp^2+p-1}{p-1}\equiv\binom{p-1}{p-1}=1$, by iterating
$\binom{a+p}r\equiv\binom ar$ ($r<p$, Lemma~\ref{lem:A-binom}) $cp$
times.\lean{[1249] Multiplicity.subset\_multiplicity\_congr, bin\_shift\_multiple}
(This step only needs $p\mid k$.)
\end{proof}

\subsubsection*{B.4 An explicit square lift}

Split the rows and the columns of the square board into the blocks
$B_0=\{0,\dots,p-2\}$, $B_1=\{p-1,\dots,2p-3\}$, $B_2=\{2p-2\}$ of sizes
$a=(p-1,p-1,1)$.\lean{[1249] Block.squarePartition, capacity, blockOf}  For a
matching $g$ let $M(g)\in\mathbb N^{3\times3}$ count its edges from row block $I$ to
column block $J$, with row sums $\mathrm{rc}_I$ and column sums $\mathrm{cc}_J$.
$M$ is \emph{feasible} if $\mathrm{rc}_I\le a_I$ and $\mathrm{cc}_J\le a_J$; its
\emph{total} is the number of edges.  $E_{IJ}$ is the unit matrix at
$(I,J)$.\lean{[1249] Block.counts, Feasible, inc}

\begin{lemma}[reduction to count matrices]\label{lem:certificate}
Let $V\colon\mathbb N^{3\times3}\to\mathbb Z$ satisfy, for feasible $M$:
(unit) $V(M)\equiv1$ if the total is $p$; and, if the total is $p-1$,
\begin{gather*}
\textstyle\sum_J(a_J-\mathrm{cc}_J(M))\,V(M+E_{IJ})=0\ \text{for every }I,\\
\textstyle\sum_I(a_I-\mathrm{rc}_I(M))\,V(M+E_{IJ})=0\ \text{for every }J .
\end{gather*}
Then $H(g)=V(M(g))$ is a square lift.\lean{[1249] Block.Certificate, Certificate.toIntegralLift, row\_formula, col\_formula, feasible\_counts}
\end{lemma}

\begin{proof}
At a matching $g$, a free pigeon in row block $I$ has $a_J-\mathrm{cc}_J(M(g))$ free
holes in column block $J$, and each extension has count matrix $M(g)+E_{IJ}$.
\end{proof}

Define, for feasible $M$ of total $p$,
\begin{multline*}
V_0(M)=-(-1)^{M_{01}+M_{10}}\,(p-2)!\;N(M)\\
\cdot(p-1-\mathrm{rc}_0)!\,(p-1-\mathrm{rc}_1)!\,(p-1-\mathrm{cc}_0)!\,(p-1-\mathrm{cc}_1)!
\end{multline*}
and $V_0(M)=0$ otherwise, where the numerator $N(M)$ is, for feasible $M$ (cases
read top to bottom; feasibility allows at most one edge in block row $2$ and one in
block column $2$),
\begin{itemize}
\item $N(M)=p-1$ if $M_{22}=1$;
\item $N(M)=M_{00}M_{11}-M_{01}M_{10}$ if $\mathrm{rc}_2=\mathrm{cc}_2=0$;
\item if $\mathrm{rc}_2=0$ and $\mathrm{cc}_2=1$: $N(M)=-(M_{10}+M_{11})$ when
$M_{02}=1$, and $-(M_{00}+M_{01})$ when $M_{12}=1$;
\item if $\mathrm{rc}_2=1$ and $\mathrm{cc}_2=0$: $N(M)=-(M_{01}+M_{11})$ when
$M_{20}=1$, and $-(M_{00}+M_{10})$ when $M_{21}=1$;
\item if $\mathrm{rc}_2=\mathrm{cc}_2=1$ and $M_{22}=0$: $N(M)=-1$ when
$M_{02}=M_{20}$, and $1$ otherwise.\lean{[1249] Block.numerator, explicitValue (ExplicitFormula.lean)}
\end{itemize}
(A comment in the source says the formula was obtained from a generating-function
coefficient; that derivation is not part of the Lean proof and is not used.)

\begin{lemma}\label{lem:weight}
Let $w(M)=\prod_I(a_I-\mathrm{rc}_I)!\prod_J(a_J-\mathrm{cc}_J)!$ and
$\varphi(M)=(-1)^{M_{01}+M_{10}}N(M)$ when $\mathrm{rc}_2,\mathrm{cc}_2\le1$, else $0$.
For feasible $M$ of total $p$, $V_0(M)=-(p-2)!\,w(M)\,\varphi(M)$; and
$(a_I-\mathrm{rc}_I)(a_J-\mathrm{cc}_J)\,w(M+E_{IJ})=w(M)$ when
$\mathrm{rc}_I<a_I$, $\mathrm{cc}_J<a_J$.\lean{[1249] Block.weight, explicit\_eq, weight\_inc; ExplicitAlgebra.phi, jCore}
\end{lemma}

\begin{lemma}\label{lem:phi-row}
Let $M$ be feasible of total $p-1$ and $\mathrm{rc}_I<a_I$.  Then
$\sum_J\varphi(M+E_{IJ})=0$.  Also $\varphi(M)=0$ whenever $M$ has total $p$ and some
$\mathrm{cc}_J>a_J$; and $\varphi(M^{\mathsf T})=\varphi(M)$.\lean{[1249] Block.phi\_row, phi\_bad\_col, phi\_transpose}
\end{lemma}

\begin{proof}
Feasibility leaves ten patterns for the entries in block row and column~$2$
(at most one edge in row block $2$ and one in column block $2$).\lean{[1249] Block.boundary\_cases}
The Lean proof checks the identity in each of the $10\times3$ cases (pattern, row
block $I$) by rewriting and linear arithmetic (\texttt{simp}, \texttt{omega},
\texttt{grind}).  One case, for illustration: no edges in block row or column $2$
and $I=0$.  Writing $\sigma=(-1)^{M_{01}+M_{10}}$, the three terms are
$\sigma\bigl((M_{00}+1)M_{11}-M_{01}M_{10}\bigr)$,
$-\sigma\bigl(M_{00}M_{11}-(M_{01}+1)M_{10}\bigr)$ and $-\sigma(M_{10}+M_{11})$,
which sum to $0$.
\end{proof}

\begin{proposition}[exact harmonicity]\label{prop:B-harm}
For feasible $M$ of total $p-1$:
$\sum_J(a_J-\mathrm{cc}_J)V_0(M+E_{IJ})=0$ for every $I$ and
$\sum_I(a_I-\mathrm{rc}_I)V_0(M+E_{IJ})=0$ for every $J$.\lean{[1249] Block.explicit\_row, explicit\_col, weighted\_term}
\end{proposition}

\begin{proof}
If $\mathrm{rc}_I=a_I$, every $M+E_{IJ}$ is infeasible and all terms vanish.
Otherwise multiply the row sum by $a_I-\mathrm{rc}_I\ne0$.  By
Lemma~\ref{lem:weight} each term becomes $-(p-2)!\,w(M)\,\varphi(M+E_{IJ})$ (when
$\mathrm{cc}_J=a_J$ the coefficient is $0$ and $\varphi(M+E_{IJ})=0$ by
Lemma~\ref{lem:phi-row}), and the sum vanishes by Lemma~\ref{lem:phi-row}.  Columns
follow from $V_0(M^{\mathsf T})=V_0(M)$.\lean{[1249] Block.explicit\_transpose}
\end{proof}

\begin{lemma}[factorial reflection]\label{lem:reflection}
For every $p\ge1$ and $r\le p-1$:
$(-1)^r\,r!\,(p-1-r)!\equiv(p-1)!\pmod p$.  For $p\ge2$:
$(p-2)!\equiv-(p-1)!$.\lean{[1249] Block.factorial\_reflection, factorial\_tail}
\end{lemma}

\begin{proof}
Induction on $r$: passing from $r$ to $r+1$ multiplies by $-(r+1)$ and divides by
$p-1-r\equiv-(r+1)$.  The difference of consecutive terms is an explicit multiple
of $p$, so no division is needed.  For the second claim,
$(p-1)!=(p-1)(p-2)!\equiv-(p-2)!$.
\end{proof}

\begin{proposition}[common residue]\label{prop:B-cong}
For feasible $M$ of total $p$, $V_0(M)\equiv-\bigl((p-1)!\bigr)^3\pmod p$.\lean{[1249] Block.explicit\_common\_congruence, commonValue, main\_pair}
\end{proposition}

\begin{proof}
The Lean proof treats the ten boundary patterns separately.  In each it rewrites
the four factorial arguments using the total $p$, reduces $V_0(M)$ mod $p$ to
$(p-2)!\,(-1)^{r+t}\,r!\,s!\,t!\,u!$ with $r+s=t+u=p-1$, and applies
Lemma~\ref{lem:reflection} twice:
$(p-2)!\,(p-1)!^2\equiv-(p-1)!^3$.  The first pattern (no edges in block row or
column $2$) runs as follows.  Then $\mathrm{rc}_0+\mathrm{rc}_1=p$ with
$\mathrm{rc}_I\le p-1$, so both are $\ge1$, and similarly for columns; the
factorial arguments are $p-1-\mathrm{rc}_0=\mathrm{rc}_1-1$, etc.  Since
$M_{00}M_{11}-M_{01}M_{10}=pM_{00}-\mathrm{rc}_0\,\mathrm{cc}_0$, and
$(-1)^{\mathrm{rc}_0+\mathrm{cc}_0}=(-1)^{M_{01}+M_{10}}$,
\[
V_0(M)\equiv(p-2)!\,(-1)^{\mathrm{rc}_0}\,\mathrm{rc}_0!\,(\mathrm{rc}_1-1)!\cdot
(-1)^{\mathrm{cc}_0}\,\mathrm{cc}_0!\,(\mathrm{cc}_1-1)!\pmod p,
\]
which has the stated form with $r=\mathrm{rc}_0$, $s=\mathrm{rc}_1-1$,
$t=\mathrm{cc}_0$, $u=\mathrm{cc}_1-1$.  The other nine patterns are verified in
Lean by the same rewriting followed by \texttt{simp}/\texttt{grind}; we have not
written them out.
\end{proof}

\begin{theorem}\label{thm:square-lift}
For every prime $p$ a square lift exists.\lean{[1249] squareLiftAssumption, Block.squareLiftAssumption}
\end{theorem}

\begin{proof}
$(p-1)!$ is a unit (units), so $-\bigl((p-1)!\bigr)^3$ has an inverse
$u$.\lean{[1249] Block.commonValue\_unit}  Then $V=u\,V_0$ satisfies the hypotheses of
Lemma~\ref{lem:certificate} by Propositions~\ref{prop:B-cong} and~\ref{prop:B-harm}.
\end{proof}

\begin{remark}
For $p=2$ and $p=3$ the library also contains explicit square lifts given as data,
whose defining properties are verified by kernel computation (\texttt{decide});
they give the statements for $p=2$, $n\ge3$ and $p=3$, $n\ge5$ independently of
B.4.\lean{[1249] integralLift2, integralLift3, no\_standard\_degree\_two, no\_standard\_degree\_three}
\end{remark}

\paragraph{Mechanism of B.}
An explicit integer function on $p$-matchings of the $(2p-1)\times(2p-1)$ board,
depending only on the $3\times3$ block-count matrix, has all free line sums at
$(p-1)$-matchings exactly $0$ and is $\equiv$ a unit on every $p$-matching (the
\emph{square lift}).  Summing its copies over all choices of $2p-1$ pigeons among
$2p-1+k$ keeps exact harmonicity and multiplies the residue by
$\binom{k+p-1}{p-1}\equiv1$.  Because the line sums are exact, the top layer
$(H-1)/p$ is balanced against $-1$; this is where $p^2\mid k$ enters, since the
column count $p+k$ must leave $k/p\equiv0$.  Spliced under the symmetric
falling-factorial levels it is a \emph{balanced unit design} of degree $p$ on
$2p-1$ holes.  The \emph{design-to-lower-bound criterion} (B.1) is the peeling of A
run at level $p-1$, with this design replacing the type-$0$ window whose value $p!$
vanishes mod $p$.

%======================================================================
\subsection{C. No degree-4 refutation over $\mathbb F_2$ at surplus $8$}
\label{sec:lean-C}

\begin{theorem}[C]\label{thm:C}
For every $n\ge7$, $\mathrm{ontoFPHP}(n,8)$ has no degree-$4$ semantic PC refutation
over $\mathbb Z/2$, and no degree-$4$ standard PC refutation over
$\mathbb F_2$.\lean{[0841] L3.Main.no\_deg4; [1156] StdPCCor.std\_no\_deg4\_k8}
\end{theorem}

By monotonicity (Lemma~\ref{lem:mono}) there is no refutation of degree at most $4$, so the PC degree is at
least $5$.\lean{[maint] DegreeMono.no\_deg\_le4, deg\_ge5}  The proof is Theorem~\ref{thm:ext-criterion} with $p=2$, $k=8$,
$n_0=7$, $D=4$: it suffices to prove $\mathrm{EXT}_2(7,8,d)$ for $d=0,1,2,3$.\lean{[1434] Ext278b.hard\_of\_ext3}
Level $0$ is explicit and levels $1$, $2$ use small kernel-checked certificates; level $3$ is reduced to a single finite
statement, Proposition~\ref{prop:C-z3}, which is established by a kernel-checked
certificate.  Throughout, $n=7$, $k=8$: $15$ pigeons, $7$ holes, $105$ variables.

\paragraph{Relabellings.}  Permutations $\sigma$ of the pigeons and $\tau$ of the
holes act on functions of matchings, and $\mathrm{Bal}_p(y,f,e)$ implies
$\mathrm{Bal}_p(y^{\sigma\tau},f^{\sigma\tau},e)$.\lean{[1402] Relabel.relab, bal\_relab}
So every relabelling of a function with an extension has one.

\begin{lemma}[level 0]\label{lem:C0}
If $p\mid k$ and $n>0$, then $\mathrm{EXT}_p(n,k,0)$.\lean{[1402] Ext0.ext0}
\end{lemma}

\begin{proof}
Given $f(\emptyset)$, put $y(x_{ii})=f(\emptyset)$ for $i<n$,
$y(x_{i0})=f(\emptyset)$ for $i\ge n$, and $y=0$ elsewhere.  Every row sums to
$f(\emptyset)$, columns $j\ne0$ sum to $f(\emptyset)$, and column $0$ sums to
$(k+1)f(\emptyset)\equiv f(\emptyset)$.
\end{proof}

\begin{lemma}[level 1]\label{lem:C1}
$\mathrm{EXT}_2(7,8,1)$.\lean{[1434] Ext278a.ext1}
\end{lemma}

\begin{proof}
Let $H$ be the degree-$3$ design of the library: $H(\emptyset)=1$ and
$\mathrm{Bal}_2(H,H,e)$ for $e\le2$.\lean{[1434] HDes.H, HDes.h\_design}  It is a
two-block product: on pigeons $<12$ $\times$ holes $<4$ it is the library design of
$\mathrm{ontoFPHP}(4,8)$ reindexed, on pigeons $12,13,14$ $\times$ holes $4,5,6$ the
indicator of the diagonal, and $0$ on cross edges.  Every function $c$ on edges
whose row and column sums at $\emptyset$ are all $\equiv\kappa$ lies in
$\mathrm{Span}_1$ of the relabellings of $H$: this is a kernel-checked peeling
certificate (85 elimination steps by sums of relabellings in echelon order, then
21 line propagations), sound by the general checker
theorem.\lean{[1434] Ext278a.span1; L1Data.cert, L1Data.run\_ok; Peel.span\_of\_run}
Each relabelling of $H$ extends itself at level $1$; Lemma~\ref{lem:ext-span}
concludes.
\end{proof}

\begin{lemma}[level 2]\label{lem:C2}
$\mathrm{EXT}_2(7,8,2)$.\lean{[1434] Ext278b.ext2, span2, g\_ext}
\end{lemma}

\begin{proof}
Let $\mathrm{Bal}_2(f,c,1)$.  By the flag lemma at $g=\emptyset$ the link $c$ is
level-$0$ balanced,\lean{[1434] Ext278a.link\_bal; [1402] RowCol.row\_col} so by the
level-$1$ certificate $c\in\mathrm{Span}_1$(relabellings of $H$).  Since each
generator extends itself, there is $u$ in $\mathrm{Span}_2$ of the same generators
with $\mathrm{Bal}_2(u,c,1)$.\lean{[1434] Ext278a.lift1}  Then $f-u$ has zero links at
level $1$, and a kernel-checked level-$2$ certificate (files \texttt{L2C0}--\texttt{L2C8})
puts it in $\mathrm{Span}_2$ of the relabellings of $R_0$, the product of the
rectangles on $\{0,1\}\times\{0,1\}$ and $\{2,3\}\times\{2,3\}$, and of $T$, the
indicator of $2$-matchings inside the $3\times3$ block on pigeons and holes
$\{0,1,2\}$.\lean{[1434] L2Base.R0; [1402] TExt.T; [1434] Ext278b.bases2, cov, p0}
Each generator has a level-$2$ extension: $R_0$ by a three-block product with
block extensions (uses $n\ge7$, $p\mid k$), $T$ by an explicit two-block product,
$H$ by itself.\lean{[1402] RectExt.rect\_ext, TImage.t\_image\_ext, TExt.t\_ext; [1434] Ext278a.gh\_ext}
\end{proof}

\begin{lemma}[lifting to level 3]\label{lem:C-lift}
Let $1<n$ and let $G$, $G_z$ be sets of functions such that
(i) each $\gamma\in G$ has $L$ with $\mathrm{Bal}_p(L,\gamma,2)$ and a level-$3$
extension of $L$;
(ii) every $c$ with $\mathrm{Bal}_p(c,h,1)$ for some $h$ lies in $\mathrm{Span}_2(G)$;
(iii) each $z\in G_z$ has a level-$3$ extension;
(iv) every $z$ with $\mathrm{Bal}_p(z,0,2)$ lies in $\mathrm{Span}_3(G_z)$.
Then $\mathrm{EXT}_p(n,k,3)$.\lean{[1434] Lift.EXT\_lift (with d = 1), Lift.lift\_of\_span}
\end{lemma}

\begin{proof}
Let $\mathrm{Bal}_p(f,c,2)$.  The link $c$ is level-$1$ balanced (flag lemma), so
$c\in\mathrm{Span}_2(G)$ by (ii).  Lifts are linear, so by (i) there is $L$ with
$\mathrm{Bal}_p(L,c,2)$ and an extension $y$ of $L$.  Then $f-L$ has zero links at
level $2$, so by (iv), (iii) and Lemma~\ref{lem:ext-span} it has an extension $y'$.
$y+y'$ extends $f$.
\end{proof}

Let $\mathcal S_{23}$ be the set of $g$ with $\mathrm{Bal}_2(g,g,2)$ and
$\mathrm{Bal}_2(g,g,3)$, and $\mathcal S_3$ those with
$\mathrm{Bal}_2(g,g,3)$.\lean{[1434] Ext278c.SB23, SB3}  For $G=\mathcal S_{23}$ take
$L=\gamma$ in (i); for $G_z=\mathcal S_3$ take the extension $z$ itself in (iii).

\begin{lemma}[condition (ii)]\label{lem:C-span2b}
Every $c$ with $\mathrm{Bal}_2(c,h,1)$ lies in
$\mathrm{Span}_2(\mathcal S_{23})$.\lean{[1434] Ext278c.span2b}
\end{lemma}

\begin{proof}
The link $h$ is level-$0$ balanced.  A second kernel-checked level-$1$ certificate,
over the library's degree-$4$ unit design $D$ (with $D(\emptyset)=1$ and
$\mathrm{Bal}_2(D,D,e)$ for $e\le3$) and a rectangle tower $R_{Bd}$, puts $h$ in
$\mathrm{Span}_1$ of functions self-balanced at levels $1,2,3$, which lifts to $u\in
\mathrm{Span}_2(\mathcal S_{23})$ with $\mathrm{Bal}_2(u,h,1)$.\lean{[1434] Ext278c.span1b, lift1b, bases1b; L1bData; K5Bal; Towers.RBd}
Then $c-u$ has zero links at level $1$, so the level-$2$ certificate of
Lemma~\ref{lem:C2} puts it in the span of relabellings of $R_0$ and $T$.  By kernel
evaluation on all pairs of variables, $R_0$ agrees mod $2$ on $2$-matchings with a
self-balanced tower $R_{RM}$, and $T$ with a relabelling of a tower $T_B$; both are in
$\mathcal S_{23}$.\lean{[1434] Ext278c.zl2, gz\_span; Cong.r0\_cong, Cong.t\_cong; Towers.rrm\_bal}
\end{proof}


\paragraph{Generator lemmas for levels 1--3.}
We prove the properties of the generators used in Lemmas~\ref{lem:C1}, \ref{lem:C2} and
\ref{lem:C-span2b}.  Throughout $n=7$, $k=8$ unless stated otherwise; pigeon $i$ and hole $j$
give the variable $x_{ij}$ of index $7i+j$.

\emph{Block products.}  Block maps $\beta_P$ on pigeons and $\beta_H$ on holes, with values
$<m$, put the edge $x_{ij}$ in block $b$ if $\beta_P(i)=\beta_H(j)=b$, and make it a
\emph{cross edge} otherwise.  For factors $\varphi_b$ on local matchings whose values depend
only on the set of edges, the block product is $\Pi(g)=0$ if $g$ has a cross edge and
$\Pi(g)=\prod_b\varphi_b(g_b)$ otherwise, $g_b$ the edges of $g$ in block $b$.  For a local
matching $\sigma$ of block $b$ and a pigeon $i$ of block $b$ free in $\sigma$, the
\emph{local row sum} $\mathrm{lrow}_{\varphi_b}(\sigma,i)$ is $\sum\varphi_b(\sigma+x_{it})$ over
the holes $t$ of block $b$ free in $\sigma$; local column sums are defined likewise.  We use:
\begin{itemize}
\item[(P1)] (stated, not reconstructed: \texttt{[1402] Win.bal\_win}; its hypotheses are read
off the call in \texttt{HDes.lean}.)  If at every matching $g$ with $|g|=e$ (under a further
hypothesis on $g$ that the proof of Lemma~\ref{lem:C-H} does not use) every local row sum
$\mathrm{lrow}_{\varphi_b}(g_b,i)$, $b=\beta_P(i)$, $i$ free, is $\equiv\varphi_b(g_b)$, and
likewise every local column sum, then $\mathrm{Bal}_p(\Pi,\Pi,e)$ ($m=3$).
\item[(P2)] The same with the weaker hypothesis that, for $g$ without cross edges,
$\bigl(\mathrm{lrow}_{\varphi_b}(g_b,i)-\varphi_b(g_b)\bigr)\prod_{c\ne b}\varphi_c(g_c)\equiv0$,
and likewise for columns.\lean{[1402] TExt.bal\_win'}  Proof: if $g$ has a cross edge, $\Pi$
vanishes at $g$ and at every one-edge extension of it; otherwise the free row sum of $\Pi$ at
$g$ is $\mathrm{lrow}_{\varphi_b}(g_b,i)\prod_{c\ne b}\varphi_c(g_c)$, since an edge from $i$ to
a hole of another block is a cross edge (stated, not reconstructed: \texttt{[1402]
ProdExt.line\_Prod}), while $\Pi(g)=\varphi_b(g_b)\prod_{c\ne b}\varphi_c(g_c)$.
\item[(P3)] (stated, not reconstructed: \texttt{[1434] ProdM.bal\_prodM}, in the form
described in the paragraph on the generators below.)  Say $\varphi_b$ is \emph{locally balanced
up to $E_b$} if the local line identities hold at every local matching of block $b$ with at
most $E_b$ edges.\lean{[1434] ProdM.LocBal (definition read off its uses in Towers.lean)}  If
each $\varphi_b$ is locally balanced up to $E_b$, vanishes on local matchings with fewer than
$s_b$ edges, and $e\le E_b+\sum_{c\ne b}s_c$ for every $b$, then $\mathrm{Bal}_p(\Pi,\Pi,e)$.
\end{itemize}
A \emph{length-only} factor is $\varphi(\sigma)=G(|\sigma|)$.  On a block with $P$ pigeons and
$Q$ holes its local row sum at $\sigma$ is exactly $(Q-|\sigma|)\,G(|\sigma|+1)$ and its local
column sum $(P-|\sigma|)\,G(|\sigma|+1)$, by counting free cells.\lean{[1402] TExt.gsym\_row, gsym\_col}
Local balance of a length-only factor is packaged as (stated, not reconstructed:
\texttt{[1434] LocLf.lf\_locbal}); in every application below its numerical hypotheses
are closed by \texttt{omega} and \texttt{decide}.

\begin{lemma}[the design $H$]\label{lem:C-H}
$H(\emptyset)=1$ and $\mathrm{Bal}_2(H,H,e)$ for $e\le2$.\lean{[1434] HDes.h\_design}
\end{lemma}

\begin{proof}
$H$ is the block product with $\beta_P(i)=[i\ge12]$, $\beta_H(j)=[j\ge4]$ and factors
$\varphi_0=B_7$, $\varphi_1=\mathrm{Dg}$ (the third slot of the Lean product is unused).  Here
$B_7(\sigma)=B(\rho(\sigma))$, where $B$ is the library design of $\mathrm{ontoFPHP}(4,8)$ and
$\rho(x_{ij})=4i+j$ renames an edge of block $0$ to the variable $x_{ij}$ of the $4$-hole
board; and $\mathrm{Dg}(\sigma)=1$ if every edge of $\sigma$ is one of $(12,4)$, $(13,5)$,
$(14,6)$, else $0$.\lean{[1434] HDes.H, B7, Dg, rho}  $H(\emptyset)=1$ is evaluated by
\texttt{decide}.  For the balance we apply (P1) and check the local identities at a matching
$g$ with $|g|=e\le2$.

\emph{Block $0$.}  $\rho$ is injective on block $0$, preserves pigeons, holes and conflicts,
and maps local matchings to matchings of $\mathrm{ontoFPHP}(4,8)$ of the same size.  So for a
pigeon $i<12$ free in $g$ the local row sum of $B_7$ at $g_0$ over the holes $t<4$ is the row
sum $\mathrm{row}^{(4,8)}_B(\rho(g_0),i)$, with $i$ free in $\rho(g_0)$; likewise for columns
$j<4$.\lean{[1434] HDes.row0, col0, red\_CM, red\_len, red\_pfree, red\_hfree}  Since
$|\rho(g_0)|\le2$, the library fact $\mathrm{Bal}^{(4,8)}_2(B,B,e')$ for $e'\le2$ (stated,
not reconstructed: \texttt{[1402] BDesMain.b\_design}) gives $\equiv B(\rho(g_0))=B_7(g_0)$.

\emph{Block $1$.}  For a pigeon $i\in\{12,13,14\}$ free in $g$, among the holes $4,5,6$ only
$t=i-8$ can give a nonzero term $\mathrm{Dg}(g_1+x_{it})$.  If $\mathrm{Dg}(g_1)=0$ all terms are
$0$; if $\mathrm{Dg}(g_1)=1$, then $g_1$ is diagonal, the hole $i-8$ is free (its diagonal
pigeon $i$ is), and the term is $1$.  So the local row sum equals $\mathrm{Dg}(g_1)$ exactly;
columns $j\in\{4,5,6\}$ likewise, with the single pigeon $j+8$.\lean{[1434] HDes.row1, col1}
\end{proof}

\begin{lemma}[the design $D$]\label{lem:C-D}
$D(\emptyset)=1$ and $\mathrm{Bal}_2(D,D,e)$ for $e\le3$.\lean{[1434] K5Bal.d\_design}
\end{lemma}

\begin{proof}
$D=\mathrm{Wt}(\mathit{sup})$ is the $\mathbb Z_{15}\times\mathbb Z_7$-invariant weight of the
library design of run 1627 (stated, not reconstructed: \texttt{[1627] K5Defs.Wt}
with its properties \texttt{Wt\_congr}, \texttt{Wt\_nm}); $\mathrm{Wt}$ depends only on the set of a list and vanishes on lists
that are not matchings.  $D(\emptyset)=1$ is the kernel-checked fact
\texttt{chk\_W0}.  Let $|g|\le3$.  By the two properties of $\mathrm{Wt}$, each term of the row
sum of $D$ at $g$ and pigeon $i$ equals $\mathrm{Wt}(x_{ij}::g)$, so
$\mathrm{row}_D(g,i)-D(g)$ is the pigeon defect $\mathrm{pigS}(\mathit{sup},i,g)$ of the run-1627
development, and likewise columns.\lean{[1434] K5Bal.term, row\_eq, col\_eq}  The library
lemma \texttt{[1627] K5Lem.lin\_all} (stated, not reconstructed), fed with the
kernel-checked certificates \texttt{chk\_all} and \texttt{chk\_empty} (\texttt{[1627] K5Chk}),
states that both defects are even for every $g$ of size $\le3$; the freeness of $i$ is not
used.\lean{[1434] K5Bal.d\_bal}
\end{proof}

\begin{lemma}[towers]\label{lem:C-towers}
The four block products of Table~\ref{tab:towers} satisfy $\mathrm{Bal}_2(\Phi,\Phi,e)$ for
every $e\le3$.\lean{[1434] Towers.towers\_bal, rbd\_bal, tb\_bal, rrm\_bal, rrrm\_bal}
\end{lemma}

\begin{table}[ht]
\centering\small
\begin{tabular}{llllll}
tower & block (pigeons $\times$ holes) & factor & $E_b$ & $s_b$\\\hline
$R_{Bd}$ & $\{0,\dots,11\}\times\{0,\dots,3\}$ & $B_7$ & $2$ & $0$\\
 & $\{12,13\}\times\{4,5\}$ & $[s\ge1]$ & $\infty$ & $1$\\
 & $\{14\}\times\{6\}$ & $1$ & $\infty$ & $0$\\\hline
$T_B$ & $\{0,\dots,11\}\times\{0,\dots,3\}$ & $B_7$ & $2$ & $0$\\
 & $\{12,13,14\}\times\{4,5,6\}$ & $[s\ge2]$ & $\infty$ & $2$\\\hline
$R_{RM}$ & $\{0,1\}^2$, $\{2,3\}^2$ & $[s\ge1]$ each & $\infty$ & $1$\\
 & $\{4,\dots,14\}\times\{4,5,6\}$ & $M$ & $1$ & $0$\\\hline
$R_{RRRM}$ & $\{0,1\}^2$, $\{2,3\}^2$, $\{4,5\}^2$ & $[s\ge1]$ each & $\infty$ & $1$\\
 & $\{6,\dots,14\}\times\{6\}$ & $1$ & $0$ & $0$
\end{tabular}
\caption{The towers of \texttt{Towers.lean} ($s$ = local size; $\infty$ is $100$ in Lean).}
\label{tab:towers}
\end{table}

\begin{proof}
Each tower is a block product, and we apply (P3) with the $E_b$ and $s_b$ of
Table~\ref{tab:towers}.\lean{[1434] Towers.RBd, TB, RRM, RRRM}
\emph{Vanishing below $s_b$}: $[s\ge1]$ at $s=0$ and $[s\ge2]$ at $s\le1$ are $0$ (\texttt{simp}).
\emph{Truncation}: $e\le E_b+\sum_{c\ne b}s_c$ for $e\le3$ and all $b$ is checked by
\texttt{decide}; the binding cases are $3\le2+1+0$ for $B_7$ in $R_{Bd}$, $3\le2+2$ for $B_7$
in $T_B$, $3\le1+1+1$ for $M$, and $3\le0+1+1+1$ for the all-ones block of $R_{RRRM}$.
\emph{Local balance}:
\begin{itemize}
\item Length-only factors, by \texttt{LocLf.lf\_locbal} with \texttt{decide}.  With the
local sums $(Q-s)G(s+1)$ and $(P-s)G(s+1)$ the identities read, mod $2$: for $[s\ge1]$ on
$2\times2$, $2\cdot1\equiv0$ at $s=0$ and $1\cdot1=1$ at $s=1$; for $[s\ge2]$ on $3\times3$,
$3\cdot0=0$, $2\cdot1\equiv0$, $1\cdot1=1$ at $s=0,1,2$; for the pin, $1\cdot1=1$ at $s=0$; for
the all-ones $9\times1$ block, $1\cdot1=1$ (rows) and $9\cdot1\equiv1$ (columns) at $s=0$.
At larger $s$ there is no free line in the block, except for rows of the $9\times1$ block at
$s=1$, where the identity would read $0\equiv1$; this is why $E_b=0$ there.
\item $B_7$ (on pigeons $<12$, holes $<4$) up to $E=2$: by the reindexing argument of
Lemma~\ref{lem:C-H} and \texttt{BDesMain.b\_design}.\lean{[1434] Towers.b7\_loc}
\item $M$ up to $E=1$: $M(\emptyset)=1$, $M(\{v\})=[v\in S_1]$, $M(\{v,w\})=[\{v,w\}\in S_2]$
for explicit lists $S_1$ ($11$ edges) and $S_2$, and $M=0$ on larger sets.  The local
identities at $\emptyset$ and at every single edge of the block are verified by kernel
computation (\texttt{decide +kernel}).\lean{[1434] Towers.M, MS1, MS2, M\_chk0, M\_chk1r, M\_chk1c, M\_loc}
\qedhere
\end{itemize}
\end{proof}

\begin{lemma}[rectangle pairs extend]\label{lem:C-rect}
Let $n\ge7$, $p\mid k$, and let $f=\rho\otimes\rho'$ be the product (\texttt{tprod},
defined in \texttt{RectBlock}) of the signed rectangles $\rho$ on $\{a,a'\}\times\{b,b'\}$ and $\rho'$ on
$\{c,c'\}\times\{d,d'\}$, with $a,a',c,c'$ distinct and $b,b',d,d'$ distinct.  Then there is $Y$
with $\mathrm{Bal}_p(Y,f,2)$.\lean{[1402] RectExt.rect\_ext, rect\_ext\_lab}
\end{lemma}

\begin{proof}
Choose fresh pigeons $x\ne y$ and fresh holes $h\ne h'$ (possible as $n\ge7$).\lean{[1402] RectExt.fresh}
Blocks: $0=\{a,a',x\}\times\{b,b',h\}$, $1=\{c,c',y\}\times\{d,d',h'\}$, and $2$ the rest, with
$n+k-6$ pigeons and $n-6\ge1$ holes.\lean{[1402] RectExt.bpf, count2}

\emph{Block $2$.}  There is $\mu$ on the block-$2$ cells with every row sum and every column sum
$\equiv1\pmod p$.  This holds for any $L_1$ pigeons and $L_2\ge1$ holes with $L_2\le L_1$ and
$p\mid L_1-L_2$ (here $L_1-L_2=k$), by induction on $L_2$: for $L_2=1$ take $\mu\equiv1$
(the column sum is $L_1\equiv L_2=1$); otherwise put $\mu=1$ at one cell $(s,q)$, $\mu=0$ on
the rest of row $s$ and column $q$, and take $\mu$ on the remaining $(L_1-1)\times(L_2-1)$
board by induction.\lean{[1402] RectExt.mu\_exists}  Let $M_2$ be $\mu$ on single edges and $0$
elsewhere; its local line sums at $\emptyset$ are $\equiv1$.

\emph{The extension.}  The library combines block data by (stated, not reconstructed:
\texttt{[1402] ProdExt.bal\_prod3}), applied with local degrees $1,1,0$.  Its data are, on
block $0$, the function $A_0$ equal to $\rho$ on single edges and $0$ elsewhere, and the
degree-$2$ function $Y_0=\mathrm{Yb}(a,a',x;b,b',h)$, a signed sum of three products of edge
functions (stated, not reconstructed: \texttt{[1402] RectBlock.Yb}); on block $1$ the same
for $\rho'$; on block $2$ the indicator $A_2$ of $\emptyset$ and $Y_2=M_2$.  The hypotheses
supplied are: the local line sums of $A_0$ at $\emptyset$ are $0$ (in row $a$,
$\rho(x_{ab})+\rho(x_{ab'})+\rho(x_{ah})=1-1+0$, and likewise elsewhere); at each block-$0$
edge $e$, every free local line sum of $Y_0$ at $\{e\}$ equals $\rho(e)$ exactly (stated, not
reconstructed: \texttt{[1402] RectBlock.r\_row, r\_col, Yb\_row, Yb\_col}); the same on block
$1$; and the line sums of $M_2$ at $\emptyset$ are $\equiv1=A_2(\emptyset)$.  The result is $Y$
with $\mathrm{Bal}_p(Y,\Pi_A,2)$, where $\Pi_A$ is the block product of $A_0,A_1,A_2$.

\emph{Agreement.}  On $2$-matchings $\Pi_A=\rho\otimes\rho'$: $\rho$ and $\rho'$ vanish off
blocks $0$ and $1$, and the remaining case split on the blocks of the two edges is closed by
\texttt{simp} and \texttt{grind}.\lean{[1402] RectExt.rect\_zero, tprod\_eq\_prod}  Since
$\mathrm{Bal}_p(Y,\cdot,2)$ depends only on residues on $2$-matchings
(Lemma~\ref{lem:ext-span}), $\mathrm{Bal}_p(Y,f,2)$.
\end{proof}

\begin{lemma}[the generator $T$ extends]\label{lem:C-T}
Let $n\ge4$ and $2\mid k$, and let $T$ be the indicator of the $2$-matchings inside the
block $\{0,1,2\}\times\{0,1,2\}$.  Then $\mathrm{Bal}_2(Y,T,2)$ for the block product $Y$
below.\lean{[1402] TExt.t\_ext}
\end{lemma}

\begin{proof}
Blocks: $0=\{0,1,2\}\times\{0,1,2\}$ and $2$ the remaining pigeons and holes (the middle
slot of the Lean product is empty, with factor $1$).  Factors: $\varphi_0(\sigma)=[|\sigma|\ge2]$;
$\varphi_2(\emptyset)=1$, $\varphi_2(\{x_{ij}\})=1$ if $i=j$ or ($i\ge n$ and $j=3$), else $0$,
and $\varphi_2=0$ on two or more edges.\lean{[1402] TExt.bt, G0, dd, dl, Y, T}  We show
$\mathrm{Bal}_2(Y,Y,2)$ by (P2); let $g$ have $|g|=2$ and no cross edge.\lean{[1402] TExt.Y\_bal}

\emph{Lines in block $0$.}  The bracket of (P2) is $(3-s)\varphi_0(s+1)-\varphi_0(s)$ with
$s=|g_0|\le2$ (three holes, resp.\ three pigeons, in the block), i.e.\ $0$, $2$, $0$ for
$s=0,1,2$: even, by \texttt{decide}.

\emph{Lines in block $2$.}  If $g_2=\emptyset$: a pigeon $i\ge3$ has exactly one block-$2$ hole
with $\varphi_2=1$, namely $i$ if $i<n$ and $3$ if $i\ge n$ (here $n\ge4$), so its local row
sum is $1=\varphi_2(\emptyset)$;\lean{[1402] TExt.row\_out} a hole $j\ge3$ has local column sum
$1+[j=3]\,k\equiv1$, using $2\mid k$.\lean{[1402] TExt.col\_out}  If $g_2\ne\emptyset$: the local
line sums of $\varphi_2$ at $g_2$ are $0$ ($\varphi_2$ vanishes on two edges), and the other
factor $\varphi_0(g_0)=0$ because $|g_0|\le1$, so the bracket times the other factor is $0$.\lean{[1402] TExt.loc\_big, G0\_small}

\emph{Agreement.}  On $2$-matchings $Y=T$: if both edges lie in block $0$, $Y=1\cdot1=T$;
otherwise either there is a cross edge, or an edge in block $2$ forces $|g_0|\le1$ and
$\varphi_0(g_0)=0$, and in both cases $Y=0=T$.\lean{[1402] TExt.Y\_eq}  Hence
$\mathrm{Bal}_2(Y,T,2)$ (Lemma~\ref{lem:ext-span}).
\end{proof}

Lemma~\ref{lem:C-H} is the design $H$ of Lemma~\ref{lem:C1}; Lemmas~\ref{lem:C-rect}
and~\ref{lem:C-T} are the extensions of $R_0$ and $T$ in Lemma~\ref{lem:C2}; and
Lemmas~\ref{lem:C-D} and~\ref{lem:C-towers} give the self-balanced functions $D$,
$R_{Bd}$, $R_{RM}$ and $T_B$ of Lemma~\ref{lem:C-span2b}.  The tower $R_{RRRM}$ (an
RRR-type product) is not cited in Lemma~\ref{lem:C-span2b}.


It remains to verify condition (iv) for $G_z=\mathcal S_3$; this is the level-$3$
certificate.

\begin{proposition}[level-3 zero-link span]\label{prop:C-z3}
Every $z$ on matchings of $\mathrm{ontoFPHP}(7,8)$ with $\mathrm{Bal}_2(z,0,2)$ lies
in $\mathrm{Span}_3(\mathcal G_3)$, where $\mathcal G_3$ is the set of valid RRR and
$T\otimes R$ generators defined below; $\mathcal G_3\subseteq\mathcal S_3$.\lean{[1434] Ext278d.z3\_of\_run, g3\_sb3; [0841] L3.Main.run\_all, cov\_all}
\end{proposition}

Given the proposition, Lemma~\ref{lem:C-lift} gives
$\mathrm{EXT}_2(7,8,3)$,\lean{[1434] Ext278c.ext3\_of\_z3, hard\_of\_z3} and with
Lemmas~\ref{lem:C0}--\ref{lem:C2} and Theorem~\ref{thm:ext-criterion} this proves
Theorem~C.\lean{[1434] Ext278d.hard\_of\_run}

\paragraph{The generators.}
Both families are block products: matchings with an edge between different blocks
get $0$, and otherwise the value is a product of local factors.  A \emph{length-only}
factor on a block with $P$ pigeons and $Q$ holes has a value $G(s)$ depending only
on the size $s$ of the local matching; it is locally balanced up to length $L$ when, for
every $s\le L$, $(Q-s)G(s+1)\equiv G(s)\pmod p$ if $s<P$ (a free pigeon exists) and
$(P-s)G(s+1)\equiv G(s)\pmod p$ if $s<Q$ (a free hole exists).\lean{[1434] LocLf.lf\_locbal}  A
product of locally balanced factors is balanced against itself at level $e$, with a
truncation rule: a factor need only be balanced up to some length $E_b$ provided it
vanishes below some length $s_b$ and $e\le E_b+\sum_{c\ne b}s_c$ for every
block~$b$.\lean{[1434] ProdM.bal\_prodM}
\begin{itemize}
\item $\mathrm{RRR}(\pi)$: three $2\times2$ rectangle blocks
$\{p_c,q_c\}\times\{j_c,h_c\}$ ($c<3$) and the all-ones block on the remaining $9$
pigeons $\times$ $1$ hole.\lean{[1434] Gen3.GRR}
\item $T\otimes R(\pi)$: the $T$ tower on a $3\times3$ block, a rectangle block, the
all-ones block on $9$ pigeons $\times$ the remaining hole, and a $1\times1$ pin.\lean{[1434] Gen3T.GTR}
\end{itemize}
For valid parameters (a decidable condition) both are self-balanced at every level
$e\le3$; in the Lean proof the local conditions and the truncation inequality are
discharged by \texttt{decide}.\lean{[1434] Gen3.grr\_bal, Gen3T.gtr\_bal}  They vanish on
matchings of size $<3$, so they have zero links at level $2$.\lean{[1434] Peel3.GS.zero2, GS.zl}
On a matching with $3$ edges each value is the number of occurrences of the
matching's code in an explicit list: the $64$ products of one edge from each
rectangle for RRR, and the $72$ products of a $2$-matching inside the $T$ block with a
rectangle edge for $T\otimes R$.\lean{[1434] Gen3.grr\_count, Gen3T.gtr\_count; Peel3.GS.count}

\paragraph{The certificate and what it asserts.}
Code the matchings with three edges, i.e.\ decreasing lists $[a,b,c]$ of variables
$<105$, by $r_3(a,b,c)=\binom a3+\binom b2+c<187460$ (injective).\lean{[1434] Code3.r3, r3\_inj}
A \emph{state} for $z$ is a pair $(r,P)$ of a residual function $r$ and a set $P$ of
\emph{finalised} codes such that
(S1) $\mathrm{Bal}_2(r,0,2)$; (S2) $z-r\in\mathrm{Span}_3(\mathcal G_3)$; (S3) $r$ is
even on every finalised matching.\lean{[1434] Peel3.State3, Inv3}
Initially $r=z$ and $P=\emptyset$.\lean{[1434] Peel3.state\_init}  Each step names a
pivot $[a,b,c]$, which must be a matching with three edges, and adds it to $P$:
\begin{itemize}
\item \texttt{k1}$(\mathit{gs};a,b,c)$: valid generators $\mathit{gs}$ such that every
code in their joint support list is the pivot or not finalised, and the pivot occurs
an odd number of times.  Replace $r$ by $r-r(\text{pivot})\sum\mathit{gs}$.  The new
residual is even at the pivot and unchanged mod $2$ at finalised codes (the
generators are $0$ there); (S1) and (S2) are preserved because the generators have
zero links and lie in $\mathcal G_3$.\lean{[1434] Peel3.chk1, kill\_sound}
\item \texttt{kx}: the same, with the parity condition checked through the XOR mask
of the support list, which allows cancellations among the
generators.\lean{[1434] Peel3.chkX; [0841] L3.Big.Part (splitting of large masks)}
\item \texttt{pr}$(u,w,\mathit{line};a,b,c)$: a free row or column at the
$2$-matching $[u,w]$ all of whose codes other than the pivot are finalised, the
pivot occurring an odd number of times.  By (S1) the line sum of $r$ is even, and
by (S3) all other terms are even, so $r(\text{pivot})$ is even; $r$ is
unchanged.\lean{[1434] Peel3.chkP, line\_par, step\_sound}
\end{itemize}
If after the run every matching with three edges is finalised, then $r\equiv0$ on all
of them and $z\equiv z-r$, so $z\in\mathrm{Span}_3(\mathcal G_3)$ by
(S2).\lean{[1434] Peel3.run3\_sound, span\_of\_cov, covA}

The Lean proof of Proposition~\ref{prop:C-z3} consists of the soundness theorem
just described, which is a conventional proof, and two finite facts, each
\emph{verified by kernel computation} (\texttt{decide +kernel}):
\begin{enumerate}
\item the checker accepts the step list: $\mathrm{run}_3(\emptyset,\mathit{ss})$ returns a
final mask $P_{\mathrm f}$.  The list is stored in files \texttt{L3/F0}--\texttt{F89};
their headers record $95{,}550$ steps in $2{,}154$ pieces, one kernel check per piece,
and the files contain $35{,}811$ \texttt{k1}, $1{,}394$ \texttt{kx} and $58{,}345$
\texttt{pr} steps.  The pieces are chained by kernel-checked equalities of the
intermediate masks;\lean{[0841] L3.F$i$.r$j$, L3.Main.lk0--lk89, run\_all; [1434] Peel3.run3\_append}
\item $P_{\mathrm f}$ covers every matching with three edges: for every $a<105$, every
$b<a$, $c<b$ with $[a,b,c]$ a matching, $r_3(a,b,c)\in P_{\mathrm f}$
($60$ kernel checks).\lean{[0841] L3.Cov.cov, L3.Main.cov\_all}
\end{enumerate}
We do not reproduce the certificate.  Its correctness rests on the kernel's
evaluation of these two facts; no part of the argument above depends on how the
step list was found.  (The finding page reports that a second agent re-checked the
certificate with an independent program; that check is outside the Lean proof.)

\paragraph{Mechanism of C.}
The same EXT criterion as in A and B, at $(p,k,n_0)=(2,8,7)$.  Levels $0$--$2$
follow the peeling pattern with fixed designs and small certificates.  Level $3$
reduces, by the lifting lemma, to one finite linear-algebra statement over
$\mathbb F_2$: the zero-link space on matchings with three edges is spanned by the
self-balanced RRR and $T\otimes R$ products.  A Gaussian-elimination-style
certificate with $95{,}550$ steps proves it.
\endgroup


% sec4-growth.tex -- Section 4 of the C3 paper (draft 2026-11-08): result D.
% Sources (run B, 2026-11-04): codex-20261104-x1-syzygies/03-FI-HOLE-AVOIDANCE.md and 04-NS-PC-PADDING.md
% (author x1, Codex); reviews notes-20261104-referee-x1-padding-c1.md (c1, Claude) and c2's review;
% adjudication codex-20261104-x2-padding-adjudication/01-ADJUDICATION.md (x2, Codex);
% cold referee notes-20261104/referee-runB-padding-transfer.md; Astra
% codex-20261104-runB-adjudication/00-ADJUDICATION.txt (C3) and 02-C3-SOURCE-AND-DEGREE.txt.
% Source excerpts: Rhoades Def 1.1 / Thm 3.12 (rhoades-source-excerpts.txt); Beame--Riis Def 8.2,
% Lemma 8.4, proof of Thm 8.1 (sources-20261101/beame-riis.txt).

\section{Growth in \texorpdfstring{$\ell$}{l}: the padding transfer}\label{sec:growth}

This section proves result D.  The mechanism is a transfer from Nullstellensatz designs to
PC lower bounds that costs $2D-2$ extra holes.  It is applied to the explicit designs of
Beame and Riis.

\begin{theorem}[padding transfer]\label{thm:padding}
Let $F$ be a field, $D\ge1$, $n_0\ge1$ and $k\ge0$.  If the board with $n_0+k$ pigeons and
$n_0$ holes has a unit degree-$D$ design over $F$, then for every
\[
  n\ \ge\ n_0+2D-2
\]
the board with $n+k$ pigeons and $n$ holes satisfies $\mathrm{ND}_D$ and $\mathrm{U}_D$.  In
particular $\mathrm{onto\text{-}FPHP}^{n+k}_n$ has no $\mathrm{PC}_F$ refutation of degree at
most $D$.
\end{theorem}

\begin{corollary}[D]\label{cor:growth}
Let $p$ be prime, $\ell\ge1$ and $F$ a field of characteristic $p$.  For
\[
  n\ \ge\ \frac{(p+2)^\ell-p^\ell}{2}+2^{\ell+1}-4,
\]
every $\mathrm{PC}_F$ refutation of $\mathrm{onto\text{-}FPHP}^{n+p^\ell}_n$ has degree at
least $2^\ell$.
\end{corollary}

\paragraph{Tier and inputs.}  T-WRITTEN.  The proof uses three published results that we do
not re-prove: Corollary~B of Church and Ellenberg~\cite{churchellenberg}, Theorem~3.12 of
Rhoades~\cite{rhoades2024}, and the design construction of Beame and
Riis~\cite[Def.~8.2, Lemma~8.4, proof of Thm.~8.1]{beame1998counting}.  It uses no
representation theory beyond these, and no result of Harman or of Hone et al.  The input we
consider most exposed is the conversion of Rhoades's theorem into filtered completeness on
the square (Lemma~\ref{lem:square-complete}); it is short, and we give it in full.
Section~\ref{sec:growth-provenance} records who wrote and checked each step.

\subsection{External inputs}\label{sec:growth-inputs}

\paragraph{FI-modules.}  An \emph{FI-module} assigns an abelian group $V(T)$ to each finite
set $T$ and a map $V(T)\to V(T')$ to each injection $T\hookrightarrow T'$, functorially.  It
is \emph{generated in degree $\le d$} if every $V(T)$ is spanned by the images of the
$V(T_0)$ with $|T_0|\le d$.  For a module $W$ that assigns a group to each finite set and
an action of bijections only, the \emph{induced} (``free'') FI-module is
\[
  M(W)(T)=\bigoplus_{T_0\subseteq T}W(T_0),
\]
with injections acting by relabelling~\cite[\S2.1, eq.~(2)]{churchellenberg}.  ``Free''
here does not mean projective, and no restriction on characteristic is involved.

\begin{theorem}[Church--Ellenberg~{\cite[Cor.~B]{churchellenberg}}]\label{thm:ce}
Let $M$ be an induced FI-module generated in degree $\le a$ and $V$ an FI-module generated
in degree $\le b$.  Then the kernel of every FI-module map $V\to M$ is generated in degree
$\le a+b+1$.
\end{theorem}

The theorem is stated for abelian groups.  We apply it to $F$-vector spaces by forgetting
the scalar structure; the images of smaller pieces are $F$-subspaces, so generation as
groups gives spanning over $F$.

\paragraph{The square.}  On an $s\times s$ board let $I_s\subseteq F[x]$ be generated by all
row sums $\sum_jx_{ij}$, all column sums $\sum_ix_{ij}$, and all products $x_{ij}x_{ij'}$,
$x_{ij}x_{i'j}$ of two variables in one row or one column, squares
included~\cite[Def.~1.1]{rhoades2024}.  Let $I(\mathcal P_s)$ be the ideal of polynomials
vanishing on the $s!$ permutation matrices, and $\operatorname{gr}I(\mathcal P_s)$ the ideal
of top-degree homogeneous parts of its elements.

\begin{theorem}[Rhoades~{\cite[Thm.~3.12]{rhoades2024}}]\label{thm:rhoades}
Over every field, $I_s=\operatorname{gr}I(\mathcal P_s)$.
\end{theorem}

Rhoades's later results on the module structure assume characteristic $0$ or $p>s$; we use
only Theorem~\ref{thm:rhoades}, which has no such assumption.

\paragraph{Beame--Riis designs.}  For $\ell\ge0$ put $N_\ell=((p+2)^\ell-p^\ell)/2$.
Beame and Riis construct, by induction on $\ell$, a $(2^\ell-1)$-design over
$\mathbb Z_p$ on the board with $N_\ell+p^\ell$ pigeons and $N_\ell$
holes~\cite[proof of Thm.~8.1]{beame1998counting}.  The step is their Lemma~8.4: a
$d$-design on a $D_0\times R_0$ board gives a $(2d+1)$-design on a board with
$(p+1)|D_0|+|R_0|$ pigeons and $|D_0|+(p+1)|R_0|$ holes.  Their $d$-designs
(Definition~8.2) are exactly our unit degree-$d$ designs with values in $\mathbb F_p$.  We
use the construction, not Theorem~8.1 itself; this avoids converting between their
multiplier degree and our product degree.

\subsection{Square filtered completeness}\label{sec:growth-square}

\begin{lemma}\label{lem:square-complete}
On an $s\times s$ board, for every $t$, $W_t$ is the set of elements of $F_t$ that vanish
at every permutation matrix.  Consequently every linear functional on $F_t$ that vanishes on
$W_t$ is the restriction of a linear combination, with coefficients in $F$, of evaluations at
permutation matrices.
\end{lemma}

\begin{proof}
Elements of $W_t$ vanish on permutation matrices, since each row or column relation does.
Conversely let $g\in F_t$ vanish on permutation matrices; lift it to a multilinear
polynomial $f$ of degree $\le t$ supported on matchings.  Then $f\in I(\mathcal P_s)$.  By
Theorem~\ref{thm:rhoades} its top homogeneous part, of degree $t'\le t$, is a combination
$\sum_r h_r\gamma_r$ of the generators $\gamma_r$ of $I_s$ with homogeneous $h_r$ and
$\deg(h_r\gamma_r)=t'$.  Let $\tilde\gamma_r$ be the corresponding axiom: $\gamma_r-1$
(a pigeon or hole axiom) for a row or column sum, $x^2-x$ (a Boolean axiom) for a square
$x^2$, and $\gamma_r$ itself (a collision or functional axiom) for a product of two distinct
variables in one line.  Each $\tilde\gamma_r$ has top part $\gamma_r$.  Then $f-\sum_rh_r\tilde\gamma_r$ has degree $<t'$ and still
vanishes on permutation matrices.  Repeating, $f$ is a combination of the axioms with
multipliers of degree at most $t-1$ for the linear axioms and $t-2$ for the quadratic ones.
Applying $\rho$, the Boolean, collision and functional terms vanish and each linear term becomes a combination of
row and column relations at bases of size $\le t-1$, so $g=\rho(f)\in W_t$.

For the second statement: a functional on $F_t$ vanishing on $W_t$ factors through the
evaluation map $F_t\to F^{\mathcal P_s}$, whose kernel is $W_t$.  Extend the induced
functional from the image to all of $F^{\mathcal P_s}$; every functional on
$F^{\mathcal P_s}$ is a combination of point evaluations.
\end{proof}

The functional obtained in this way need not be invariant, nonnegative or unital.  Every
evaluation at a permutation matrix $w$, viewed as the function $M\mapsto[M\subseteq w]$ on
matchings of every size, satisfies all row and column relations at every base: if
$M\subseteq w$ and $i$ is free in $M$, then $w(i)$ is free in $M$, so exactly one term of
the free row sum is $1$.  So do their linear combinations.

\subsection{Balanced arrays have small support}\label{sec:growth-support}

\begin{lemma}\label{lem:small-support}
Let $e\ge1$.  Every $f\in H_e(P,S)$ is a sum of arrays $f_\alpha\in H_e(P,S)$, each
supported on matchings inside $P_\alpha\times S_\alpha$ with $|P_\alpha|\le2e$ and
$|S_\alpha|\le2e$.
\end{lemma}

\begin{proof}
Fix the pigeon set $P$ and regard everything as a functor of the hole set.  Let $C_e(T)$ be
the space of arrays on $e$-matchings $P\leftrightarrow T$ all of whose free column sums
vanish.  Each column condition involves a single $e$-element hole set
$\mathrm{holes}(N)\cup\{j\}$, so
$C_e(T)=\bigoplus_{T_0\subseteq T,\,|T_0|=e}U(T_0)$, where $U(T_0)$ is the space of
column-balanced arrays using exactly the holes $T_0$.  With injections acting by
relabelling and extension by zero, $C_e$ is the induced module $M(U)$, generated in degree
$e$.  Let $B_{e-1}(T)$ have a basis of pairs $(N,i)$, $N$ an $(e-1)$-matching
$P\leftrightarrow T$ and $i$ a pigeon free in $N$; it is induced, generated in degree
$e-1$, since $(N,i)$ uses exactly the holes of $N$.  The row-sum map
$(Rf)(N,i)=\sum_{j}f(N\cup\{(i,j)\})$ is a map of FI-modules $C_e\to B_{e-1}$: after
extension by zero, the new terms vanish, and coordinates $(N,i)$ with $N$ outside the old
hole set get the value $0$ on both sides.  Its kernel is $H_e(P,\cdot)$.  By
Theorem~\ref{thm:ce} with $a=e-1$, $b=e$, the kernel is generated in degree
$\le 2e$: every $f\in H_e(P,S)$ is a sum of zero extensions of balanced arrays on hole sets
of size $\le2e$.  Extension by zero preserves both marginal conditions.

Apply the same argument with pigeons and holes exchanged to each summand, keeping its hole
set fixed.  This splits each summand further into balanced arrays on at most $2e$ pigeons,
without enlarging its hole set.
\end{proof}

\subsection{Extending one balanced layer}\label{sec:growth-tensor}

\begin{lemma}[square realisation]\label{lem:square-real}
Let $|A|=|B|=2e$ and $f\in H_e(A,B)$.  There is a function $\mu$ on the matchings of the
$A\times B$ board (of all sizes) that satisfies every row and column relation, vanishes on
matchings of size $<e$, and equals $f$ on matchings of size $e$.
\end{lemma}

\begin{proof}
The functional on $F_e$ that is $0$ below level $e$ and $f$ at level $e$ vanishes on $W_e$:
relations at bases of size $<e-1$ involve only zeros, and those at bases of size $e-1$ are
the vanishing marginals of $f$.  Lemma~\ref{lem:square-complete} and the remark after it
give $\mu$.
\end{proof}

\begin{lemma}[tensor extension]\label{lem:tensor}
Let $1\le e\le D-1$.  Split the board into pigeon sets $A\sqcup A'$ and hole sets
$B\sqcup B'$ with $|A|=|B|=2e$.  Let $\mu$ be as in Lemma~\ref{lem:square-real} on
$A\times B$ and let $\nu$ be a unit degree-$(D-e)$ design on $A'\times B'$.  For a matching
$M$ with $|M|\le D$ put
\[
  G(M)=\begin{cases}
  0 & \text{if $M$ has an edge in $A\times B'$ or $A'\times B$,}\\
  0 & \text{if $|M\cap(A\times B)|<e$,}\\
  \mu(M\cap(A\times B))\,\nu(M\cap(A'\times B')) & \text{otherwise.}
  \end{cases}
\]
Then $G$, as a functional on $F_D$, vanishes on $W_D$ and on all matchings of size $<e$,
and $G(M)=f(M)$ for matchings $M\subseteq A\times B$ of size $e$ (zero on other matchings of
size $e$).
\end{lemma}

\begin{proof}
$G$ is defined: in the third case the residual part has size at most $D-e$.  At level $e$
the residual part of a non-zero term is empty and $\nu(\emptyset)=1$.

Fix a base $M$ with $|M|\le D-1$ and a free pigeon $i$; columns are symmetric.  Write
$a=|M\cap(A\times B)|$ and $b=|M\cap(A'\times B')|$.  We check that the free row sum of $G$
at $(M,i)$ equals $G(M)$.
\begin{itemize}
\item $M$ has a crossing edge: every extension has one too; both sides are $0$.
\item $i\in A$, $a\le e-2$: extensions into $B$ have core size $\le e-1$ and extensions into
  $B'$ cross; both sides are $0$.
\item $i\in A$, $a=e-1$: $G(M)=0$.  Only extensions into $B$ can be non-zero, and their sum is
  $\nu(M\cap(A'\times B'))\cdot\sum_{j\in B\text{ free}}\mu(M_{\mathrm{core}}\cup\{(i,j)\})
  =\nu(\cdot)\,\mu(M_{\mathrm{core}})=0$, by the row relation of $\mu$ and $\mu=0$ below $e$.
  Here $b\le D-1-a=D-e$, so $\nu$ is evaluated within its degree.  This is the only place
  where the cutoff at level $e$ matters.
\item $i\in A$, $a\ge e$: the same row relation of $\mu$ gives the sum $G(M)$.
\item $i\in A'$, $a<e$: every extension into $B'$ still has core size $<e$, and extensions
  into $B$ cross; both sides are $0$.
\item $i\in A'$, $a\ge e$: then $b\le D-1-e$, so the row relation of $\nu$ at the base
  $M\cap(A'\times B')$ is one of its defining relations; extensions into $B$ cross; the sum
  is $\mu(M_{\mathrm{core}})\,\nu(M\cap(A'\times B'))=G(M)$.
\end{itemize}
So $G$ satisfies every relation at bases of size $\le D-1$, that is, $G(W_D)=0$.
\end{proof}

\begin{lemma}[padding a design]\label{lem:pad}
A unit degree-$D$ design on the board $(n_0+k)\times n_0$ gives one on $(r+k)\times r$ for
every $r\ge n_0$, and, by restriction, unit designs of every smaller degree.
\end{lemma}

\begin{proof}
Add $r-n_0$ pigeons and $r-n_0$ holes and fix a bijection $\sigma$ between them.  Define the
new design to be the old one times $[M_{\mathrm{new}}\subseteq\sigma]$ on matchings without
edges between old and new parts, and $0$ on others.  The proof of Lemma~\ref{lem:tensor}
applies with the old board as residual and $\sigma$ as core; there is no cutoff, and the
value at $\emptyset$ stays $1$.
\end{proof}

\subsection{Proof of Theorem~\ref{thm:padding}}\label{sec:growth-proof}

Let $n\ge n_0+2D-2$, and let $U$ be the unit degree-$D$ design on $(n+k)\times n$ given by
Lemma~\ref{lem:pad}.  Then $\mathrm{U}_D$ holds.

For $\mathrm{ND}_D$ we show that every functional $\lambda$ on $F_{D-1}$ vanishing on
$W_{D-1}$ extends to a functional on $F_D$ vanishing on $W_D$
(Section~\ref{sec:defs-written}).  Start with $\Lambda_0=\lambda(1)\,U$; the remainder
$\lambda-\Lambda_0|_{F_{D-1}}$ vanishes on $W_{D-1}$ and at $\emptyset$.  Suppose that for
some $1\le e\le D-1$ the remainder $r_e$, a functional on $F_{D-1}$ vanishing on $W_{D-1}$,
vanishes on all matchings of size $<e$.  Its level-$e$ layer $f_e$ lies in $H_e$ by the last
paragraph of Section~\ref{sec:defs-written} with $t=D-1$.  By
Lemma~\ref{lem:small-support}, $f_e=\sum_\alpha f_\alpha$ with each $f_\alpha$ balanced on at
most $2e$ pigeons and $2e$ holes.  Since $n\ge 2D-1>2e$, enlarge each support to exactly
$2e$ pigeons and $2e$ holes; the array stays balanced.  The complement has $n-2e\ge n_0$
holes and surplus $k$, so by Lemma~\ref{lem:pad} it has a unit degree-$(D-e)$ design.
Lemmas~\ref{lem:square-real} and~\ref{lem:tensor} give $G_\alpha$ vanishing on $W_D$ and
below level $e$, with level-$e$ layer $f_\alpha$.  Put
$\Lambda_e=\Lambda_{e-1}+\sum_\alpha G_\alpha$.  The new remainder
$r_{e+1}=r_e-\sum_\alpha G_\alpha|_{F_{D-1}}$ vanishes below level $e+1$ and still vanishes
on $W_{D-1}$.  After the step $e=D-1$ the remainder is zero, so $\Lambda_{D-1}$ extends
$\lambda$ and vanishes on $W_D$.  (For $D=1$ there are no correction steps.)

The PC conclusion is Proposition~\ref{prop:nd-pc}.  \qed

\subsection{Proof of Corollary~\ref{cor:growth}}\label{sec:growth-cor}

Take $D=2^\ell-1$, $n_0=N_\ell$ and $k=p^\ell$.  Beame and Riis's construction gives a unit
degree-$D$ design over $\mathbb F_p$ on $(n_0+k)\times n_0$; its defining conditions are
linear with coefficients in $\mathbb F_p$, so it is also a design over every field $F$ of
characteristic $p$.  For $\ell\ge1$, $n_0=N_\ell\ge N_1=1$.  Theorem~\ref{thm:padding}
excludes refutations of degree at most $2^\ell-1$ whenever
$n\ge N_\ell+2(2^\ell-1)-2=N_\ell+2^{\ell+1}-4$.  \qed

\begin{corollary}[exact degree at $p=2$, without E's inputs]\label{cor:growth-p2}
Let $\ell\ge1$.  Over every field of characteristic $2$, the PC degree of
$\mathrm{onto\text{-}FPHP}^{n+2^\ell}_n$ is exactly $2^\ell$ for
$n\ge(4^\ell-2^\ell)/2+2^{\ell+1}-4$.
\end{corollary}

\begin{proof}
At $p=2$ the bound $2^\ell$ of Corollary~\ref{cor:growth} equals the upper bound $p^\ell$ of
Proposition~\ref{prop:riis}.
\end{proof}

So at $p=2$ the exact degree of Corollary~\ref{cor:exact} has another proof, at a larger
threshold ($n\ge40$ instead of $15$ for $\ell=3$), using neither of E's
representation-theoretic inputs; both proofs use Riis's upper bound.  This corollary was
deduced while drafting the paper and checked in the model review of the paper, not in the
run that produced D.  For odd $p$ the two bounds differ, because Beame and Riis's designs stop at
degree $2^\ell-1$ while the Nullstellensatz degree may be as large as $p^\ell-1$; designs of
degree $p^\ell-1$, if they exist, would give the exact degree for every $p$ through this
section alone (Section~\ref{sec:open}).

For $\ell=1$ Corollary~\ref{cor:growth} gives degree at least $2$ for $k=p$ and every $n\ge1$, consistent with
Theorem~A.  For $p=2$, $\ell=2$ it gives degree at least $4$ for surplus $4$ from $n\ge10$;
the Lean result gives exactly $4$ from $n\ge5$.

\paragraph{The Nullstellensatz route, for comparison.}  Theorem~8.1 of
\cite{beame1998counting} states its bound in multiplier degree.  It also yields the
design used above: a certificate of product degree at most $D$, reduced modulo the Boolean
axioms, has every non-zero multiplier of degree at most $D-1$, since every axiom has positive
degree; with $D=2^\ell-1$ this contradicts Theorem~8.1, and duality gives the unit design.
The two routes give the same threshold; we cite the construction because it needs no
convention conversion.

\subsection{What was checked}\label{sec:growth-provenance}

The proof of Theorem~\ref{thm:padding} was written by a Codex agent (x1) in run 20261104-0759
(\texttt{03-FI-HOLE-AVOIDANCE.md}, \texttt{04-NS-PC-PADDING.md}).  Two Claude agents reviewed
it in the run (c1 re-did every case of Lemma~\ref{lem:tensor} by hand); a second Codex agent
adjudicated it against the primary sources of Church--Ellenberg and Rhoades; a separate
model referee with no contact with the run re-read both sources and found the transfer
consistent with its own small computations over $\mathbb F_2$ and $\mathbb F_3$.  The
adjudicator Astra checked the Beame--Riis construction in the author-hosted PDF
(Definitions~7.1, 7.2, 8.2, Lemma~8.4, proof of Theorem~8.1) and the conversion above.
The tensor construction of Lemma~\ref{lem:tensor} was also checked by exact computation at
$D=3$ over $\mathbb F_2$, with supports of size $2$ and $4$, on the full $12\times8$ board
(all $60{,}884$ free-line equations per case).  None of this is a Lean proof or a human
review, and none of it re-proves Corollary~B of \cite{churchellenberg}, Theorem~3.12 of
\cite{rhoades2024} or Lemma~8.4 of \cite{beame1998counting}.

% sec5-periodicity.tex -- Section 5 of the C3 paper (draft 2026-11-08): result E.
% Sources (run B 20261104-0759): claude-20261104-c3-closure/01-PC-CLOSURE-PERIODICITY.md (c1: closure formulation,
% tilting proof, R1-R3, computations); codex-20261104-c3-monoidal/05-CONSOLIDATED-EXACT-PC-PERIODICITY.md (x2: primary
% proof), 03 (equality-boundary counterexample), 04 (exact-quotient repair); codex-20261104-x1-adjudication/01;
% notes-20261104/referee-runB-closure-periodicity.md (cold referee); Astra codex-20261104-runB-adjudication/00 (C1) and 03.
% Hone et al. statements checked against the archived arXiv:2509.07377v2 text (Def 6.1.1, Lemma 7.3.2, Thm 8.4.5,
% Thm 8.5.3, Lemma 8.5.1) on 2026-11-08.

\section{Periodicity and the exact degree}\label{sec:periodicity}

This section proves result E: whether $\mathrm{onto\text{-}FPHP}^{m}_n$ has a PC refutation of
degree at most $D$ depends, for $m$ large compared with $D$, only on $m$ modulo the least
power of $p$ above $D$.  Since the square board $m=n$ is satisfiable, this fixes the PC degree
exactly.

\begin{theorem}[periodicity]\label{thm:periodicity}
Let $F$ be a field of characteristic $p$, $n\ge1$, $D\ge1$, and $q$ a power of $p$ with
$q>D$.  Let $m\equiv m'\pmod q$.
\begin{enumerate}
\item[(a)] If $m,m'>2D$, then $\mathrm{onto\text{-}FPHP}^{m}_n$ has a $\mathrm{PC}_F$
  refutation of degree at most $D$ if and only if $\mathrm{onto\text{-}FPHP}^{m'}_n$ has one.
\item[(b)] If Theorem~8.4.5 of \cite{hone2025} holds (see below), the same holds for
  $m,m'\ge2D$.
\end{enumerate}
\end{theorem}

Here $m$ is the number of pigeons and $n$ the number of holes; $n$ is arbitrary, and $m<n$ is
allowed.

\begin{corollary}[E, exact degree]\label{cor:exact}
Let $p$ be prime, $k\ge1$, $n\ge1$ and $q=p^{v_p(k)}$.  Over every field of characteristic
$p$, the PC degree of $\mathrm{onto\text{-}FPHP}^{n+k}_n$ is exactly $q$ for every
$n\ge2q-1$.  If Theorem~8.4.5 of \cite{hone2025} holds, the same holds for $n\ge2q-2$.
\end{corollary}

\begin{proof}
If $q=1$ the degree is $1$ (Section~\ref{sec:intro-question}).  Let $q\ge p$ and $D=q-1$.
The board with $n$ pigeons and $n$ holes has a solution, a permutation matrix, so it has no
refutation of any degree.  Since $n\equiv n+k\pmod q$ and $n>2D$ when $n\ge2q-1$,
Theorem~\ref{thm:periodicity}(a) gives no refutation of degree at most $q-1$ for $n+k$
pigeons.  Proposition~\ref{prop:riis} gives one of degree $q$.  With (b), $n\ge2D=2q-2$
suffices.
\end{proof}

In particular, for surplus $p^\ell$ the PC degree is exactly $p^\ell$ once $n\ge2p^\ell-1$:
Riis's upper bound is attained, and the degree grows with $\ell$ as fast as it can.  This is
the strongest statement in this paper.  It is also the one that rests on the deepest
inputs, so we state its status with care.

\paragraph{Tier and inputs.}  T-WRITTEN.  Part (a) is proved in Appendix~\ref{app:tilting}
(outlined in Section~\ref{sec:periodicity-tilting}), using Harman's
equivalences~\cite{harman2015} with the repairs of Section~\ref{sec:errata}, and standard facts
on tilting modules (tensor products and Levi restrictions of tilting modules are tilting;
bounded tilting complexes compute the derived category).  The input we consider most exposed
there is the compatibility of Harman's equivalence with induction and restriction of
arbitrary submodules (Lemma~\ref{lem:B}), whose tilting-theoretic foundations we cite and do
not re-prove.

Section~\ref{sec:periodicity-oti} gives a shorter argument through the semisimplification
functors of Hone, Klein, Pauwels, Sherman, Yacobi and Zhang~\cite{hone2025}, which covers (b)
as well.  It uses their Theorem~8.4.5.  A model reconstruction of the proof of that theorem
(Astra, November 2026) found that the published proof does not establish it: their
Lemma~8.4.4 rests on Lemma~8.2.1, which is false as stated (Section~\ref{sec:errata}).  This is
not a counterexample to Theorem~8.4.5, and the reconstruction found every other step of
Section~\ref{sec:periodicity-oti}, and their Lemma~7.3.2, to hold.  Until that step is
repaired, (b) and the threshold $n\ge2q-2$ are conditional, and (a) rests on
Appendix~\ref{app:tilting}.

\paragraph{Comparison.}  For $q=p$, Theorem~A gives the exact degree $p$ at the Lean tier
from $n\ge2p-3$.  For $q=p^2$, result F gives the exact degree $p^2$ from $n\ge2p^2-3$ for
every prime, on different inputs (Section~\ref{sec:bands}).  Result D gives $2^\ell$, which is
exact when $p=2$, on the inputs and thresholds of Section~\ref{sec:growth}.
Below the threshold the degree can be smaller: over $\mathbb F_2$ with $k=4$, exact
computation of the closure below gives degree $3$ at $n=3$ and $n=4$, while Lean gives
exactly $4$ from $n=5$ (\texttt{[1156] StdPCExact.std\_f2\_k4\_exact}), below the threshold
$2q-1=7$ of Corollary~\ref{cor:exact}.

\subsection{The closure}\label{sec:periodicity-closure}

We use the notation of Section~\ref{sec:defs-written} on the board with $m$ pigeons and $n$
holes, and write $\mathcal B_m$, $F_d(m)$.  Let $C_0(m)\subseteq F_D(m)$ be the span of the
reduced pigeon and hole axioms, and recursively
\[
  S_t=C_t\cap F_{D-1},\qquad
  C_{t+1}=C_t+\sum_{i,j}\rho(x_{ij})\,S_t .
\]
The chain is increasing in the finite-dimensional space $F_D(m)$, so it stabilises; let
$C_\infty(m)$ be its limit.  All $C_t$ are invariant under permutations of the pigeons.

\begin{lemma}\label{lem:closure}
$\mathrm{onto\text{-}FPHP}^m_n$ has a $\mathrm{PC}_F$ refutation of degree at most $D$ if and
only if $1\in C_\infty(m)$.
\end{lemma}

\begin{proof}
($\Rightarrow$) As in Proposition~\ref{prop:nd-pc}, every line $f$ of a degree-$D$
derivation has $\rho(f)\in C_\infty$: axioms reduce into $C_0$ or to $0$, linear combinations
stay in $C_\infty$, and if $xf$ is a legal multiplication then $\deg f\le D-1$, so
$\rho(f)\in C_\infty\cap F_{D-1}=S_t$ for large $t$ and $\rho(xf)\in C_{t+1}$.

($\Leftarrow$) We show by induction on $t$ that every element of $C_t$, written as a
multilinear matching polynomial, is a line of a degree-$D$ derivation.  For $t=0$ this is
clear.  Let $s\in S_t$; by induction $s$ is derivable, and $\deg s\le D-1$ as a polynomial,
so $x_{ij}s$ is a legal multiplication.  Its reduction $\rho(x_{ij}s)$ differs from
$x_{ij}s$ by multiples of Boolean, collision and functional axioms: reduce one repeated
variable or forbidden pair at a time; each multiple has degree at most $\deg(x_{ij}s)\le D$,
as in Lemma~\ref{lem:ml-reduce}.  Linear combinations give $C_{t+1}$.
\end{proof}

On the square board ($m=n$), every element of $C_\infty$ vanishes at every permutation
matrix (true for $C_0$ and preserved by the recursion), so $1\notin C_\infty(n)$.

The recursion is a finite sequence of linear-algebra operations on matrices with entries in
$\mathbb F_p$.  Extending scalars to any field containing $\mathbb F_p$ commutes with sums,
images and intersections, so $1\in C_\infty$ over $F$ if and only if $1\in C_\infty$ over
$\mathbb F_p$ if and only if over $\overline{\mathbb F}_p$.  We may therefore work over
$k=\overline{\mathbb F}_p$.

\paragraph{The closure as representation theory.}  The symmetric group $S_m$ acts on
pigeons, and all spaces above are $kS_m$-modules.  For $e\ge0$ and a fixed set of $e$ holes
listed in increasing order, the $e$-matchings onto those holes form the permutation module
$E_e(m)$ of ordered $e$-tuples of distinct pigeons; so $F_d(m)$ is a direct sum of
$\binom ne$ copies of $E_e(m)$, $e\le d$.  For a hole $j$ let $U_j=E_1(m)$, the natural
permutation module on pigeons, and let $\mu_j\colon U_j\otimes F_{D-1}(m)\to F_D(m)$ send
$u_i\otimes x_M$ to $\rho(x_{ij}x_M)$.  Then
\[
  C_{t+1}=C_t+\textstyle\sum_{j}\mu_j\bigl(U_j\otimes S_t\bigr),
\]
and the constant $1$ is the image of the map $k\to F_D(m)$ onto the empty matching.  So
$C_\infty(m)$ is built from permutation modules by tensoring with $U_j$, images,
intersections and sums, and the question is whether the composite
$k\to F_D(m)\to F_D(m)/C_\infty(m)$ is zero.

\subsection{A conditional proof through semisimplification}\label{sec:periodicity-oti}

This argument is shorter than the one in Appendix~\ref{app:tilting} and covers part (b), but
it depends on \cite[Thm.~8.4.5]{hone2025}, whose published proof is incomplete
(Section~\ref{sec:errata}).  We give it because every other step has been checked, so a
repair of that one step would make (b) unconditional.

It suffices to prove one step: $M=L+q$ pigeons against $L$ pigeons, with $L\ge2D$ (for (b))
or $L>2D$ (for (a)).  Applying the step repeatedly in either direction gives the theorem.

\paragraph{The functor.}  Let $H\subseteq S_q$ be an elementary abelian $p$-group of order
$q$ acting regularly on the last $q$ pigeons.  Let $\Phi$ be restriction from $S_M$ to
$S_L\times H$ followed by semisimplification in the $H$-variable: an $H$-module
decomposes into indecomposables, and semisimplification discards the summands of dimension
divisible by $p$ (the \emph{negligible} ones).  This is the OTI functor of \cite{hone2025}
for the semisimplification functor $\varphi_H$, which satisfies the hypotheses of their
Theorem~8.4.5 (stated there).  $\Phi$ is strong symmetric monoidal \cite[\S2.4]{hone2025}.  It
is \emph{not} exact on all $kS_M$-modules.

\paragraph{External input.}  Call $\lambda\vdash M$ \emph{$q$-stable} if
$\sum_{i\ge2}\lambda_i<q$ and, when $M<3q$, $\lambda_1\ge(M+q)/2$
\cite[Def.~6.1.1]{hone2025}.  Let $\mathrm{Pol}^{\mathrm{st}}_{M,q}$ be the Serre subcategory
of polynomial functors of degree $M$ whose simples are the $L_{\lambda^t}$ with $\lambda$
$q$-stable, and $\mathrm{Pol}^{\mathrm{tr}}_{L,q}$ its truncated counterpart in degree $L$, with
simples $L_{(\lambda-qe_1)^t}$.

\begin{theorem}[{\cite[Thm.~8.4.5]{hone2025}}]\label{thm:hone}
For $\varphi$ satisfying the two hypotheses there (semisimplification does), the OTI
functor corestricts to an equivalence of highest-weight categories
$\mathrm{Pol}^{\mathrm{st}}_{M,q}\to\mathrm{Pol}^{\mathrm{tr}}_{L,q}$.
\end{theorem}

By \cite[Lemma~7.3.2]{hone2025}, this functor commutes with Schur--Weyl duality: for the
Schur functor $\mathcal F$, $(\mathcal F\boxtimes1)\Phi\simeq\Phi\mathcal F$.  We use the
twisted Schur functor $\tilde{\mathcal F}=\mathcal F(-)\otimes\mathrm{sgn}$; twisting costs
nothing on the $H$-side, since $H$ is a $p$-group and its sign character is trivial in
characteristic $p$.

\paragraph{Step 1: exactness on the bounded-tail range.}  Let $\mathcal A_D(M)$ be the Serre
subcategory of polynomial functors whose simple labels are $\lambda^t$ with
$\lambda=(M-e,\tau)$, $|\tau|=e\le D$, and $\mathcal A_D(L)$ the same with $L$.  These
$\lambda$ are $q$-stable: $e\le D<q$, and $M-e\ge(M+q)/2$ is equivalent to $L\ge2e$.
Theorem~\ref{thm:hone} restricts to an exact equivalence
$\Theta\colon\mathcal A_D(M)\to\mathcal A_D(L)$ preserving $e$ and $\tau$.  The kernel
$K_D(M)$ of $\tilde{\mathcal F}$ on $\mathcal A_D(M)$ is generated by the simples whose
$\lambda$ is $p$-singular (some part repeated $p$ times).  Since
$M-e=L+q-e>e\ge\tau_1$, a repetition in $\lambda$ lies inside $\tau$ and remains one in
$(L-e,\tau)$, where $L-e\ge e\ge\tau_1$.  Hence $\Theta(K_D(M))\subseteq K_D(L)$.  (For
$L>2D$ this is an equality; at $L=2D$ it can be strict.)  The exact functor
$\tilde{\mathcal F}\Theta$ therefore factors through the Serre quotient
$\mathcal A_D(M)/K_D(M)$.

This quotient is the full image $\mathcal R_D(M)$ of $\tilde{\mathcal F}$ on
$\mathcal A_D(M)$, a full subcategory of $kS_M$-modules whose kernels, cokernels and
subobjects are the ordinary ones: for $X\in\mathcal A_D(M)$ and a submodule
$T\le\tilde{\mathcal F}X$, the pullback $X_T=X\times_{G\tilde{\mathcal F}X}GT$ along the
unit of the right adjoint $G$ of $\tilde{\mathcal F}$ is a subobject of $X$ with
$\tilde{\mathcal F}X_T=T$; lifting the graph of a module map in the same way identifies
morphisms.  (Compare \cite[Lemma~8.5.1]{hone2025}.)  By Lemma~7.3.2 the induced functor
$\mathcal R_D(M)\to\mathcal R_D(L)$ is $\Phi$.  So $\Phi$ is exact on $\mathcal R_D(M)$.  It
need not be faithful there, and we never use faithfulness.

\paragraph{Step 2: every object used lies in the range.}  The module $E_e(M)$ is
$\tilde{\mathcal F}$ of $\Lambda^{M-e}V\otimes V^{\otimes e}$, whose highest weight is
$(M-e,1^e)^t$; every composition factor has a label $\nu$ dominated by it, so $\nu^t$
dominates $(M-e,1^e)$, has first part $\ge M-e$ and tail of size $\le e$.  Hence $E_e(M)$,
$e\le D$, and so $F_d(M)$, $d\le D$, lie in $\mathcal R_D(M)$.  On a fixed hole set,
$U_j\otimes E_e\cong E_{e+1}\oplus E_e^{\oplus e}$ (the extra pigeon is new, or equals the
pigeon at one of the $e$ holes), so $U_j\otimes F_{D-1}$ lies in the range too, and with it
$U_j\otimes S$ for every submodule $S\le F_{D-1}$.  All intersections, images and sums in the
recursion therefore take place in $\mathcal R_D(M)$.

\paragraph{Step 3: the maps are the obvious ones.}  For a permutation module $k[X]$,
restricted to $H$, the $H$-orbits of size $>1$ give summands whose every $H$-endomorphism has
constant diagonal on the orbit and hence trace $0$; they are negligible and are killed.  So
$\Phi(k[X])=k[X^H]$, and an equivariant map between permutation modules goes to its block
from fixed basis vectors to fixed basis vectors.  A matching fixed by $H$ uses none of the
last $q$ pigeons (an edge at a moved pigeon would move), so $\Phi(F_d(M))=F_d(L)$ and
$\Phi(U_j(M))=U_j(L)$, compatibly with inclusions and with the monoidal structure.  On fixed
basis vectors: a pigeon axiom of a remaining pigeon is unchanged; a hole axiom
$\sum_{i}x_{ij}-1$ loses exactly the terms at moved pigeons, which form a non-trivial orbit,
and becomes the hole axiom of the smaller board; and $\mu_j$ sends $u_i\otimes x_N$ to the
same matching, the same new matching, or $0$, on both boards.

\paragraph{Step 4: the closure transfers.}  By induction on $t$, $\Phi$ carries
$C_t(M)\le F_D(M)$ to $C_t(L)\le F_D(L)$: exactness transfers the intersection with
$F_{D-1}$, the images and the sums; monoidality and Step~3 transfer
$\mu_j(U_j\otimes S_t)$.  Both chains stabilise after finitely many steps.

\paragraph{Step 5: detecting the constant.}  Let $c\colon k\to F_D(M)/C_\infty(M)$ be the
image of the constant.  If $c\ne0$ it is injective, since $k$ is simple; $\Phi$ is exact, so
$\Phi(c)$ is injective, and $\Phi(k)=k\ne0$, so $\Phi(c)\ne0$.  If $c=0$ then $\Phi(c)=0$.
By Step~4, $\Phi(c)$ is the corresponding map for $L$ pigeons.  So $1\in C_\infty(M)$ if and
only if $1\in C_\infty(L)$, and Lemma~\ref{lem:closure} finishes the step.  \qed

\subsection{The proof of part (a), through Harman's equivalences and tilting modules}\label{sec:periodicity-tilting}

For $m,m'>2D$, Theorem~\ref{thm:periodicity}(a) is proved without \cite{hone2025}; the proof is
given in full in Appendix~\ref{app:tilting}.  It transports the closure
along Harman's equivalence between the bounded-tail image categories of $S_m$ and $S_{m'}$ in weight
$\le D$~\cite[Thm.~2.9]{harman2015}.  The multiplication step is written as induction and
restriction of submodules, so the argument needs the equivalence to commute with induction
and restriction of arbitrary submodules, not only of permutation modules.  This is proved on
the Schur-algebra side, where induction and restriction become tensoring with the natural
module and a weight space of Levi restriction, both exact and preserving tilting modules; the
compatibility holds on tilting modules because the twisted Schur functor is fully faithful
there, and extends to all modules through $K^b(\mathrm{Tilt})\simeq D^b$.  The full
faithfulness is proved directly in every characteristic, because the identification Harman's
proof uses for it is false at $p=2$; the period must be a power of $p$ strictly greater than
$D$; and the strict inequality $m,m'>2D$ is needed (Section~\ref{sec:errata}).

\subsection{Checks}\label{sec:periodicity-checks}

The statements were tested by exact computation of $C_\infty$ (full row reduction over
$\mathbb F_p$, no use of symmetry).  For $D=2$ the predicate ``$1\in C_\infty$'' is periodic in
$m$ with period $4$ for $p=2$ ($m=5,\dots,12$), $3$ for $p=3$ ($m=5,\dots,13$) and $5$ for
$p=5$ ($m=5,\dots,14$), for $n=3,4$, and also for $n>m$.  Period $2$ fails for $p=2$, $D=2$
($n=3$: refutable at $m=5$, not at $m=7$), so the condition $q>D$ cannot be weakened to
$q\ge D$.  At the boundary $m=2D$ over $\mathbb F_2$ ($D=1,2,3$, $n\le7$) every congruent
triple agreed.  The mechanism itself was tested: for $q=p$ and $H$ generated by a $p$-cycle,
the Jordan type of the generator on $C_\infty(M)$ and on $F_D(M)/C_\infty(M)$ has no blocks of
size strictly between $1$ and $p$, and the number of blocks of size $1$ equals the dimension
of the corresponding space for $L=M-p$ pigeons (for example $(p,M,n,D)=(3,9,3,2)$:
$\dim C_\infty(9)=199$, with $88$ blocks of size $1$, and $\dim C_\infty(6)=88$).  These are
finite checks, not evidence for the general statement beyond the absence of a counterexample
at these sizes.

\paragraph{Who did what.}  All in run 20261104-0759; all reviews are model reviews.  The
closure formulation (Lemma~\ref{lem:closure}), the computations of
Section~\ref{sec:periodicity-checks} and the proof of Appendix~\ref{app:tilting} are by a Claude
agent (c1).  A Claude referee with no contact with the run reviewed that proof and found the
characteristic-$2$ error in \cite{harman2015}; a Codex agent (x2) reviewed it and supplied the
two-way scalar extension.  The proof of Section~\ref{sec:periodicity-oti} is by x2, who also found
the counterexample to \cite[Thm.~8.5.3]{hone2025} at the boundary and the exact-quotient argument
that avoids it; c1 reviewed it, and a second Codex agent (x1) adjudicated it, including the
boundary range.  After the run, a model referee re-read \cite{harman2015} and \cite{hone2025}
and re-ran the closure computations with its own code, and Astra adjudicated both proofs.  Astra
then reconstructed the proof of \cite[Thm.~8.4.5]{hone2025} and Lemma~7.3.2 from the source:
Lemma~7.3.2 holds, the proof of Theorem~8.4.5 has the gap described in
Section~\ref{sec:errata}, and Steps~1--5 hold conditional on that theorem.  A second model
re-checked the counterexample to Lemma~8.2.1 by hand.

% sec6-bands.tex -- Section 6 of the C3 paper (draft 2026-11-08): result F and the F_3 target.
% Sources (run B 20261104-0759): claude-20261104-c2-direct-nd/01-DIRECT-ND.md (c2: Lemmas 1-5, (A), (I), data) and
% 02/04/06/07 (consequences by band); codex-20261104-x1-syzygies/03 (L'), 08 (S(4) independent check), 09, 10, 20-22
% (bands); x2 12-14; notes-20261104/referee-runB-direct-route.md (cold referee); Astra
% codex-20261104-runB-adjudication/00 (C2), 01-C2-GLOBAL-AUDIT.txt, 05 (S(4) numbers). Lean: K100 (run 1627).

\section{Fixed-degree bands and the \texorpdfstring{$\mathbb F_3$}{F3} target}\label{sec:bands}

This section gives a third mechanism, which works degree by degree.  It reduces
$\mathrm{ND}_d$ to a statement about the homogeneous ideals generated by row sums and by
column sums, and proves that statement for every $d<p^2$.  It cannot give growth in $\ell$ at
fixed $p$, and it is not a third proof of it; what it gives is exactness at
$v_p(k)\le2$ from a slightly lower threshold, and the $\mathbb F_3$ target, on inputs that
differ from those of Sections~\ref{sec:growth} and~\ref{sec:periodicity}.

\begin{theorem}[F]\label{thm:bands}
Let $p$ be prime, $k\ge1$, $q=p^{v_p(k)}$, and $1\le D<\min(q,p^2)$.  Over every field of
characteristic $p$, $\mathrm{onto\text{-}FPHP}^{n+k}_n$ has no PC refutation of degree at
most $D$ for any $n\ge2D-1$.  Consequently, if $v_p(k)=2$, its PC degree is exactly $p^2$
for every $n\ge2p^2-3$.
\end{theorem}

\begin{corollary}[the $\mathbb F_3$ target]\label{cor:f3}
If $9\mid k$ and $n\ge7$, $\mathrm{onto\text{-}FPHP}^{n+k}_n$ has no PC refutation of degree
at most $4$ over $\mathbb F_3$.
\end{corollary}

The corollary is the case $p=3$, $D=4$.  It has two proofs: through the identity (I) below in
degrees $d\le4$, which for $p=3$ needs only the first two bands and the connectivity theorem;
and through exact computations on three square boards (Section~\ref{sec:bands-f3}).

\paragraph{Tier and inputs.}  T-WRITTEN.  The ideal-identity proof uses the connectivity
theorem of Bj\"orner, Lov\'asz, Vre\'cica and \v Zivaljevi\'c for chessboard complexes, as stated by
Shareshian and Wachs~\cite[Thm.~1.1]{shareshianwachs}.  Degrees $d<2p$ need nothing else.
Degrees $2p\le d<p^2$ (odd $p$) also use Yal\c c\i n's theorem on equivariant Moore spaces and
the Dade group~\cite[Def.~1.1, Props.~3.8, 5.7]{yalcin2016}, Proposition~3.3 of
Miller~\cite{miller2024}, and uniqueness of Dade caps.  The input we consider most exposed is
Miller's Proposition~3.3 at its interface with the truncated chessboard cochain complex,
where the output needed is an equality of integer dimensions of Brauer quotients
(Section~\ref{sec:bands-high}).  The second proof of Corollary~\ref{cor:f3} replaces
connectivity by Church--Ellenberg's Corollary~B and exact computation.

\subsection{From ND to a kernel statement}\label{sec:bands-reduction}

We use the notation of Section~\ref{sec:defs-written} on a board with $m$ pigeons and $n$
holes.  Let $V_d$ be the span of matchings of size exactly $d$.  The \emph{labels} at level
$d$ are the pairs $(N,\mathrm{row}\ i)$ and $(N,\mathrm{col}\ j)$ with $|N|=d-1$ and $i$, $j$
free in $N$.  Let $a_d$ send a label to the top part of its relation
($\sum_jx_{N\cup\{(i,j)\}}$, resp.\ $\sum_ix_{N\cup\{(i,j)\}}$), a map from the label space
to $V_d$, and let $\beta$ send a label to its base $x_N$.

\begin{lemma}\label{lem:nd-ker}
$W_d\cap F_{d-1}=W_{d-1}+\beta(\ker a_d)$.  Hence $\mathrm{ND}_d$ holds if and only if
$\beta(\ker a_d)\subseteq W_{d-1}$.
\end{lemma}

\begin{proof}
An element of $W_d$ is $w+\sum_gc_g\,\mathrm{rel}(g)$ with $w\in W_{d-1}$ and $g$ over the
level-$d$ labels.  Its degree-$d$ part is $a_d(c)$, so it lies in $F_{d-1}$ iff
$c\in\ker a_d$, and then it equals $w-\beta(c)$.
\end{proof}

Fix total orders on pigeons and holes.  The following elements of the label space lie in
$\ker a_d$:
\begin{itemize}
\item $rr(N;a,b)=\sum_j(N\cup\{(b,j)\},\mathrm{row}\ a)-\sum_j(N\cup\{(a,j)\},\mathrm{row}\ b)$
  for $|N|=d-2$ and free pigeons $a<b$; $cc(N;a,b)$ likewise for holes;
\item $rc(N;i,j)=\sum_{i'\ne i}(N\cup\{(i',j)\},\mathrm{row}\ i)-\sum_{j'\ne j}(N\cup\{(i,j')\},\mathrm{col}\ j)$;
\item the \emph{anchored} elements $Z^{(t)}_N$, $1\le t\le d$, $|N|=d-t$:
  \[
  Z^{(t)}_N=\sum_E\Bigl(\sum_{\substack{i\ \text{free}\\ i>\max\mathrm{pig}(E)}}(N\cup E,\mathrm{row}\ i)
     -\sum_{\substack{j\ \text{free}\\ j>\max\mathrm{hol}(E)}}(N\cup E,\mathrm{col}\ j)\Bigr),
  \]
  the outer sum over $(t-1)$-matchings $E$ on the pigeons and holes free in $N$
  ($\max\emptyset=-\infty$).  Both halves map under $a_d$ to the sum of all $d$-matchings
  containing $N$, each once.
\end{itemize}

\begin{lemma}\label{lem:lower-images}
$\beta(rr)$, $\beta(cc)$, $\beta(rc)$ lie in $W_{d-1}$, and with $s=|N|$,
\[
  \beta(Z^{(t)}_N)\equiv\Bigl(\tbinom{m-s}{t}-\tbinom{n-s}{t}\Bigr)x_N\pmod{W_{d-1}}.
\]
\end{lemma}

\begin{proof}
$\beta(rr(N;a,b))$ is the difference of the row relations at $(N,b)$ and $(N,a)$, and
similarly for $cc$; $\beta(rc(N;i,j))$ is the column relation at $(N,j)$ minus the row
relation at $(N,i)$.  For $Z$, group the $E$ by their pigeon set $A$, $|A|=t-1$; the row
coefficient is the number of free pigeons above $\max A$.  For fixed $A$,
$\sum_{\mathrm{pig}(E)=A}x_{N\cup E}\equiv x_N$ modulo $W_{d-1}$, by applying one row
relation at a time at bases of size $\le d-2$.  Since $t$-subsets of the free pigeons
correspond to pairs ($A$, a new largest element), the row half contributes
$\binom{m-s}{t}x_N$; the column half contributes $\binom{n-s}{t}x_N$.
\end{proof}

If $m-n=k$ and $t<q=p^{v_p(k)}$, then $\binom{m-s}t\equiv\binom{n-s}t\pmod p$, by
$(1+z)^k=(1+z^q)^{k/q}\equiv1+O(z^q)$ over $\mathbb F_p$.  So:

\begin{proposition}\label{prop:A-to-nd}
Let $D<q$ and suppose that for every $d\le D$
\begin{center}
(A)$(m,n,d)$: \ $\ker a_d$ is spanned by the elements $rr$, $cc$, $rc$ and $Z^{(t)}_N$.
\end{center}
Then $\mathrm{ND}_1,\dots,\mathrm{ND}_D$ hold on the board, and since $W_0=0$ also
$\mathrm{U}_D$; by Proposition~\ref{prop:nd-pc} there is no PC refutation of degree at most
$D$.
\end{proposition}

(A) is homogeneous and does not see $k$; the surplus enters only through the binomial
differences.  (A) fails off the stable range, for example at $(p,m,n,d)=(3,7,4,3)$, where
$\ker a_3$ has dimension $944$ and the listed elements span $916$ (computed twice,
independently).

\subsection{The ideal identity}\label{sec:bands-identity}

Let $\mathcal A$ be the homogeneous matching ring: basis $x_M$ over matchings,
$x_Mx_{M'}=x_{M\cup M'}$ if $M,M'$ are edge-disjoint and $M\cup M'$ is a matching, and $0$
otherwise.  (So $x_{ij}^2=0$; this is the top-degree ring, not the Boolean algebra
$\mathcal B$.)  Let $R_i=\sum_jx_{ij}$, $C_j=\sum_ix_{ij}$, $I=(R_i)$, $J=(C_j)$, and
$\sigma_t$ the sum of all $t$-matchings.  The identity in degree $d$ is
\[
  \text{(I)}(m,n,d):\qquad (I\cap J)_d=(IJ)_d+\sigma_1\mathcal A_{d-1}+\sigma_p\mathcal A_{d-p},
\]
with the last term absent for $d<p$.

\begin{lemma}\label{lem:I-to-A}
If $\min(m,n)\ge2d-1$, then (I)$(m,n,d)$ implies (A)$(m,n,d)$.
\end{lemma}

\begin{proof}
To $s\in\ker a_d$ with row part $\alpha$ and column part $\gamma$ associate
$y(s)=\sum\alpha_iR_i=-\sum\gamma_jC_j\in(I\cap J)_d$ (the products are taken in
$\mathcal A$, with the base of each label as coefficient).  Then $y(rr)=y(cc)=0$,
$y(rc(N;i,j))=x_NR_iC_j$ (because $R_ix_{ij}=0$), and $y(Z^{(t)}_N)=x_N\sigma_t$.  So every
element on the right-hand side of (I) is $y$ of a combination of $rc$'s and $Z$'s: $(IJ)_d$ is
spanned by the $x_NR_iC_j$, which vanish unless $i$ and $j$ are free in $N$, and
$\sigma_1\mathcal A_{d-1}$, $\sigma_p\mathcal A_{d-p}$ by the $x_N\sigma_1$, $x_N\sigma_p$.
(The run's first form of (I) listed all $x_N\sigma_t$; the form above, with $\sigma_1$ and
$\sigma_p$ only, is the one proved.)  By (I), subtracting suitable $rc$'s and $Z$'s leaves $s'$ with $y(s')=0$, that is
$\sum\alpha'_iR_i=0$ and $\sum\gamma'_jC_j=0$ separately.  For a fixed set of $d$ pigeons
the row-only relation is a cocycle of the augmented cochain complex of the chessboard
complex $\Delta(d,n)$ in dimension $d-2$, and the $rr$'s are its coboundaries.  For
$n\ge2d-1$ the connectivity theorem~\cite[Thm.~1.1]{shareshianwachs} gives
$\tilde H_{d-2}=\tilde H_{d-3}=0$ over $\mathbb Z$, hence $\tilde H^{d-2}=0$ over every field.
So $\alpha'$ is a sum of $rr$'s, and likewise $\gamma'$ of $cc$'s (using $m\ge2d-1$).
\end{proof}

\paragraph{Degrees below $p$.}  Let $G=S_d$ act on ordered $d$-tuples by positions.  The
coefficient spaces of $(I\cap J)_d$ identify with $V^G$, where $V=P_m\otimes P_n$ and $P_m$ is
the span of functions on injective $d$-tuples in $[m]$ that are independent of at least one
position; this identification needs no division.  On the spanning tensors of deletion
indicators, the norm $N_G=\sum_{g\in G}g$ gives $x_N\sigma_1$ when the row and column
indicators omit the same position and $x_NR_iC_j$ when they omit different positions.  Hence
\[
  N_GV=(IJ)_d+\sigma_1\mathcal A_{d-1},\qquad
  (I\cap J)_d\big/\bigl((IJ)_d+\sigma_1\mathcal A_{d-1}\bigr)\cong\hat H^0(G,V).
\]
For $d<p$, $|G|=d!$ is invertible, the Tate group vanishes, and (I) holds.

\subsection{The bands \texorpdfstring{$p\le d<p^2$}{p <= d < p^2}}\label{sec:bands-high}

We give the argument in these bands in outline.  Each step was written out in full in the
run, reviewed there, and reconstructed in Astra's audit of the local-to-global argument; the
outline below follows that audit.  It is the part of the paper whose written form is furthest
from what is printed here.
For $p\le d<p^2$ write $d=ap+r$ with $0\le r<p$; then $a<p$, $v_p(d!)=a$, and
$Q=(C_p)^a$ acting on $a$ disjoint blocks of $p$ positions is a Sylow $p$-subgroup of $S_d$.
Since $[S_d:Q]$ is prime to $p$, restriction to $Q$ followed by transfer is multiplication
by a unit on $\hat H^0$, so it suffices to show that $\hat H^0(Q,V)$ is generated by
transfers of classes that become $(p-1)!\,x_N\sigma_p$ in $\mathcal A$.
\begin{enumerate}
\item \emph{One cycle} ($a=1$, any number $r$ of fixed positions).  Projecting the augmented
  chessboard cochains of $\Delta(d,m)$ to the fixed-position cochains of $\Delta(r,m)$ has a
  kernel of free modules; connectivity gives exactness below the top; dimension shifting along
  the cyclic period shows that the actual map from functions on fixed-position $r$-tuples is
  onto on $\hat H^0$.  This covers $p\le d<2p$ for every prime, including $p=2$, $d=2,3$.
\item \emph{Full-vertex summands.}  The fixed complexes of subgroups of $Q$ are again
  chessboard complexes, tight Moore complexes by connectivity; Yal\c c\i n's theorem identifies
  $P_m$ as capped endo-permutation with an explicit Dade class.  For the tensor product
  $V=P_m\otimes P_n$ the relative terms cancel, giving $[V]=2[\Omega_Q]$; its cap is the
  minimal second syzygy $C=\ker(R^a\to R)$, $R=\mathbb F_p[x_1,\dots,x_a]/(x_i^p)$.  Dade cap
  uniqueness makes every full-vertex indecomposable summand of $V$ a copy of $C$, and $\hat H^0(Q,C)$ is
  generated by transfers from the cyclic factors.
\item \emph{Proper-vertex summands.}  Miller's Proposition~3.3, applied to the truncated
  chessboard cochain complexes, gives for every non-trivial $L\le Q$ the integer identity
  $\dim\mathrm{End}(V)(L)=(h_s(m)h_s(n))^2$ for the Brauer quotients, where $h_s$ is the top
  cohomology dimension of the $L$-fixed chessboard complex.  With Higman's criterion this forces
  every proper vertex into a coordinate hyperplane $H_i$, so those summands contribute only
  transfers from the $H_i$.
\item \emph{Induction on $a$.}  Restricted to $H_i$, $V$ is the same kind of module with $a-1$
  moving blocks and $r+p$ fixed positions; the one-cycle step allows any number of fixed
  positions, so induction on $a$ reduces everything to single $p$-cycles, where the classes are
  fixed-label indicators whose transfer to $S_d$ is $(p-1)!\,x_N\sigma_p$.
\end{enumerate}
At $d\ge p^2$ the subgroup $Q$ is no longer Sylow and the argument stops; nothing here
concerns that range.

\begin{proof}[Proof of Theorem~\ref{thm:bands}]
For $d<\min(q,p^2)$ and $m,n\ge2d-1$, (I) holds (above), hence (A)
(Lemma~\ref{lem:I-to-A}), hence $\mathrm{ND}_d$ (Lemmas~\ref{lem:nd-ker},
\ref{lem:lower-images}); Proposition~\ref{prop:A-to-nd} excludes refutations of degree
$\le D$ once $n\ge2D-1$.  All spaces are defined over $\mathbb F_p$, so the statement holds
over every field of characteristic $p$.  If $v_p(k)=2$, take $D=p^2-1$; the upper bound is
Proposition~\ref{prop:riis}.
\end{proof}

\subsection{The \texorpdfstring{$\mathbb F_3$}{F3} target by finite certificates}\label{sec:bands-f3}

The second proof of Corollary~\ref{cor:f3} avoids the ideal identity.  Two lemmas move (A)
from squares to rectangles.

\begin{lemma}[hole avoidance; from Corollary~B of \cite{churchellenberg}]\label{lem:hole-avoid}
If $n\ge2e+1$, every array in $H_e$ is a sum of arrays in $H_e$ each of which avoids some
hole.  Dually, if the restriction of $v\in V_e$ to every board with one hole removed lies in
the image of $a_e$ there, then $v$ lies in the image of $a_e$.
\end{lemma}

This is the case $|S|>2e$ of Lemma~\ref{lem:small-support}.  The threshold is sharp: exact
computation shows that the statement fails at $n=2e$ on several small boards.

\begin{lemma}[pigeon induction]\label{lem:pigeon-ind}
Order the pigeons so that the removed pigeon $\omega$ is the largest.  Then
(A)$(m-1,n,d)$, (A)$(m-1,n-1,d-1)$ and Lemma~\ref{lem:hole-avoid} on $(m-1,n)$ at level $d-1$
imply (A)$(m,n,d)$.  The mirror statement removes a hole.
\end{lemma}

\begin{proof}
Split the labels of the board $(m,n)$ into those avoiding $\omega$ ($G_1$), the row labels of
$\omega$ ($G_2$), and those whose base uses $\omega$ ($G_3$), and decompose
$V_d(m)=V_d(m-1)\oplus\bigoplus_hx_{\omega h}V_{d-1}(m-1,n\setminus h)$.  Let
$s=(g_1,g_2,g_3)\in\ker a_d$.  (a) The first component of $a_d(s)$ is $a_d(g_1)$ on the board
$m-1$, so $g_1$ is in that kernel; each listed element of the smaller board is the
$G_1$-part of the element of the same name on the larger board (the extra terms involve
$\omega$ only in $G_2$, $G_3$, and $\omega$ largest preserves the order conditions of $Z$).
Subtract.  (b) If $g_1=0$, the component at $h$ reads
$g_2|_{\text{avoid }h}+a_{d-1}(g_{3,h})=0$, so by Lemma~\ref{lem:hole-avoid} $g_2$ lies in
the image of $a_{d-1}$ on $(m-1,n)$; the elements $rr(N;a,\omega)$ and $rc(N;\omega,j)$ have
$G_1$-part $0$ and $G_2$-parts exactly the needed extension sums.  Subtract.  (c) If
$g_1=g_2=0$, each $g_{3,h}$ lies in $\ker a_{d-1}$ on $(m-1,n\setminus h)$, and
$x_{\omega h}$ times a listed element at base $N'$ is the listed element of the same name at
base $N'\cup\{(\omega,h)\}$.
\end{proof}

Starting from (A) at level $1$ (which holds on every board), Lemma~\ref{lem:pigeon-ind} and
its mirror give (A) on all boards with $m,n\ge2d-1$ from the single square statement
$S(d)=$ (A)$(2d-1,2d-1,d)$.  For $p=3$ and $D=4$ the squares are $3\times3$, $5\times5$ and
$7\times7$.  $S(2)$, $S(3)$ and $S(4)$ were verified over $\mathbb F_2$ and $\mathbb F_3$ by
exact rank computations, by three implementations written by different agents; at $S(4)$ the
top matrix has rank $27{,}068$ on $58{,}800$ labels, and the $47{,}972$ listed elements span its
$31{,}732$-dimensional kernel.  With Proposition~\ref{prop:A-to-nd} this proves
Corollary~\ref{cor:f3} at tier T-WRITTEN with T-MEASURED certificates.

\paragraph{Lean.}  The library contains a conditional Lean theorem for the same target
(published as finding K100, DOI 10.5281/zenodo.24056819): if $9\mid k$, $n\ge7$, and the
spanning statement (A)$(n+k,n,r)$ holds for $r=2,3,4$, then there is no standard PC
refutation over $\mathbb F_3$ of degree at most $4$; (A) at level $1$ is proved on every board with positive numbers of pigeons and holes.\lean{\texttt{Programme.\guillemotleft 20261104-1627-proof-complexity\guillemotright}:
RemainingSpanning.no\_stdpc\_le; level 1: MatchingLabels.spanning\_one; the
semantic form ConcreteDegreeFour.no\_ppc}
The Lean theorem does not formalise Lemma~\ref{lem:pigeon-ind} or
Lemma~\ref{lem:hole-avoid}, so it needs (A) on every board $(n+k)\times n$ rather than on
three squares.  On the smallest board with positive surplus divisible by $9$, $16\times7$, (A) is verified by exact
computation at levels $2$ (kernel dimension $343$) and $3$ (kernel dimension $24{,}780$, two
independent programs); level $4$ was not decided within the time limit.

\paragraph{Who did what.}  All in run 20261104-0759 unless stated; all reviews are model
reviews.  The reduction (Lemmas~\ref{lem:nd-ker}--\ref{lem:I-to-A},
Proposition~\ref{prop:A-to-nd}, Lemma~\ref{lem:pigeon-ind}) and the form of (I) are by a Claude
agent (c2), with a Codex agent (x1), and were reviewed by a second Claude agent (c1).  The
hole-avoidance lemma is x1's, checked by c2 and adjudicated by a second Codex agent (x2).  The
identity (I) was proved band by band: in degree $2$ and whenever $d!$ is invertible by x1
(reviewed by c2); at $d=p$ by x2 (reviewed by x1 and c2); for $p<d<2p$ by x1 (reviewed by c2,
accepted by x2); in the higher bands by x1 and x2 together, with c2 reviewing each step.  The
finite certificates $S(2)$--$S(4)$ were computed by c1 and, independently, by x1, and
regenerated after the run by Astra with a third implementation.  A model referee with no
contact with the run re-derived Lemmas~1--5 and (I) for $d<2p$, re-read the sources of the
higher bands, and checked (I) and the spanning statement on ten boards with its own code.
Astra's global audit reconstructed the local-to-global argument of
Section~\ref{sec:bands-high}.  The Lean theorem is from run 20261104-1627 and was re-proved by a
second agent there.

% sec7-errata.tex -- Section 7 of the C3 paper (draft 2026-11-08): result G.
% Sources: codex-20261104-runB-adjudication/03-C1-INTERFACES-AND-ERRATA.txt \S5 (wording limits);
% notes-20261104/referee-runB-closure-periodicity.md; claude-20261104-c3-closure/01 \S4 R1, \S6;
% codex-20261104-c3-monoidal/03-BOUNDARY-AND-STRICT-REPAIR.md.

\section{Errata in two sources}\label{sec:errata}

Writing the proofs of Section~\ref{sec:periodicity} required checking two recent sources
closely.  We record five points.  For each we say whether it is a counterexample or an
omitted justification, and we do not draw conclusions beyond the stated point; in
particular none of them suggests that the other results of either paper fail.  The
first four points were sent on 8 November 2026 to Harman and to Yacobi, one of the authors
of \cite{hone2025}.  The fifth was found afterwards, in the reconstruction of
\cite[Thm.~8.4.5]{hone2025} described in Section~\ref{sec:periodicity}.

\paragraph{Harman, Corollary~1.8(4), at $p=2$ (counterexample).}  The corollary identifies
the value of a right adjoint on a Young module with a tilting module,
$\tilde G(Y(\lambda))=T(\lambda')$~\cite[Thm.~1.7(4), Cor.~1.8(4)]{harman2015}.  At $p=2$ in
polynomial degree $2$, $Y(2)$ is the trivial module $k$; the right-adjoint value has
dimension $N(N+1)/2$, while $T(1,1)=\Lambda^2V$ has dimension $N(N-1)/2$ (at $N=2$: $3$ against
$1$).  Harman's proof of his Theorem~2.8 uses this identification to conclude that the
twisted Schur functor is fully faithful on tilting modules.  That conclusion is true in every
characteristic: the tilting modules in question are the summands of the $\Lambda^\lambda V$;
the ordinary Schur functor sends $\Lambda^\lambda V$ to $M(\lambda)\otimes\mathrm{sgn}$ and
the twisted one to $M(\lambda)$; the
Hom spaces on both sides have the dimension of the number of $(\lambda,\mu)$ double cosets;
and the functor is faithful on them because $\Lambda^\mu V\subseteq V^{\otimes d}$ and
$V^{\otimes d}$ is projective and injective with socle detected by the Schur functor.  The
second proof of Section~\ref{sec:periodicity-tilting} uses this argument instead of
Corollary~1.8(4).

\paragraph{Harman, Theorem~2.5, the modulus (counterexample at the printed period).}  The
theorem's printed modulus allows, at $r=p=3$, the period $3$.  But the stable incidence
composite through $3$-subsets has scalar $\binom73\equiv2$ at rank $7$ and
$\binom{10}3\equiv0$ at rank $10$ in $\mathbb F_3$; both ranks exceed $2r$, so no functor
identifying these stable maps can respect this composite at period $3$.  A period that is a
power of $p$ strictly greater than $r$ avoids the problem, and that is what we use.  We do not know whether this has been noted before; we record it because our statement
depends on the correct modulus.

\paragraph{Harman, Theorem~2.9 (omitted justification).}  The proof applies Lemma~2.10 to a
restricted right adjoint without showing that it lands in the Serre subcategory where the
lemma applies.  The conclusion needed, that the Serre quotient by the kernel of the twisted
Schur functor is its full image with ambient kernels and subobjects, follows from the
pullback-and-graph argument of Section~\ref{sec:periodicity-oti} (compare
\cite[Lemma~8.5.1]{hone2025}).  This is not a counterexample to Theorem~2.9.

\paragraph{Hone et al., Theorem~8.5.3, at the boundary of Definition~6.1.1
(counterexample).}  Take $p=q=2$ and $d=4$; the label $(3,1)$ is $2$-stable, since
$\lambda_1=3\ge(4+2)/2$, and it truncates to the $2$-singular label $(1,1)$.  Let $W$ be the
sum-zero subspace of $\mathbb F_2^4$ (a submodule of the permutation module of $S_4$) and $C$
the line of the all-ones vector.  In $W/C$ with basis $a=e_1+e_2$, $b=e_1+e_3$, the
transposition $(34)$ acts by $a\mapsto a$, $b\mapsto b+a$.  So $W/C$ is free over
$\mathbb F_2[\langle(34)\rangle]$ and is killed by semisimplification, although it is a
non-zero object of the stable image category: it is a subquotient of $E_1(4)$, which belongs to
that category, and the category is closed under ambient subquotients.  The equivalence asserted in
\cite[Thm.~8.5.3]{hone2025} is therefore not faithful at that boundary.  The example does not
refute Theorem~8.4.5 of the same paper: the descended functor can be exact while killing this
object.  This is why the proof of Section~\ref{sec:periodicity-oti} detects the constant
through a map out of the simple unit and never uses faithfulness.  With strict inequality in
Definition~6.1.1(2b) the boundary case does not arise.

\paragraph{Hone et al., Lemma~8.2.1 (counterexample), and the proof of Theorem~8.4.5.}
Lemma~8.2.1 of \cite{hone2025} states: for a highest-weight category $(\mathcal C,\Lambda)$, a
maximal $\lambda$, and an additive functor $\varphi$ to an abelian category that is exact on
the Serre subcategory of $\Lambda\setminus\{\lambda\}$ and whose comparison maps
$\varphi(\Delta_\lambda)\to\tilde\varphi(\Delta_\lambda)$ and
$\tilde\varphi(\nabla_\lambda)\to\varphi(\nabla_\lambda)$ are isomorphisms ($\tilde\varphi$
the tilting derived functor), the functors $L\varphi$, $\tilde\varphi$, $R\varphi$ agree and
$\tilde\varphi$ is a direct summand of $\varphi$.  This is false.  Let $\mathcal C$ be the
representations of the Kronecker quiver (vertices $1,2$, two arrows $1\to2$) over any field,
ordered $1<2$.  Then $\Delta_1=\nabla_1=S_1$, $\Delta_2=S_2$ is projective, $\nabla_2=I_2$ is
the injective of dimension vector $(2,1)$, and the tilting objects are the sums of $S_1$ and
$I_2$.  Let $\psi(X)$ be the sum of the images of all maps from $S_1\oplus S_2\oplus I_2$ to
$X$.  It is additive, exact on the Serre subcategory generated by $S_1$, and the identity on
$S_1$, $S_2$, $I_2$ and the maps between them, so both hypotheses hold and $\tilde\psi$ is the
identity.  For the projective cover $P_1$ of $S_1$, of dimension vector $(1,2)$,
$\mathrm{Hom}(S_1,P_1)=\mathrm{Hom}(I_2,P_1)=0$ and $\mathrm{Hom}(S_2,P_1)=k^2$, so
$\psi(P_1)=S_2^2$, of which $P_1=\tilde\psi(P_1)$ is not a direct summand; and
$L\psi(S_1)=0\ne S_1$.  The error in the proof is the claim that $P_\bullet\to
i_*i^*P_\bullet$ is a homotopy equivalence: a short exact sequence of complexes with
contractible kernel gives a quasi-isomorphism, which an additive functor need not preserve.
Lemma~8.4.4 uses Lemma~8.2.1 to identify the OTI functor with its tilting derived functor on
all stable objects, and Theorem~8.4.5 rests on Lemma~8.4.4.  So the published proof of
Theorem~8.4.5 is incomplete.  This is not a counterexample to Theorem~8.4.5, which may hold
for the OTI functor by an argument specific to it.  The point was found by Astra in a model
reconstruction of the proof and checked by hand by a second model.

Further minor points found in the same reading (a misprint in \cite[\S7.3]{hone2025},
where $M^{\lambda^t}\otimes\mathrm{sgn}$ should read $M^\lambda\otimes\mathrm{sgn}$; the
justification of \cite[Prop.~2.7]{harman2015} at $\ell=1$) and, in the same reconstruction, two points
in \cite[\S8.3, \S7.3]{hone2025} (the genericity condition before Definition~8.3.1 is not
sufficient for Lemma~8.3.2 in characteristic $2$, while coefficients linearly independent
over $\mathbb F_p$, which the paper gives as an example, suffice; and in the proof of
Lemma~7.3.4 the dominance inequality should read $\ge$, with target degree $M-q$) do not
affect anything here and were not included in those notes.

% sec8-10-tail.tex -- Sections 8-10 of the C3 paper (draft 2026-11-08).
% Sources: requests-20261106/c3-publication-prior-art-v2.md; notes-20261025/c3-sources.md;
% publish/summaries/c3-*.md; OUTLINE.md section 9; PLAN.md rows G-C3.1-G-C3.3.

\section{Relation to the literature}\label{sec:literature}

\paragraph{The same family.}  Three sources treat the PC or Nullstellensatz degree of
$\mathrm{onto\text{-}FPHP}^{n+k}_n$ with $p\mid k$.
\begin{itemize}
\item Riis's criterion~\cite{riis1993independence} (Theorem~8 of \cite{razborov2001php}) gives
  the upper bound $p^{v_p(k)}$ for every $n$ (Proposition~\ref{prop:riis}).  Our results show
  that it is attained: at the Lean tier for $v_p(k)=1$ from $n=2p-3$ (A), and at the written
  tier for every $k$ from $n=2p^{v_p(k)}-1$ (E).
\item Beame and Riis~\cite{beame1998counting} prove Nullstellensatz degree at least
  $2^\ell-1$ (multiplier degree) for surplus $p^\ell$ and $n\ge((p+2)^\ell-p^\ell)/2$, and at
  most $p^\ell-1$ for $n\ge p^\ell$, and list PC lower bounds for this family as open.  Their designs are the input of D.
\item Riis and Sitharam~\cite[Thm.~33]{riis1997nonconstant} state PC degree at least $2^\ell-1$
  for surplus $p^\ell$ once $N\ge(21+3p^2)(2^\ell-2)-p$.  We have read the statement and the
  preceding lemmas (pp.~36--38, same linear encoding) but not checked the proof, which
  delegates to a tree translation (Lemmas~31, 32) that we have not reconstructed.  We draw no
  conclusion about its correctness, and none from the fact that Razborov's later survey asks
  Question~7 without citing it.  For comparison: at $\ell=2$ it gives $3$ for
  $N\ge2(21+3p^2)-p$, against $p+1$ from $n=2p-1$ (B, Lean) and $p^2$ from $n=2p^2-3$ (F,
  written); at $p=2$, $\ell=3$ it gives $7$ for $N\ge196$, against $5$ from $n=7$ (C, Lean)
  and $8$ from $n=15$ (E, written).
\end{itemize}

\paragraph{Razborov's Question~7.}  The mathematical statements here concern the exact
eventual degree and, for surplus $p^\ell$, growth as $\ell$ varies; under E's named inputs
that degree is $p^\ell$.  We make no claim of resolving the historical formulation of
Question~7 or of priority over Theorem~33 of \cite{riis1997nonconstant}.  None of the Lean
results bears on growth.

\paragraph{Other characteristics and other principles.}  Buss, Grigoriev, Impagliazzo and
Pitassi~\cite{bgip2001} prove linear PC degree for Tseitin formulas mod $p$ in
characteristic $q\ne p$ (when the field contains a root of unity of order $p$) and for counting principles
in other characteristics; Alekhnovich and Razborov~\cite{alekhnovich2003nonbinomial} extend
linear bounds to non-binomial cases; Buss, Impagliazzo, Kraj\'\i\v cek, Pudl\'ak, Razborov
and Sgall~\cite{bikprs1996} prove Nullstellensatz lower bounds for counting principles mod
$q$ over $\mathbb Z_m$ with $q\nmid m$.  All concern the regime in which the characteristic
does not divide the relevant modulus, and none treats the onto principle with $p\mid k$.
Impagliazzo, Pudl\'ak and Sgall~\cite{ips1999} give Razborov's bound for the injective
principle, where the onto and functional axioms are absent and the bound does not depend on
$p$.  The applications of general degree methods~\cite{miksa2015generalized} that we examined do
not yield bounds for the linear onto functional system when $p$ divides $k$, and in the
sources we examined we found no suitable low-degree reduction from the cited counting or
Tseitin principles to this system.

\paragraph{The algebraic inputs.}  The proofs of D, E and F use published results from
representation stability and modular representation theory: FI-modules~\cite{churchellenberg},
matrix loci~\cite{rhoades2024}, stable periodicity~\cite{harman2015},
semisimplification~\cite{hone2025}, chessboard complexes~\cite{shareshianwachs}, Moore spaces
and the Dade group~\cite{yalcin2016}, and endosplit resolutions~\cite{miller2024}.  We make no claim about whether they have been
used in proof complexity before.

\paragraph{Search.}  The searches behind this section (web, ECCC, and the citations of the
sources; October and November 2026), with the list of sources opened and of those read only as
abstracts, are published with each of the findings behind A--C, under \emph{prior art}.

\section{How these results were produced, and what was checked}\label{sec:process}

\paragraph{Who did what.}  All results in this paper were found and proved by AI research
agents (Claude models by Anthropic and Codex models by OpenAI) working in an automated
research harness in October and November 2026.  The author built and ran the harness, chose
the problems, directed the runs, read the outputs and decided what to publish.  He is not a
specialist in this area and does not follow every proof in detail himself.  He publishes the
results as they stand, with the evidence for each stated by its tier, and takes
responsibility for publishing them.  The harness records every claim together with its
evidence, its reviews and the files behind it; those records are the basis of the roles stated
below.  Results A--C come from runs that produce Lean proofs; D--G from one exploratory run
(20261104-0759) with two Claude and two Codex agents, followed by adjudication.

\paragraph{The Lean results.}  The formal statements are supported by kernel-checked proof
terms.  Their interpretation as the claims of this paper also depends on the definitions and
the correspondence established in Section~\ref{sec:defs}.  The library is public, with the exact toolchain and a
manifest; each finding records the commit it was checked at.  Its checks rebuild the proof modules with
the recorded toolchain and then check the published declarations' types and axiom
dependencies with \texttt{\#check} and \texttt{\#print axioms}.  Every theorem uses
only \texttt{propext}, \texttt{Classical.choice} and \texttt{Quot.sound}.  On 6 November 2026 a separate model reproduction rebuilt the named proof targets from the
public repository and checked the published declarations; hardware, timings and memory are
recorded in the ``How to verify'' section of each finding.
For A and B, a second agent in the same run supplied Lean re-proofs, reusing between one and
nine of the first agent's lemmas for A and six for each statement of B.  For C, a second agent
independently checked the level-3 certificate with a program written from the Lean
definitions; this was not a second Lean re-proof.

\paragraph{The written results.}  Each proof of D--G was written by one agent and
reconstructed by at least one other, of the other vendor in most cases; a separate model
referee with no contact with the run re-read the primary sources for each of the three
mechanisms; and the adjudicator Astra (a Codex model outside the run) checked each against the
primary sources and wrote the scoped verdicts we report.  Section~\ref{sec:growth-provenance}
and the corresponding paragraphs of Sections~\ref{sec:periodicity}--\ref{sec:bands} give the
roles per result.  Where reviewers co-developed an argument before reviewing it, the run
records say so.  The paper's later prose reconstructions (Section~\ref{sec:lean},
Appendix~\ref{app:tilting}) and added corollaries were checked in a separate model review of
the paper, which led to corrections.  This is a demanding process, but it is model review,
and we do not present it as more.

\paragraph{What was not done.}  No human expert has reviewed any proof in this paper.  The
external theorems named as inputs were checked at the level of their statements and the
interfaces we use, not re-proved; the exception is \cite[Thm.~8.4.5]{hone2025}, whose proof
was reconstructed, with the result described in Sections~\ref{sec:periodicity}
and~\ref{sec:errata}.  D--G are not formalised.  The finite computations are
exact, but are evidence only for the cases computed.  We would welcome corrections, and will
record them with the published findings.

\section{Open problems}\label{sec:open}

\begin{enumerate}
\item \emph{Formalise E.}  A first step is planned for a revision of this paper: a Lean proof that
  periodicity, and hence the exact degree $p^{v_p(k)}$, follows from the existence of a functor
  with the properties used in Section~\ref{sec:periodicity-oti} or Appendix~\ref{app:tilting}, so that the remaining trust is
  isolated in explicit functor properties and the listed external inputs.  The inputs are deep; a formalisation of F in
  the bands $d<2p$, which needs only the connectivity theorem, is closer.
\item \emph{Complete the proof of \cite[Thm.~8.4.5]{hone2025}} for the semisimplification
  functor, without their Lemma~8.2.1.  This would make Theorem~\ref{thm:periodicity}(b) and
  the threshold $n\ge2q-2$ unconditional.
\item \emph{The uniform spanning statement.}  Prove the spanning statements
  (A)$(n+k,n,r)$, $r=2,3,4$, uniformly over $\mathbb F_3$ for $n\ge7$ and $9\mid k$ ((A) fails
  off the stable range, Section~\ref{sec:bands-reduction}).  Alternatively, formalise
  Lemmas~\ref{lem:hole-avoid} and~\ref{lem:pigeon-ind} and import kernel-checked square
  certificates $S(2)$, $S(3)$, $S(4)$ in the formal spanning definitions; this would discharge
  the hypotheses of the Lean theorem behind Corollary~\ref{cor:f3}.
\item \emph{Thresholds.}  The degree is $p^{v_p(k)}$ from $n=2p^{v_p(k)}-1$ (E) and, for
  $v_p(k)\le2$, from $2p^{v_p(k)}-3$ (A, F).  What is it below?  Over $\mathbb F_2$ with $k=4$
  it is $3$ at $n=3,4$ and $4$ from $n=5$, so the threshold of E is not sharp.
\item \emph{The ideal identity beyond $p^2$.}  Investigate the full identity
  $(I\cap J)_d=(IJ)_d+\sum_{t=1}^d\sigma_t\mathcal A_{d-t}$, the form tested by exact
  computation in the run, for $d\ge p^2$.  For $d<p^2$ the proof of
  Section~\ref{sec:bands-high} establishes the stronger form with $\sigma_1$ and $\sigma_p$
  only; at $d=p^2$ the elementary abelian subgroup is no longer Sylow and the transfer
  argument stops.
\item \emph{Designs of degree $p^\ell-1$.}  Beame and Riis leave the Nullstellensatz degree for
  surplus $p^\ell$ between $2^\ell-1$ and $p^\ell-1$ for odd $p$.  Unit designs of degree
  $p^\ell-1$ would close that gap and, through Theorem~\ref{thm:padding}, give the exact PC
  degree $p^\ell$ on the inputs of Section~\ref{sec:growth} alone.
\item \emph{Riis--Sitharam, Theorem~33.}  Check its proof.
\end{enumerate}


\appendix
% appA-tilting.tex -- Appendix A of the C3 paper (draft 2026-11-08): the second proof of Theorem 5.1(a).
% Sources: claude-20261104-c3-closure/01-PC-CLOSURE-PERIODICITY.md (c1: sections 1-2, Lemma B, repairs R1-R3);
% notes-20261104-referee-c3-closure.md (cold referee, incl. the p=2 counterexample and the fixes);
% notes-20261104-referee-c3-closure-x2.md (x2: scalar extension); Astra
% codex-20261104-runB-adjudication/03-C1-INTERFACES-AND-ERRATA.txt section 4. Harman statements checked against
% arXiv:1509.06414v3 (harman-local.txt) on 2026-11-08.

\section{The second proof of periodicity}\label{app:tilting}

This appendix proves Theorem~\ref{thm:periodicity}(a), for $m,m'>2D$.  It is the proof on
which part (a), and hence the exact degree for $n\ge2q-1$, rests; it does not use
\cite{hone2025}.  It extends natural isomorphisms only between exact functors, so the flaw in
\cite[Lemma~8.2.1]{hone2025} (Section~\ref{sec:errata}), which concerns additive functors, does
not affect it.  It transports the closure of Section~\ref{sec:periodicity-closure} along
Harman's stable equivalences~\cite{harman2015}.  Those equivalences are compatible with
induction and restriction of permutation modules; the closure needs compatibility for
\emph{arbitrary} submodules, which we prove on the Schur-algebra side with standard facts on
tilting modules (Lemma~\ref{lem:B}).

\paragraph{Status.}  T-WRITTEN.  The proof was written by a Claude agent (c1) in run
20261104-0759, reviewed by a model referee with no contact with the run (who found the
characteristic-$2$ error of Section~\ref{sec:errata} and supplied the replacement argument of
Lemma~\ref{lem:R1}) and by a Codex agent (x2, who supplied the two-way scalar extension), and
adjudicated by Astra, who accepted the coherence argument of Section~\ref{app:coherence}
conditional on the named inputs and without re-proving their foundations.  The inputs are:
Harman's Theorem~2.5 with the corrected modulus and his Corollary~1.6 and Proposition~2.7
(for $\ell=1$, which the referee checked directly); Ringel duality and the equivalence
$K^b(\mathrm{Tilt})\simeq D^b$ for highest-weight categories with finitely many
weights~\cite{ringel1991,happel1988}; the facts that tensor products of tilting modules for
$GL_N$ are tilting and that restriction to a Levi subgroup preserves tilting
modules~\cite{wang1982,donkin1998,mathieu1990}; the Hom-dimension formula for permutation
modules~\cite[13.19]{james1978}; and standard properties of the Schur
functor~\cite{green1980}.

Throughout, $k=\overline{\mathbb F}_p$ (Section~\ref{sec:periodicity-closure} reduces to this
case), $q$ is a power of $p$ with $q>D$, and $m\equiv m'\pmod q$ with $m,m'>2D$.  For a
partition $\lambda$ and $m\ge|\lambda|+\lambda_1$ put $\lambda(m)=(m-|\lambda|,\lambda)$.

\subsection{Harman's equivalences, corrected}\label{app:harman}

\paragraph{Permutation and Young modules.}  Let $\mathrm{Perm}^{\le r}_m$ be the category of
direct sums of permutation modules $M(\lambda(m))$ of $S_m$ with $|\lambda|\le r$, and
$\mathrm{Young}^{\le r}_m$ its Karoubi envelope, the direct sums of Young modules
$Y(\lambda(m))$.  Homomorphisms between permutation modules have the stable Carter--Lusztig
bases of \cite[\S2]{harman2015}; a \emph{stable map} is one whose coefficients in these bases do
not depend on $m$.

\begin{theorem}[Harman~{\cite[Thm.~2.5, Cor.~2.6]{harman2015}}, modulus corrected]\label{thm:harman25}
Let $r\ge0$ and let $q$ be a power of $p$ with $q>r$.  (For $r=0$ the categories consist of
sums of the trivial module and the equivalence is the evident one.)  For $m,m'>2r$ with $m\equiv m'\pmod q$,
sending $M(\lambda(m))$ to $M(\lambda(m'))$ and each stable map to the stable map with the
same coefficients defines an equivalence
$E^{\mathrm{P}}\colon\mathrm{Perm}^{\le r}_m\to\mathrm{Perm}^{\le r}_{m'}$, and hence an
equivalence $E^{\mathrm Y}$ of Young categories.
\end{theorem}

The structure constants of composition are binomial coefficients in $m$ of degree at most
$r$, and $\binom md\bmod p$ depends on $m$ only modulo the least power of $p$ above $d$
(Lucas).  The printed modulus $p^{\lceil\log_pr\rceil}$ is too small when $r$ is a power of
$p$ (Section~\ref{sec:errata}); with $q>r$ the proof of \cite{harman2015} applies unchanged.
For $\ell=1$, \cite[Prop.~2.7]{harman2015} states that $E^{\mathrm P}$ commutes with
restriction to $S_{m-1}$ and induction from $S_{m-1}$, up to natural isomorphism, on permutation
modules.  The relevant maps are explicit: restriction sends the basis element $h_\tau$ to
$\sum h_{\tau-e_{ij}}$ and induction adjoins a singleton block, with coefficients $0$ and $1$
that do not depend on $m$; so this can be checked directly.

\paragraph{Schur algebras.}  For $N\ge m$ let $S(N,m)$ be the Schur algebra, so that
$S(N,m)$-modules are the polynomial representations of $GL_N$ of degree $m$, and let
$\mathcal S^{\le r}(N,m)$ be the Serre subcategory generated by the simples
$L(\lambda(m)')$, $|\lambda|\le r$.  It is a highest-weight category whose tilting objects
are the $T(\lambda(m)')$~\cite[Prop.~1.3 and proof of Thm.~2.8]{harman2015}.  Let $\mathcal F$
be the Schur functor and $\tilde{\mathcal F}=\mathcal F(-)\otimes\mathrm{sgn}$; it is exact and
$\tilde{\mathcal F}(T(\lambda'))=Y(\lambda)$~\cite[Cor.~1.8(1),(2)]{harman2015}.

\begin{lemma}[full faithfulness on tilting modules]\label{lem:R1}
In every characteristic, $\tilde{\mathcal F}$ restricts to an equivalence
$\mathrm{Tilt}(\mathcal S^{\le r}(N,m))\to\mathrm{Young}^{\le r}_m$.
\end{lemma}

\begin{proof}
Harman deduces this from \cite[Cor.~1.8(4)]{harman2015}, which is false at $p=2$
(Section~\ref{sec:errata}); here is a proof that does not use it.  The exterior powers
$\Lambda^\mu V=\Lambda^{\mu_1}V\otimes\Lambda^{\mu_2}V\otimes\cdots$ are tilting, and
$T(\mu')$ is a summand of $\Lambda^\mu V$ with multiplicity one; so the tilting objects of
$\mathcal S^{\le r}(N,m)$ are the summands of the $\Lambda^{\lambda(m)}V$, $|\lambda|\le r$,
and $\tilde{\mathcal F}(\Lambda^\mu V)=M(\mu)$.

\emph{Dimensions.}  Since $\Lambda^\lambda V$ has a Weyl filtration and $\Lambda^\mu V$ a
good filtration,
$\dim\mathrm{Hom}(\Lambda^\lambda V,\Lambda^\mu V)=\sum_\nu(\Lambda^\lambda V:\Delta(\nu))
(\Lambda^\mu V:\nabla(\nu))=\sum_\rho K_{\rho\lambda}K_{\rho\mu}$, the number of
$(S_\lambda,S_\mu)$ double cosets in $S_m$, which is
$\dim\mathrm{Hom}_{S_m}(M(\lambda),M(\mu))$ in every
characteristic~\cite[13.19]{james1978}.

\emph{Faithfulness.}  $\Lambda^\mu V$ embeds in $V^{\otimes m}$ (antisymmetrisation is injective
in every characteristic).  For $N\ge m$, $V^{\otimes m}$ is projective and injective, and every
simple submodule $L$ of it has $\mathcal F(L)\ne0$~\cite[\S6]{green1980}.  If
$f\colon\Lambda^\lambda V\to\Lambda^\mu V$ is non-zero, its image is a non-zero submodule of
$V^{\otimes m}$, so it contains such an $L$, and since $\mathcal F$ is exact,
$\mathcal F(\operatorname{im}f)\ne0$, so $\tilde{\mathcal F}(f)\ne0$.

A faithful functor between Hom spaces of equal finite dimension is bijective on them.  So
$\tilde{\mathcal F}$ is fully faithful on the $\Lambda^{\lambda(m)}V$, hence on their summands;
its essential image is $\mathrm{Young}^{\le r}_m$.
\end{proof}

\begin{proposition}[Harman's Theorems~2.8 and~2.9, with the repairs]\label{prop:harman-eq}
Let $N\ge m$, $N'\ge m'$.  There is an equivalence of highest-weight categories
$\Theta\colon\mathcal S^{\le r}(N,m)\to\mathcal S^{\le r}(N',m')$ with
$\Theta(T(\lambda(m)'))\cong T(\lambda(m')')$, together with a natural isomorphism
$\zeta\colon\tilde{\mathcal F}\Theta\cong E^{\mathrm Y}\tilde{\mathcal F}$ on tilting objects.
$\Theta$ descends to an equivalence of abelian categories
$E\colon\mathcal R^{\le r}_m\to\mathcal R^{\le r}_{m'}$ between the essential images
$\mathcal R^{\le r}_m=\tilde{\mathcal F}(\mathcal S^{\le r}(N,m))\subseteq\mathrm{Rep}(S_m)$,
whose kernels, cokernels and subobjects are the ordinary ones, and $E$ restricts to
$E^{\mathrm P}$ on permutation modules up to the isomorphisms $\zeta$.
\end{proposition}

\begin{proof}
By Lemma~\ref{lem:R1} and Theorem~\ref{thm:harman25}, composing $\tilde{\mathcal F}$, then
$E^{\mathrm Y}$, then a quasi-inverse of $\tilde{\mathcal F}$ at $(N',m')$ gives an equivalence
$\mathrm{Tilt}(\mathcal S^{\le r}(N,m))\to\mathrm{Tilt}(\mathcal S^{\le r}(N',m'))$ that
preserves labels.  Ringel duality~\cite[Cor.~1.6]{harman2015} extends it to $\Theta$; on
tilting objects $\Theta$ agrees with the given equivalence up to natural isomorphism, since
both are determined by the image of a full tilting object and its endomorphism ring.  (This is
a property of the construction, not its definition; we use only this property.)

For $|\lambda|\le r<m/2$ the first part $m-|\lambda|$ of $\lambda(m)$ exceeds every other
part, so $\lambda(m)$ is $p$-singular if and only if $\lambda(m')$ is.  Hence $\Theta$ maps the
Serre subcategory generated by the $p$-singular simples, which is the kernel of
$\tilde{\mathcal F}$, onto its counterpart.  It remains to identify the Serre quotient
$\mathcal S^{\le r}/\ker\tilde{\mathcal F}$ with $\mathcal R^{\le r}_m$.  Harman applies
Gabriel's lemma to $\mathcal S^{\le r}$, which needs a right adjoint landing in
$\mathcal S^{\le r}$; that is not shown (Section~\ref{sec:errata}).  Instead: the right adjoint
$\tilde{\mathcal G}$ of $\tilde{\mathcal F}$ on all of $S(N,m)$-mod has invertible counit, so
$S(N,m)\text{-mod}/\ker\tilde{\mathcal F}\simeq\mathrm{Rep}(S_m)$ by Gabriel's
lemma~\cite{gabriel1962}; and since $\mathcal S^{\le r}$ is a Serre subcategory, for
$X\in\mathcal S^{\le r}$ and $T\le\tilde{\mathcal F}X$ the pullback
$X_T=X\times_{\tilde{\mathcal G}\tilde{\mathcal F}X}\tilde{\mathcal G}T$ is a subobject of $X$
in $\mathcal S^{\le r}$ with $\tilde{\mathcal F}X_T=T$, and lifting graphs the same way gives
every morphism as a roof inside $\mathcal S^{\le r}$.  So the quotient is the full subcategory
$\mathcal R^{\le r}_m$, closed under submodules and quotients in $\mathrm{Rep}(S_m)$.  The
exact equivalence $\Theta$ descends to an exact equivalence $E$.
\end{proof}

\subsection{Compatibility with induction and restriction}\label{app:lemmaB}

Write $E_m$ for the equivalence of Proposition~\ref{prop:harman-eq} at weight $\le D$ between
$S_m$ and $S_{m'}$, and $E'_{m-1}$ for the one at weight $\le D-1$ between $S_{m-1}$ and
$S_{m'-1}$ (note $m-1>2(D-1)$ and $m-1\equiv m'-1$).

\begin{lemma}[B]\label{lem:B}
There are natural isomorphisms
\[
  E_m\circ\mathrm{Ind}\cong\mathrm{Ind}\circ E'_{m-1}\ \text{ on }\mathcal R^{\le D-1}_{m-1},\qquad
  E'_{m-1}\circ\mathrm{Res}\cong\mathrm{Res}\circ E_m\ \text{ on }\mathcal R^{\le D-1}_{m},
\]
which on permutation modules are the isomorphisms of \cite[Prop.~2.7]{harman2015} with
$\ell=1$.
\end{lemma}

\begin{proof}
Fix $N\ge\max(m,m')+1$.  On the Schur side take
$\Phi_{\mathrm{Ind}}=-\otimes V\colon S(N,m-1)\text{-mod}\to S(N,m)\text{-mod}$ and let
$\Phi_{\mathrm{Res}}\colon S(N,m)\text{-mod}\to S(N-1,m-1)\text{-mod}$ be the part of weight
$1$ for the last factor of the Levi subgroup $GL_{N-1}\times GL_1$.

\emph{(i)}  $\tilde{\mathcal F}(M\otimes V)\cong\mathrm{Ind}\,\tilde{\mathcal F}(M)$ and
$\tilde{\mathcal F}(\Phi_{\mathrm{Res}}M)\cong\mathrm{Res}\,\tilde{\mathcal F}(M)$, naturally:
the $(1^m)$-weight space of $M\otimes V$ is $\bigoplus_i M_{(1^m)-\varepsilon_i}$, which is
the induced module, and the transposition $(m\ N)$ carries the $(1^m)$-weight space of $M$ to
the weight $(1^{m-1},0,\dots,0,1)$, $S_{m-1}$-equivariantly; the sign twist commutes with both.
\emph{(ii)}  Both functors are exact and send tilting modules to tilting
modules~\cite{wang1982,donkin1998,mathieu1990}; a weight summand of a tilting module for the
Levi subgroup is tilting.
\emph{(iii)}  By branching, $-\otimes V$ maps $\mathcal S^{\le D-1}(N,m-1)$ into
$\mathcal S^{\le D}(N,m)$ (composition factors of $\Delta(\mu)\otimes V$ lie below
$\mu+\varepsilon_i$), and $\Phi_{\mathrm{Res}}$ maps $\mathcal S^{\le D-1}(N,m)$ into
$\mathcal S^{\le D-1}(N-1,m-1)$.
\emph{(iv)}  By (i), both functors preserve the kernels of $\tilde{\mathcal F}$.

\emph{Step 1 (tilting level).}  Let $\phi\in\{\mathrm{Ind},\mathrm{Res}\}$ with Schur-side
partner $\Phi$, and $T$ a tilting object.  Then, naturally in $T$,
\[
  \tilde{\mathcal F}\Theta\Phi T\cong E^{\mathrm Y}\tilde{\mathcal F}\Phi T
  \cong E^{\mathrm Y}\phi\tilde{\mathcal F}T\cong\phi E^{\mathrm Y}\tilde{\mathcal F}T
  \cong\phi\tilde{\mathcal F}\Theta T\cong\tilde{\mathcal F}\Phi\Theta T,
\]
using $\zeta$, (i), \cite[Prop.~2.7]{harman2015} (on permutation modules, hence on their
summands), $\zeta$ again, and (i).  Both $\Theta\Phi T$ and $\Phi\Theta T$ are tilting by
(ii), and $\tilde{\mathcal F}$ is fully faithful on tilting objects (Lemma~\ref{lem:R1}), so
$\Theta\Phi\cong\Phi\Theta$ naturally on tilting objects.

\emph{Step 2 (all objects).}  The categories involved are highest-weight with finitely many
weights, so $K^b(\mathrm{Tilt})\to D^b$ is an equivalence~\cite{ringel1991,happel1988}.  An
exact functor induces a functor on $D^b$ computed termwise on tilting complexes, and the natural
isomorphism of Step~1 applied termwise gives one on $K^b(\mathrm{Tilt})$, hence on $D^b$.
Taking $H^0$ on the heart gives $\Theta\Phi\cong\Phi\Theta$ on all objects, restricting to the
isomorphism of Step~1 on tilting objects.

\emph{Step 2$'$ (restriction in weight).}  The same argument shows that $\Theta$ at weight
$\le D$, restricted to weight $\le D-1$, is isomorphic to $\Theta$ at weight $\le D-1$: both
are exact and agree with the same Young-module equivalence on tilting objects.  This is needed
because $C_t\cap F_{D-1}$ is transported at weight $D$ and then restricted.

\emph{Step 3 (descent).}  By (iv) and the universal property of Serre quotients for exact
functors and natural transformations, the isomorphism descends to $E\circ\phi\cong\phi\circ E$
on the image categories.  On permutation modules it is the isomorphism of Step~1, that is,
Harman's.
\end{proof}

\subsection{Coherence}\label{app:coherence}

Two points make the isomorphisms usable for subobjects; Astra accepted both, conditional on the
inputs above.

\emph{Independence of $N$.}  Induction uses $S(N,m-1)$ and restriction lands in
$S(N-1,m-1)$.  Truncation to the weights with last coordinate $0$ is an exact Morita
equivalence $S(N,m-1)\text{-mod}\to S(N-1,m-1)\text{-mod}$ for $N-1\ge m-1$; it preserves
tilting objects and labels and commutes with $\tilde{\mathcal F}$.  Applying Step~2 to
$\Theta^{(N)}$ and to $\Theta^{(N-1)}$ conjugated by truncation (both exact, both isomorphic
to the same Young functor on tilting objects) gives a natural isomorphism between them that is
the identity on Young modules.  So $E'_{m-1}$ does not depend on $N$, and the output of
restriction can be fed into induction.

\emph{Fixed identifications.}  Fix, on each side, the descent isomorphism
$\delta\colon E\tilde{\mathcal F}\to\tilde{\mathcal F}\Theta$ and the tilting comparison
$\alpha\colon\tilde{\mathcal F}\Theta\to E^{\mathrm Y}\tilde{\mathcal F}$ of
Proposition~\ref{prop:harman-eq}.  On tilting objects their composite descends, by full
faithfulness of $\tilde{\mathcal F}$ there (Lemma~\ref{lem:R1}), to a natural isomorphism
$\eta\colon E|_{\mathrm{Young}}\to E^{\mathrm Y}$; write $\eta'$ for the same at $m-1$.  The same
$\delta$ and $\alpha$ are used in constructing the isomorphisms of Lemma~\ref{lem:B}.  Let
$c\colon E^{\mathrm P}_m\mathrm{Ind}\to\mathrm{Ind}\,E^{\mathrm P}_{m-1}$ and
$c'\colon E^{\mathrm P}_{m-1}\mathrm{Res}\to\mathrm{Res}\,E^{\mathrm P}_m$ be Harman's
isomorphisms on permutation modules.  Unwinding the descended isomorphisms $\beta$ (induction)
and $\gamma$ (restriction), the Schur-functor isomorphisms of (i) cancel by naturality of
$\delta$ and $\alpha$, and for permutation objects $P$, $Q$
\[
  \beta_P=(\mathrm{Ind}\,\eta'_P)^{-1}\circ c_P\circ\eta_{\mathrm{Ind}\,P},\qquad
  \gamma_Q=(\mathrm{Res}\,\eta_Q)^{-1}\circ c'_Q\circ\eta'_{\mathrm{Res}\,Q}.
\]
All maps used below (the multiplication pieces $g_j$, $\varepsilon_j$, the inclusions, the
axiom maps and the constant) are stable maps defined from the same stable decompositions as
$c$.  For each inclusion $T\to P$ of a submodule into a permutation object, naturality gives a
commuting square whose ambient map is the fixed component above; taking images identifies the
transported subobject.  So in the induction below every subobject is carried to the
corresponding subobject on the smaller board, and no unknown automorphism enters.

\subsection{Transport of the closure}\label{app:transport}

\paragraph{The multiplication step as induction.}  Fix a hole $j$ and let $S_{m-1}$ fix pigeon
$m$.  Multiplication by $x_{mj}$ maps $\mathrm{Res}\,F_{D-1}(m)$ into the summand
$P_j\subseteq\mathrm{Res}\,F_D(m)$ spanned by matchings containing the edge $(m,j)$; call this
map $g_j$.  $P_j$ is a permutation module of $S_{m-1}$ of weight $\le D-1$ (matchings of size
$\le D-1$ on pigeons $[m-1]$ and holes other than $j$).  Let
$\varepsilon_j\colon\mathrm{Ind}\,P_j\to F_D(m)$ be adjoint to the inclusion.  Both are stable
maps with coefficients $0$ and $1$.  For an $S_m$-submodule $S\subseteq F_{D-1}(m)$,
\[
  \textstyle\sum_i\rho(x_{ij})S=\varepsilon_j\bigl(\mathrm{Ind}\,g_j(\mathrm{Res}\,S)\bigr),
\]
because $\sum_i\rho(x_{ij})S=\sum_{\sigma\in S_m}\sigma\,\rho(x_{mj})\,S$ and the image of
$\mathrm{Ind}\,T\to F_D$ is $S_m\cdot T$.  So every $C_t(m)$ is obtained from permutation
modules of weight $\le D$ (for $S_m$) or $\le D-1$ (for $S_{m-1}$) and stable maps by images,
preimages ($C_t\cap F_{D-1}=\iota^{-1}(C_t)$ for the inclusion $\iota$), finite sums, and
induction and restriction of submodules.  The constant is the stable map $k\to F_D(m)$.

\begin{proof}[Proof of Theorem~\ref{thm:periodicity}(a)]
By induction on $t$, $\eta\bigl(E_m(C_t(m))\bigr)=C_t(m')$ as subobjects of $F_D(m')$.  For
$t=0$, $C_0$ is the sum of the images of the axiom maps, which are stable.  For the step: $E_m$
is exact, so it preserves images, preimages and sums of subobjects, and the subobjects are the
ordinary ones (Proposition~\ref{prop:harman-eq}); stable maps go to the same stable maps up to
$\eta$ (Section~\ref{app:coherence}); induction and restriction of the subobjects
$C_t\cap F_{D-1}$ and $g_j(\mathrm{Res}\,S)$ are transported by Lemma~\ref{lem:B}, with
Step~2$'$ for the change of weight.  Both chains stabilise.  Since $E_m$ is an equivalence, the
composite $k\to F_D(m)\to F_D(m)/C_\infty(m)$ is zero if and only if its image
$k\to F_D(m')\to F_D(m')/C_\infty(m')$ is.  Lemma~\ref{lem:closure} and the reduction to
$\overline{\mathbb F}_p$ (Section~\ref{sec:periodicity-closure}) finish the proof.
\end{proof}

The strict inequality $m,m'>2D$ is used twice: in Theorem~\ref{thm:harman25} and in the
comparison of $p$-singular labels in Proposition~\ref{prop:harman-eq}.  At $m=2D$ this proof
does not apply, and Theorem~\ref{thm:periodicity}(b) rests on the conditional proof of
Section~\ref{sec:periodicity-oti}.


\bibliographystyle{plain}
\bibliography{refs}
\end{document}
